% =====================================================================
%  Poglavje: Analiza in vrednotenje rešitve
%  Vključeno iz diploma-predloga.tex prek % =====================================================================
%  Poglavje: Analiza in vrednotenje rešitve
%  Vključeno iz diploma-predloga.tex prek % =====================================================================
%  Poglavje: Analiza in vrednotenje rešitve
%  Vključeno iz diploma-predloga.tex prek % =====================================================================
%  Poglavje: Analiza in vrednotenje rešitve
%  Vključeno iz diploma-predloga.tex prek \input{podporno/analiza.tex}
%
%  Številke so vpisane iz izhodov skript v CarStockMarketEvaluation/analysis/.
%  Preostale oznake \res{...} (rdeče) je treba še izpolniti:
%    - obremenitveni test (benchmarks/loadtest.py)
%    - sledljivost testov, Jest coverage, Lighthouse
%    - tabeli udeležencev in najdenih težav
%  Slike (PDF) iz analysis/out/ prekopiraj v mapo slike/.
% =====================================================================

\newcommand{\res}[1]{\textcolor{red}{\textbf{[#1]}}}

\chapter{Analiza in vrednotenje rešitve}
\label{ch:analiza}

V tem poglavju sistematično pregleadamo in analiziramo razvito rešitev. Najprej iz zahtev in
izhodišč, ki smo si jih zastavili na začetku, izpeljemo raziskovalna vprašanja in za vsako določimo merilo
uspešnosti (razdelek~\ref{sec:rq}). Nato opišemo merilno okolje in postopke
(razdelek~\ref{sec:metodologija}) ter natančno definiramo vse uporabljene mere
(razdelek~\ref{sec:mere}). Sledijo rezultati: značilnosti zbranih
podatkov~(\ref{sec:podatki}), natančnost ocene
vrednosti~(\ref{sec:natancnost}), občutljivost algoritma na
parametre~(\ref{sec:obcutljivost}), latenca posameznega
vrednotenja~(\ref{sec:latenca}) in obnašanje pod
obremenitvijo~(\ref{sec:obremenitev}). Zatem preverimo funkcionalne in
nefunkcionalne zahteve (\ref{sec:funkcionalne}, \ref{sec:nefunkcionalne}) ter
predstavimo uporabniški preizkus~(\ref{sec:uporabniki}) in simulacijo
uporabe~(\ref{sec:simulacija}). Poglavje sklenemo z obravnavo groženj
veljavnosti~(\ref{sec:groznje}) in povzetkom ugotovitev glede na
zastavljena izhodišča~(\ref{sec:povzetek-analize}).

% =====================================================================
\section{Raziskovalna vprašanja in merila uspešnosti}
\label{sec:rq}

Na začetku diplomskega dela (poglavje~\ref{ch:uvod}) smo si zastavili, da je
mogoče uporabniku na podlagi zbranih tržnih podatkov in primerjave s podobnimi
vozili ponuditi \emph{uporabno} in \emph{dovolj natančno} oceno tržne vrednosti
vozila in da vključitev takega sistema izboljša uporabniško izkušnjo. Iz tega
izhodišča in iz nefunkcionalnih zahtev (razdelek~\ref{sec:zahteve}) izpeljemo
štiri raziskovalna vprašanja (RV1--RV4) in za vsako določimo merilo uspešnosti.

\emph{Mediana absolutne odstotne napake} (MdAPE) je osrednja mera natančnosti v
tem poglavju; nanjo se sklicujejo merila uspešnosti in vsi nadaljnji razdelki.
Izračunamo jo tako, da za vsak testni primer določimo absolutno odstotno napako,
torej za koliko odstotkov se ocenjena cena razlikuje od dejanske, nato pa čez
vse primere vzamemo mediano teh vrednosti. MdAPE $=15\,\%$ tako pomeni, da je pri
polovici vozil ocena znotraj $15\,\%$ od dejanske cene, pri drugi polovici pa je
odstopanje večje; manjša kot je vrednost, natančnejše so ocene. Mediano
uporabljamo namesto povprečja, ker je manj občutljiva na posamezne skrajne
napake. Natančne definicije MdAPE in preostalih mer so v
razdelku~\ref{sec:mere}.

Raziskovalna vprašanja so:

\begin{itemize}
  \item \textbf{RV1 --- natančnost ocene.} Kako natančno sistem oceni tržno
        vrednost vozila v primerjavi z dejanskimi oglasi na trgu?
        \emph{Merilo:} MdAPE je pri vozilih z vsaj $10$ primerljivimi
        vozili manjša od $15\,\%$, ocena pa je bistveno natančnejša od
        preproste (naivne) referenčne metode.
  \item \textbf{RV2 --- odvisnost natančnosti od podatkov.} Od česa je
        natančnost ocene najbolj odvisna (velikost vzorca primerljivih vozil,
        razpršenost ponudbe, segment vozila)?
  \item \textbf{RV3 --- hitrost in razširljivost.} Ali sistem zadošča
        nefunkcionalnim zahtevam glede odzivnosti tudi pod sočasno
        obremenitvijo, in kje je zgornja meja trenutne postavitve?
  \item \textbf{RV4 --- uporabnost.} Ali ciljni uporabniki znajo brez navodil
        opraviti glavne naloge (registracija, objava oglasa, vrednotenje)?
\end{itemize}
% =====================================================================
\section{Metodologija in merilno okolje}
\label{sec:metodologija}

\paragraph{Merilno okolje.}
Meritve hitrosti (razdelka~\ref{sec:latenca} in~\ref{sec:obremenitev}) so
izvedene v \emph{produkcijskem} okolju, torej na navideznem zasebnem strežniku
z dvema procesorskima jedroma in $4$~GB pomnilnika ter s tehnologijami in
njihovimi različicami, opisanimi v poglavju~\ref{ch:razvoj}
(razdelek~\ref{sec:namestitev}). Analitične in merilne skripte tečejo nad isto
podatkovno bazo kot je na produkciji.

\paragraph{Posnetek podatkov.}
Meritve v tem poglavju temeljijo na posnetku podatkovne baze iz septembra
$2026$, ki obsega približno $86.000$ normaliziranih oglasov
(razdelek~\ref{sec:podatki}). Ker se posnetek med analizo ne spreminja, so
rezultati ponovljivi.

\paragraph{Orodja.}
Za meritve in analizo v tem poglavju smo napisali namenske skripte, ločene od
produkcijske kode; ena od njih je dokumentirana kopija ocenjevalnega
algoritma, parametrizirana za eksperimente. Statistično obdelavo in izris
grafov smo izvedli s knjižnicama \texttt{NumPy}~\cite{harris2020numpy} in
\texttt{Matplotlib}~\cite{hunter2007matplotlib}.

\paragraph{Natančnost.}
Natančnost ocene (razdelek~\ref{sec:natancnost}) merimo z naključno izbranim
normaliziranim oglasom: njegove atribute (znamka, model, letnik, kilometrina)
uporabimo kot vhod v algoritem, sam oglas pa izločimo iz primerjalne množice. Nato izračunamo oceno in jo
primerjamo z \emph{dejansko oglaševano} ceno tega oglasa. Postopek ponovimo za
$3000$ naključnih oglasov.

% =====================================================================
\section{Uporabljene mere}
\label{sec:mere}

\paragraph{Latenca}
 je čas, ki preteče od trenutka tik pred oddajo zahtevka do
trenutka, ko odjemalec prejme celoten odgovor. Glede na to, kolikšen del te
poti zajamemo, ločimo tri vrste:
\begin{itemize}
  \item[(a)] \emph{latenca obdelave na strežniku} --- čas izvajanja funkcije
        za vrednotenje, od začetka do konca, brez omrežja;
  \item[(b)] \emph{latenca poizvedbe SQL} --- čas izvajanja poizvedbe v podatkovni bazi,
   od začetka do konca, brez omrežja;
  \item[(c)] \emph{latenca od konca do konca} (angl.\ \emph{end-to-end}) ---
        (a) skupaj z omrežnim obhodom med odjemalcem in strežnikom ter
        obdelavo zahtevka HTTP; to je čas, ki ga zazna uporabnik.
\end{itemize}
V razdelku~\ref{sec:latenca} merimo vrsti (a) in (b), v
razdelku~\ref{sec:obremenitev} pa vrsto (c).

\paragraph{Percentili.}
\textbf{Percentil} $p_k$ je vrednost, pod katero pade $k\,\%$ meritev. Poročamo
$p_{50}$ (mediana), $p_{95}$ in $p_{99}$.

\paragraph{Prepustnost, sočasnost, nasičenje.}
Obnašanje sistema pod obremenitvijo opišemo s tremi količinami:
\begin{itemize}
  \item \textbf{prepustnost} --- število uspešno obdelanih zahtevkov na
        sekundo;
  \item \textbf{sočasnost} --- število odjemalcev, ki hkrati držijo odprt
        zahtevek;
  \item \textbf{nasičenje} (\emph{koleno}) --- točka, v kateri dodatna
        sočasnost ne poveča več prepustnosti, latenca pa začne naraščati
        približno linearno.
\end{itemize}
Notranjo skladnost meritev preverimo z \textbf{Littlovim zakonom}~\cite{little1961},
\begin{equation}
  \label{eq:little}
  L = \lambda \cdot W,
\end{equation}
kjer je $L$ povprečno število zahtevkov v sistemu, $\lambda$ prepustnost in
$W$ povprečna latenca.

\paragraph{Napaka ocene.}
Za oglas $i$ z dejansko ceno $p_i$ in ocenjeno ceno $\hat p_i$ definiramo
absolutno odstotno napako
\begin{equation}
  \mathrm{APE}_i = \frac{\lvert \hat p_i - p_i \rvert}{p_i}.
\end{equation}
\textbf{MAPE} je povprečje $\mathrm{APE}_i$, \textbf{MdAPE} pa njihova mediana;
MdAPE je robustnejša na osamelce in je naša glavna mera~\cite{hyndman2006}.
Poročamo tudi \textbf{MAE} (povprečna absolutna napaka) in \textbf{RMSE}
(koren povprečne kvadratne napake) v evrih ter \emph{predznačeno} odstotno
napako, ki pokaže sistematičen odmik ocene navzgor ali navzdol.

\paragraph{Pokritost.}
\textbf{Pokritost} je delež poizvedb, za katere sistem vrne oceno, torej za
katere najde vsaj eno (oziroma vsaj $5$ oz.\ $10$) primerljivo vozilo.

\paragraph{Raznolikost ponudbe.}
\textbf{Raznolikost} (v izvorni kodi \texttt{ratio}) je relativna širina
cenovnega razpona primerljivih vozil,
\begin{equation}
  \label{eq:raznolikost}
  \mathrm{raznolikost} = \frac{P_{75} - P_{25}}{\hat{p}},
\end{equation}
kjer sta $P_{25}$ in $P_{75}$ (v kodi \texttt{goodPrice} in \texttt{badPrice})
uteženi $25.$ oziroma $75.$ percentil cen primerljivih vozil, $\hat{p}$
(v kodi \texttt{estimatedPrice}) pa ocenjena cena vozila. Vrednost izražamo v
odstotkih; enoten zapis uporabljamo v vseh tabelah tega poglavja.

% =====================================================================
\section{Značilnosti zbranih podatkov}
\label{sec:podatki}
\paragraph{Viri in obseg.}
Velika večina oglasov je bila zajeta z največjega slovenskega spletnega portala
za prodajo rabljenih vozil (razdelek~\ref{sec:zajem-podatkov}); manjši del
predstavlja vzorec z večjega nemškega oglasnega portala, ki smo ga uporabili za
preverjanje delovanja cevovoda na drugačni strukturi vira. Zajem je zajel
javno dostopne oglase za nekomercialno raziskovalno rabo, z omejeno hitrostjo
zahtevkov in brez shranjevanja kontaktnih ali osebnih podatkov prodajalcev.
Po normalizaciji baza obsega $85.846$ oglasov ($83.493$ s slovenskega in
$2.353$ z nemškega portala) za $10$ znamk vozil; razčlenitev po znamki je v
tabeli~\ref{tab:podatki-znamke}.

\begin{table}[htbp]
  \centering
  \caption{Število normaliziranih oglasov po znamki.}
  \label{tab:podatki-znamke}
  % vir: analysis/out/dataset_per_brand.csv
  \begin{tabular}{lr lr}
    \toprule
    \textbf{Znamka} & \textbf{Oglasi} & \textbf{Znamka} & \textbf{Oglasi} \\
    \midrule
    Volkswagen    & 21.191 & Peugeot     & 7.513 \\
    Audi          & 11.842 & Citroën     & 6.729 \\
    BMW           & 10.724 & Ford        & 5.783 \\
    Renault       & 9.615  & Toyota      & 2.648 \\
    Mercedes-Benz & 9.355  & Alfa Romeo  & 446 \\
    \bottomrule
  \end{tabular}
\end{table}

\paragraph{Porazdelitve ključnih atributov.}
Slika~\ref{fig:porazdelitve} prikazuje porazdelitve oglaševanih cen, prevožene
kilometrine in letnika. Porazdelitev cen je izrazito desno asimetrična (dolg
rep dragih vozil), zato jo prikazujemo na logaritemski osi. Mediana cene je
$14.890$~EUR (povprečje $18.137$~EUR), mediana kilometrine $159.950$~km,
mediana letnika pa $2017$; podrobnosti so v tabeli~\ref{tab:podatki-porazdelitve}.
Zaradi berljivosti so na grafih odrezane skrajne vrednosti (nekaj oglasov ima
očitno napačno ceno, kilometrino ali letnik); njihovo število je označeno v
kotu vsakega grafa, celoten razpon pa je razviden iz stolpca \emph{maks.}\ v
tabeli~\ref{tab:podatki-porazdelitve} in iz obravnave osamelcev v
tabeli~\ref{tab:osamelci}.

\begin{table}[htbp]
  \centering
  \caption{Povzetek porazdelitev ključnih numeričnih atributov.}
  \label{tab:podatki-porazdelitve}
  % vir: analysis/out/dataset_distribution_summary.json
  \setlength{\tabcolsep}{5pt}
  \begin{tabular}{lrrrrr}
    \toprule
    & \textbf{$p_{5}$} & \textbf{mediana} & \textbf{povprečje} & \textbf{$p_{95}$} & \textbf{maks.} \\
    \midrule
    cena [EUR]        & 1.550  & 14.890  & 18.137  & 49.014  & 374.000 \\
    kilometrina [km]  & 26.000 & 159.950 & 169.464 & 334.000 & 1.974.001 \\
    letnik            & 2005   & 2017    & 2016    & 2024    & 2029 \\
    moč [kW]          & 56     & 110     & 117     & 221     & 920 \\
    \bottomrule
  \end{tabular}
\end{table}

\begin{figure}[htbp]
  \centering
  \begin{subfigure}{0.48\textwidth}
    \includegraphics[width=\textwidth]{dataset_price_hist.pdf}
    \caption{oglaševana cena (log)}
  \end{subfigure}\hfill
  \begin{subfigure}{0.48\textwidth}
    \includegraphics[width=\textwidth]{dataset_km_hist.pdf}
    \caption{prevoženi kilometri}
  \end{subfigure}
  \\[1ex]
  \begin{subfigure}{0.48\textwidth}
    \includegraphics[width=\textwidth]{dataset_year_hist.pdf}
    \caption{letnik prve registracije}
  \end{subfigure}\hfill
  \begin{subfigure}{0.48\textwidth}
    \includegraphics[width=\textwidth]{dataset_brand_bar.pdf}
    \caption{oglasi po znamki}
  \end{subfigure}
  \caption{Porazdelitve ključnih atributov v normalizirani bazi.}
  \label{fig:porazdelitve}
  % vir: analysis/out/dataset_*_hist.pdf, dataset_brand_bar.pdf
\end{figure}

\paragraph{Pokritost po segmentih.}
Za pošteno oceno mora imeti algoritem dovolj primerljivih vozil v ožjem
segmentu (ista znamka in model, podoben letnik in kilometrina).
Slika~\ref{fig:pokritost} prikazuje toplotno karto za znamki Renault in
Volkswagen: vrstice so modeli, stolpci letniki, barva pa
število oglasov v tej kombinaciji. Vidno je, da so pogosti modeli srednjih
letnikov (npr.\ Volkswagen Golf ali Renault Clio) gosto pokriti, redki modeli
in skrajni letniki pa le redko, kar neposredno pojasni razlike v zanesljivosti
ocen (razdelek~\ref{sec:natancnost}).

\begin{figure}[htbp]
  \centering
  \includegraphics[width=\textwidth]{dataset_coverage_heatmap.pdf}
  \caption{Pokritost: število oglasov po modelu in letniku za znamki
  Renault in Volkswagen (po ena znamka na vrstico).}
  \label{fig:pokritost}
  % vir: analysis/out/dataset_coverage_heatmap.pdf
\end{figure}

\paragraph{Osamelci in kakovost podatkov.}
Postopek normalizacije trenutno ne izvaja filtriranja osamelcev, zato baza
vsebuje tudi vozila z očitno napačnimi vrednostmi.
Tabela~\ref{tab:osamelci} prikazuje njihov delež. Kar $10{,}7\,\%$ oglasov ima
vsaj eno vrednost zunaj smiselnih mej, največ pri kilometrini ($9{,}9\,\%$;
najvišja zabeležena vrednost je skoraj dva milijona kilometrov) in letniku
($8{,}1\,\%$; nekaj oglasov ima celo letnik v prihodnosti). Preprost filter po
interkvartilnem razponu (IQR) na ceni znotraj skupine
$(\text{znamka},\text{model})$ bi odstranil dodatnih $2{,}2\,\%$ oglasov.
Vpliv doslednega filtriranja osamelcev na natančnost obravnavamo v
razdelku~\ref{sec:obcutljivost}, uvedbo takega filtra pa v
poglavju~\ref{ch:sklep}.

\begin{table}[H]
  \centering
  \caption{Osamelci v normalizirani bazi ($85.846$ oglasov).}
  \label{tab:osamelci}
  % vir: analysis/out/dataset_outliers.csv
  \begin{tabularx}{\textwidth}{@{}X r r@{}}
    \toprule
    \textbf{Merilo} & \textbf{Št.\ vrstic} & \textbf{Delež} \\
    \midrule
    cena zunaj $[500,\ 200.000]$~EUR         & 427   & $0{,}5\,\%$ \\
    kilometrina zunaj $[1000,\ 500.000]$~km  & 8.495 & $9{,}9\,\%$ \\
    letnik zunaj $[1990,\ 2027]$             & 6.979 & $8{,}1\,\%$ \\
    vsaj ena vrednost zunaj mej              & 9.169 & $10{,}7\,\%$ \\
    dodatno po filtru IQR na ceni            & 1.878 & $2{,}2\,\%$ \\
    \bottomrule
  \end{tabularx}
\end{table}

% =====================================================================
\section{Natančnost ocene vrednosti}
\label{sec:natancnost}

To je osrednja meritev diplomskega dela: neposredno preverja zastavljeno
izhodišče o \emph{dovolj natančni} oceni (RV1). Protokol je opisan v
razdelku~\ref{sec:metodologija}, rezultate pa izdela skripta
\texttt{analysis/accuracy\_cv.py} nad vzorcem $3000$ oglasov (od tega jih je
$2986$ dobilo oceno).

\paragraph{Merilo resnice in namerni odmik ocene.}
Kot merilo resnice uporabljamo \emph{oglaševano} ceno, saj dejanske prodajne
cene niso javno dostopne. Sistem pa namenoma ne ocenjuje oglaševane cene,
temveč \emph{pričakovano prodajno ceno}: ker se rabljena vozila v pogajanjih
praviloma prodajo pod oglaševano ceno, oceno zniža s stopenjskim popustom
($9$--$15\,\%$ glede na cenovni razred). Zato je pri primerjavi z oglaševano
ceno pričakovan rahel \emph{negativni} odmik ocene, kar je zaželeno in ne
pomeni slabše ocene. To lastnost upoštevamo pri branju vseh rezultatov v
nadaljevanju: napako merimo dvakrat --- proti oglaševani ceni z vključenim
popustom (kar dejansko vrne sistem) in brez popusta (čista primerjalna
ocena).

\paragraph{Skupni rezultat.}
Brez popusta znaša MdAPE predlaganega pristopa $13{,}5\,\%$, MAPE $23{,}8\,\%$,
MAE $2798$~EUR in RMSE $5151$~EUR; z vključenim namernim popustom je MdAPE
proti oglaševani ceni $14{,}5\,\%$ (mediana predznačene napake $-6{,}8\,\%$).
Pokritost je $99{,}5\,\%$ pri $n \ge 1$ in $95{,}7\,\%$ pri $n \ge 10$.
Slika~\ref{fig:pred-vs-actual} prikazuje napovedano ceno proti dejanski
($R^2 = 0{,}87$), slika~\ref{fig:ape-hist} pa porazdelitev odstotne napake.

\begin{figure}[htbp]
  \centering
  \begin{subfigure}{0.48\textwidth}
    \includegraphics[width=\textwidth]{accuracy_pred_vs_actual.pdf}
    \caption{napovedano proti dejanskemu}
    \label{fig:pred-vs-actual}
  \end{subfigure}\hfill
  \begin{subfigure}{0.48\textwidth}
    \includegraphics[width=\textwidth]{accuracy_ape_hist.pdf}
    \caption{porazdelitev napake}
    \label{fig:ape-hist}
  \end{subfigure}
  \caption{Natančnost ocene vrednosti (leave-one-out, brez popusta).}
  % vir: analysis/out/accuracy_pred_vs_actual.pdf, accuracy_ape_hist.pdf
\end{figure}

\paragraph{Primerjava z osnovnimi modeli.}
Da presodimo, ali je predlagani pristop boljši od trivialnega, ga primerjamo s
štirimi referenčnimi modeli na istem protokolu in isti množici kandidatov
(tabela~\ref{tab:baselines}, slika~\ref{fig:baselines}): B0~globalna mediana
cene za dani model, B1~navadno povprečje cen v filtrirnem oknu, B1m~mediana cen
v oknu, B2~linearna regresija $\text{cena} \sim \text{letnik} + \text{km}$,
B3~metoda $k$~najbližjih sosedov ($k = 10$) po standardiziranih letniku in
kilometrini, ter B4~naš pristop (uteženo povprečje s podobnostnim jedrom, brez
popusta).

\begin{table}[htbp]
  \centering
  \caption{Primerjava natančnosti z osnovnimi modeli (leave-one-out).
  MdAPE in MAPE sta v odstotkih, MAE in RMSE v evrih.}
  \label{tab:baselines}
  % vir: analysis/out/accuracy_overall.csv
  \setlength{\tabcolsep}{5pt}
  \begin{tabular}{l r r r r}
    \toprule
    \textbf{Model} & \textbf{MdAPE} & \textbf{MAPE} & \textbf{MAE} & \textbf{RMSE} \\
    \midrule
    B0 globalna mediana       & 37{,}8 & 95{,}7 & 7395 & 10883 \\
    B1 povprečje v oknu       & 14{,}2 & 24{,}9 & 2930 & 5333 \\
    B1m mediana v oknu        & 13{,}1 & 22{,}5 & 2877 & 5561 \\
    B2 linearna regresija     & 12{,}9 & 22{,}4 & 2677 & 5032 \\
    B3 $k$-najbližjih sosedov & \textbf{11{,}2} & 21{,}5 & 2497 & 4883 \\
    B4 naš pristop            & 13{,}5 & 23{,}8 & 2798 & 5151 \\
    \bottomrule
  \end{tabular}
\end{table}

\begin{figure}[htbp]
  \centering
  \includegraphics[width=0.7\textwidth]{accuracy_baselines.pdf}
  \caption{MdAPE osnovnih modelov in predlaganega pristopa.}
  \label{fig:baselines}
  % vir: analysis/out/accuracy_baselines.pdf
\end{figure}

Predlagani pristop je \emph{bistveno} boljši od naivne globalne mediane
($13{,}5\,\%$ proti $37{,}8\,\%$) in primerljiv s preostalimi metodami, ki se
gibljejo med $11$ in $14\,\%$. Ni pa najboljši: metoda $k$~najbližjih sosedov
doseže $11{,}2\,\%$, tudi navadna mediana istega okna ($13{,}1\,\%$) je
malenkost natančnejša. To je pričakovano za zasnovo, katere cilj je
\emph{razložljivost}: za razliko od naštetih metod predlagani pristop poleg
točkovne ocene vrne tudi cenovni razpon in indeks zaupanja, uteži pa so
neposredno interpretabilne. Rezultat hkrati potrjuje smiselnost prehoda na
metode strojnega učenja, ki jo predlagamo v poglavju~\ref{ch:sklep}: če že
preprost $k$-NN preseže sedanji pristop, je od zmogljivejših modelov mogoče
pričakovati še boljše rezultate.

\paragraph{Odvisnost od velikosti vzorca in segmenta (RV2).}
Slika~\ref{fig:mdape-n} prikazuje MdAPE glede na število primerljivih vozil
$n$: napaka je pri $n < 10$ okoli $20\,\%$, pri $n \ge 20$ pa se ustali pri
približno $13\,\%$. To empirično podpira izbiro konstante $20$ v indeksu
zaupanja (enačba~\ref{eq:confidence}) --- pri toliko primerljivih vozilih je
ocena že blizu svoje najboljše natančnosti. Slika~\ref{fig:mdape-brand} prikazuje MdAPE po znamki: najbolje se
odrežejo pogosti, gosto pokriti modeli (Audi $9{,}8\,\%$, Toyota $10{,}8\,\%$,
Volkswagen $12{,}1\,\%$), najslabše pa Mercedes-Benz ($24{,}9\,\%$), pri
katerem se pri normalizaciji več različnih modelov zlije v isto kanonično
ime, kar poslabša izbor primerljivih vozil.
Tabela~\ref{tab:natancnost-segmenti} kaže še razčlenitev po starosti,
kilometrini in raznolikosti ponudbe.

\begin{figure}[htbp]
  \centering
  \begin{subfigure}{0.48\textwidth}
    \includegraphics[width=\textwidth]{accuracy_mdape_vs_n.pdf}
    \caption{glede na velikost vzorca $n$}
    \label{fig:mdape-n}
  \end{subfigure}\hfill
  \begin{subfigure}{0.48\textwidth}
    \includegraphics[width=\textwidth]{accuracy_mdape_by_brand.pdf}
    \caption{po znamki}
    \label{fig:mdape-brand}
  \end{subfigure}
  \caption{Odvisnost natančnosti ocene od podatkov.}
  % vir: analysis/out/accuracy_mdape_vs_n.pdf, accuracy_mdape_by_brand.pdf
\end{figure}

\begin{table}[htbp]
  \centering
  \caption{Natančnost predlaganega pristopa po segmentih (MdAPE).}
  \label{tab:natancnost-segmenti}
  % vir: analysis/out/accuracy_by_age.csv, accuracy_by_km.csv, accuracy_by_ratio.csv
  \begin{tabularx}{\textwidth}{@{}l X r r@{}}
    \toprule
    \textbf{Razčlenitev} & \textbf{Koš} & \textbf{Primeri} & \textbf{MdAPE [\%]} \\
    \midrule
    starost [leta] & 0--3   & 367 & 9{,}0 \\
                   & 4--7   & 921 & 10{,}6 \\
                   & 8--12  & 932 & 14{,}5 \\
                   & $13+$  & 766 & 22{,}9 \\
    \midrule
    kilometrina    & $<50.000$        & 280  & 9{,}6 \\
                   & $50.000$--$100.000$  & 449  & 11{,}9 \\
                   & $100.000$--$150.000$ & 634  & 12{,}4 \\
                   & $150.000$--$200.000$ & 629  & 13{,}7 \\
                   & $>200.000$       & 1021 & 16{,}7 \\
    \midrule
    raznolikost    & $0$--$0{,}10$    & 88   & 7{,}5 \\
                   & $0{,}10$--$0{,}25$ & 1002 & 8{,}4 \\
                   & $0{,}25$--$0{,}50$ & 1613 & 15{,}5 \\
                   & $>0{,}50$        & 447  & 29{,}7 \\
    \bottomrule
  \end{tabularx}
\end{table}

Vzorec pokaže jasne zakonitosti: napaka narašča s starostjo vozila ($9\,\%$ pri
najnovejših, $23\,\%$ pri vozilih, starejših od $13$~let), s prevoženo
kilometrino in --- najmočneje --- z \emph{raznolikostjo ponudbe}. Pri modelih z
enotno ponudbo (raznolikost pod $0{,}1$) je MdAPE le $7{,}5\,\%$, pri razpršeni
ponudbi (nad $0{,}5$) pa $29{,}7\,\%$. To potrjuje, da je raznolikost, ki jo
sistem uporabniku tudi prikaže, dober vgrajen pokazatelj zanesljivosti
posamezne ocene.

\paragraph{Vpliv popusta.}
Slika~\ref{fig:popust} prikazuje predznačeno napako z in brez popusta. Brez
popusta je ocena rahlo optimistična (mediana $+3{,}8\,\%$ nad oglaševano ceno),
z namernim popustom pa se premakne na $-6{,}8\,\%$ pod oglaševano ceno. Ta
premik je skladen z zasnovo sistema, ki cilja pričakovano prodajno in ne
oglaševane cene; dokončno umerjanje popusta bi bilo mogoče šele s podatki o
dejanskih prodajah (razdelek~\ref{sec:groznje}).

\begin{figure}[htbp]
  \centering
  \includegraphics[width=0.7\textwidth]{accuracy_discount_bias.pdf}
  \caption{Predznačena napaka proti oglaševani ceni, z in brez namernega popusta.}
  \label{fig:popust}
  % vir: analysis/out/accuracy_discount_bias.pdf
\end{figure}

\paragraph{Časovni razrez.}
Pri napovedovanju najnovejših $10\,\%$ oglasov zgolj iz starejših $90\,\%$
znaša MdAPE predlaganega pristopa $13{,}0\,\%$ (tabela~\ref{tab:temporal}), kar
je celo malenkost bolje od rezultata leave-one-out in potrjuje, da ocena ne
izhaja iz uhajanja informacije.

\begin{table}[htbp]
  \centering
  \caption{Natančnost pri časovnem razrezu (napoved najnovejših $10\,\%$).
  MdAPE in MAPE sta v odstotkih, MAE in RMSE v evrih.}
  \label{tab:temporal}
  % vir: analysis/out/accuracy_temporal.csv
  \setlength{\tabcolsep}{5pt}
  \begin{tabular}{l r r r r}
    \toprule
    \textbf{Model} & \textbf{MdAPE} & \textbf{MAPE} & \textbf{MAE} & \textbf{RMSE} \\
    \midrule
    B0 globalna mediana       & 35{,}9 & 74{,}5 & 7596 & 11222 \\
    B1m mediana v oknu        & 12{,}8 & 19{,}7 & 3088 & 5805 \\
    B3 $k$-najbližjih sosedov & 10{,}5 & 18{,}2 & 2667 & 4890 \\
    B4 naš pristop            & 13{,}0 & 20{,}6 & 2994 & 5347 \\
    \bottomrule
  \end{tabular}
\end{table}

% =====================================================================
\section{Občutljivost algoritma na parametre}
\label{sec:obcutljivost}

Ocenjevalni algoritem vsebuje več hevristično določenih konstant: dovoljeno
odstopanje kilometrine ($\pm 40.000$~km) in letnika ($\pm 3$~leta), uteži
delnih ocen ($0{,}6$ in $0{,}4$), oblika padajoče funkcije in stopenjski
popust. V tem razdelku vsako spreminjamo ob nespremenjenih ostalih in merimo
MdAPE ter pokritost (vzorec $2000$ oglasov, leave-one-out). Rezultate izdela
skripta \texttt{analysis/sensitivity.py}.

\paragraph{Okno primerljivosti.}
Tabela~\ref{tab:sens-okno} in slika~\ref{fig:sens-km} kažeta, da je algoritem
na širino \emph{kilometrinskega} okna skoraj neobčutljiv: MdAPE se med
$\pm 10.000$~km in $\pm 80.000$~km giblje le med $13{,}2$ in $13{,}6\,\%$.
Nasprotno je odvisen od \emph{letnikovega} okna: pri $\pm 1$~letu je MdAPE
$12{,}6\,\%$, pri $\pm 5$~letih pa $15{,}2\,\%$. Ožje okno da bolj primerljiva
vozila, a manjši vzorec (nižja pokritost); gre za znano razmerje med
pristranskostjo in varianco~\cite{bias-variance}. Toplotna karta
(slika~\ref{fig:sens-heatmap}) potrjuje, da natančnost določa skoraj izključno
letnikovo okno. Sedanja izbira $\pm 3$~leta je razumen kompromis med
natančnostjo in pokritostjo; morebitna zožitev na $\pm 2$~leti bi MdAPE
znižala za približno $0{,}3$~odstotne točke ob še vedno visoki pokritosti.

\begin{table}[htbp]
  \centering
  \caption{Občutljivost na širino okna primerljivosti. Pokritost je delež
  poizvedb z vsaj $5$ primerljivimi vozili.}
  \label{tab:sens-okno}
  % vir: analysis/out/sensitivity_km.csv, sensitivity_year.csv
  \setlength{\tabcolsep}{5pt}
  \begin{tabular}{l r r r}
    \toprule
    \textbf{Parameter} & \textbf{MdAPE [\%]} & \textbf{Pokritost [\%]} & \textbf{Mediana $n$} \\
    \midrule
    kilometrina $\pm 10.000$~km & 13{,}2 & 92{,}3 & 53 \\
    kilometrina $\pm 40.000$~km & 13{,}4 & 98{,}1 & 196 \\
    kilometrina $\pm 80.000$~km & 13{,}6 & 99{,}0 & 362 \\
    \midrule
    letnik $\pm 1$~leto  & 12{,}6 & 95{,}6 & 97 \\
    letnik $\pm 3$~leta  & 13{,}4 & 98{,}1 & 196 \\
    letnik $\pm 5$~let   & 15{,}2 & 98{,}5 & 262 \\
    \bottomrule
  \end{tabular}
\end{table}

\begin{figure}[htbp]
  \centering
  \begin{subfigure}{0.48\textwidth}
    \includegraphics[width=\textwidth]{sensitivity_km.pdf}
    \caption{okno kilometrine}
  \end{subfigure}\hfill
  \begin{subfigure}{0.48\textwidth}
    \includegraphics[width=\textwidth]{sensitivity_year.pdf}
    \caption{okno letnika}
  \end{subfigure}
  \caption{Občutljivost na širino okna primerljivosti (MdAPE in pokritost).}
  \label{fig:sens-km}
  % vir: analysis/out/sensitivity_km.pdf, sensitivity_year.pdf
\end{figure}

\begin{figure}[htbp]
  \centering
  \includegraphics[width=0.7\textwidth]{sensitivity_km_year_heatmap.pdf}
  \caption{MdAPE [\%] po mreži dovoljenih odstopanj kilometrine in letnika.}
  \label{fig:sens-heatmap}
  % vir: analysis/out/sensitivity_km_year_heatmap.pdf
\end{figure}

\paragraph{Uteži, oblika jedra in popust.}
Tabela~\ref{tab:sens-ostalo} povzema preostale parametre. Razmerje uteži med
kilometrino in letnikom na natančnost praktično ne vpliva ($13{,}4$ do
$13{,}6\,\%$ za vsa preizkušena razmerja), prav tako ne zamenjava linearne
padajoče funkcije z eksponentno ($13{,}4$ proti $13{,}3\,\%$). Edina parametra,
ki opazno premakneta rezultat, sta torej širina letnikovega okna in namerni
popust: ta predznačeno napako proti oglaševani ceni premakne iz $+3{,}8\,\%$ v
$-6{,}8\,\%$, MdAPE pa zviša z $13{,}4$ na $14{,}5\,\%$. Ta premik je posledica
zavestne odločitve, da sistem ocenjuje pričakovano prodajno in ne oglaševane
cene, in ne slabše ocene.

\begin{table}[htbp]
  \centering
  \caption{Občutljivost na uteži, obliko jedra in popust.}
  \label{tab:sens-ostalo}
  % vir: analysis/out/sensitivity_weights.csv, sensitivity_kernel.csv, sensitivity_discount.csv
  \begin{tabularx}{\textwidth}{@{}l X r >{\raggedleft\arraybackslash}p{2.3cm}@{}}
    \toprule
    \textbf{Parameter} & \textbf{Nastavitev} & \textbf{MdAPE [\%]} & \textbf{Predznačena napaka [\%]} \\
    \midrule
    uteži (km / letnik) & $0{,}5 / 0{,}5$              & 13{,}4 & $+3{,}7$ \\
    uteži (km / letnik) & $0{,}6 / 0{,}4$ (privzeto)   & 13{,}4 & $+3{,}8$ \\
    uteži (km / letnik) & $0{,}8 / 0{,}2$              & 13{,}6 & $+3{,}8$ \\
    \midrule
    jedro               & linearno (privzeto)         & 13{,}4 & $+3{,}8$ \\
    jedro               & eksponentno                 & 13{,}3 & $+3{,}8$ \\
    \midrule
    popust              & brez                        & 13{,}4 & $+3{,}8$ \\
    popust              & stopenjski (privzeto)        & 14{,}5 & $-6{,}8$ \\
    \bottomrule
  \end{tabularx}
\end{table}

\paragraph{Sklep o robustnosti.}
Algoritem je robusten: na večino svojih konstant se skoraj ne odziva, kar
pomeni, da njihova hevristična izbira ni sporna. Njegova natančnost je odvisna
predvsem od kakovosti in gostote podatkov (razdelek~\ref{sec:natancnost}), ne
od finega nastavljanja parametrov.

% =====================================================================
\section{Latenca posameznega vrednotenja}
\label{sec:latenca}

Merimo latenco vrste (a) in (b) iz razdelka~\ref{sec:mere}. Meritev izvede
\texttt{benchmarks/bench\_algo.py}: $2000$ neposrednih klicev funkcije nad
naključnimi realnimi vhodi (brez HTTP), po $100$ ogrevalnih klicih, ki se ne
merijo. Ločeno merimo čas poizvedbe SQL, čas točkovanja v Pythonu in čas
serializacije odgovora.

\paragraph{Ozko grlo: pregled celotne tabele.}
Prvotna izvedba poizvedbe ni imela ustreznega indeksa.
Izsek~\ref{lst:explain-before} prikazuje izpis \texttt{EXPLAIN (ANALYZE,
BUFFERS)}: načrt je \texttt{Seq Scan}, baza pri vsakem vrednotenju prebere
celotno tabelo in zavrže $85.273$ vrstic, ki ne ustrezajo pogojem
(\texttt{Rows Removed by Filter}). Ker so podatki v celoti v predpomnilniku
(\texttt{shared hit}, brez \texttt{read}), je čas stabilen, a raste linearno z
velikostjo tabele.

\begin{center}
\captionof{listing}{\texttt{EXPLAIN (ANALYZE, BUFFERS)} pred dodajanjem indeksa.}
\label{lst:explain-before}
\begin{minted}{text}
Seq Scan on normalized_market_listing
    (cost=0.00..3314.15 rows=12 width=97)
    (actual time=0.555..6.847 rows=573 loops=1)
  Filter: ((kilometers >= 18500) AND (kilometers <= 98500)
       AND (registration_year >= 2019) AND (registration_year <= 2025)
       AND (brand_name = 'Citroën') AND (model_name = 'C3'))
  Rows Removed by Filter: 85273
  Buffers: shared hit=1168
Execution Time: 6.880 ms
\end{minted}
\end{center}

\paragraph{Učinek sestavljenega indeksa.}
Po dodajanju sestavljenega indeksa na
$(\text{brand\_name},\ \text{model\_name},\ \text{registration\_year},\
\text{kilometers})$ se načrt spremeni v \texttt{Bitmap Index Scan}
(izsek~\ref{lst:explain-after}): baza prebere le nekaj strani indeksa in
pripadajoče vrstice, čas izvajanja pa pade z $6{,}88$~ms na $0{,}69$~ms.
Latenca poizvedbe SQL se v mediani zmanjša z $9{,}7$~ms na $2{,}0$~ms
(tabela~\ref{tab:indeks-pred-po}, slika~\ref{fig:indeks-pred-po}). Pomembneje
od trenutnega pohitrenja je, da postane čas poizvedbe neodvisen od velikosti
tabele: namesto $O(\text{tabela})$ postane $O(\text{zadetki})$.

\begin{table}[htbp]
  \centering
  \caption{Latenca poizvedbe SQL pred in po dodajanju indeksa, v milisekundah.}
  \label{tab:indeks-pred-po}
  % vir: benchmarks/report.md (results_algo_before/after.json)
  \begin{tabular}{lrrrr}
    \toprule
    & \textbf{p50} & \textbf{p95} & \textbf{p99} & \textbf{Povprečje} \\
    \midrule
    pred (Seq Scan)  & 9{,}7 & 21{,}0 & 24{,}8 & 11{,}3 \\
    po (Index Scan)  & 2{,}0 & 8{,}6  & 12{,}5 & 3{,}0 \\
    \bottomrule
  \end{tabular}
\end{table}

\begin{figure}[htbp]
  \centering
  \includegraphics[width=0.6\textwidth]{latency_index_before_after.pdf}
  \caption{Latenca poizvedbe SQL pred in po dodajanju sestavljenega indeksa.}
  \label{fig:indeks-pred-po}
  % vir: analysis/out/latency_index_before_after.pdf
\end{figure}

\begin{center}
\captionof{listing}{\texttt{EXPLAIN (ANALYZE, BUFFERS)} po dodajanju indeksa.}
\label{lst:explain-after}
\begin{minted}{text}
Bitmap Heap Scan on normalized_market_listing
    (cost=5.84..50.49 rows=12 width=97)
    (actual time=0.202..0.615 rows=573 loops=1)
  Recheck Cond: (brand_name = 'Citroën' AND model_name = 'C3'
       AND registration_year BETWEEN 2019 AND 2025
       AND kilometers BETWEEN 18500 AND 98500)
  Buffers: shared hit=195
  ->  Bitmap Index Scan on idx_nml_valuation
        (actual time=0.166..0.166 rows=573 loops=1)
        Buffers: shared hit=6
Execution Time: 0.688 ms
\end{minted}
\end{center}

\paragraph{Razčlenitev latence po komponentah.}
Tabela~\ref{tab:latenca-vrednotenja} in slika~\ref{fig:latency-breakdown}
prikazujeta latenco celotnega vrednotenja po dodajanju indeksa. Tipično
vrednotenje traja $2{,}0$~ms, v najslabšem odstotku primerov (p99) pa
$13{,}5$~ms. Skoraj celotna latenca ($\approx 88\,\%$ povprečja) odpade na
poizvedbo SQL; točkovanje v Pythonu ($0{,}4$~ms povprečno) in serializacija
odgovora ($\approx 0{,}01$~ms) sta zanemarljiva. Poizvedba v povprečju vrne
$287$ primerljivih vozil (mediana $175$, največ $1495$).

\begin{table}[htbp]
  \centering
  \caption{Latenca posameznega vrednotenja po dodajanju indeksa
  ($N = 2000$), v milisekundah.}
  \label{tab:latenca-vrednotenja}
  % vir: benchmarks/report.md (results_algo.json)
  \begin{tabular}{lrrrr}
    \toprule
    \textbf{Komponenta} & \textbf{p50} & \textbf{p95} & \textbf{p99} & \textbf{Povprečje} \\
    \midrule
    Skupaj                & 2{,}0  & 9{,}7 & 13{,}5 & 3{,}1 \\
    Poizvedba SQL         & 1{,}8  & 8{,}4 & 11{,}5 & 2{,}7 \\
    Točkovanje (Python)   & 0{,}2  & 1{,}2 & 1{,}8  & 0{,}4 \\
    Serializacija (JSON)  & 0{,}01 & 0{,}02 & 0{,}03 & 0{,}01 \\
    \bottomrule
  \end{tabular}
\end{table}

\begin{figure}[htbp]
  \centering
  \includegraphics[width=0.6\textwidth]{latency_breakdown.pdf}
  \caption{Razčlenitev latence posameznega vrednotenja po komponentah.}
  \label{fig:latency-breakdown}
  % vir: analysis/out/latency_breakdown.pdf
\end{figure}

Vrednotenje se torej izvede daleč pod nefunkcionalno zahtevo $2$~s
(razdelek~\ref{sec:zahteve}); prostora za rast podatkovne baze je ob
prisotnosti indeksa dovolj.

% =====================================================================
\section{Obnašanje pod obremenitvijo}
\label{sec:obremenitev}

Merimo end-to-end latenco (vrsta (c)) in prepustnost končne točke
\texttt{POST /valuate} pri naraščajoči sočasnosti. Obremenitev ustvarja
skripta \texttt{benchmarks/loadtest.py} z več vzporednimi odjemalci, ki v
fiksnem $20$-sekundnem oknu neprekinjeno pošiljajo zahtevke z naključnimi
realnimi vhodi. Test se samodejno prekine ob deležu napak nad $5\,\%$ ali
$p_{95}$ nad $5$~s. Med testom vzorčimo tudi zasedenost baze
(\texttt{pg\_stat\_activity}).

\begin{table}[htbp]
  \centering
  \caption{Prepustnost in latenca pri naraščajoči sočasnosti.
  Latence so v milisekundah, $L$ je ocena po Littlovem zakonu.}
  \label{tab:obremenitev}
  % vir: benchmarks/results_load.json (loadtest.py --db-sample)
  \begin{tabular}{r r r r r r r}
    \toprule
    \textbf{Sočasnost} & \textbf{Prep.\ [z/s]} & \textbf{p50} & \textbf{p95} & \textbf{p99}
      & \textbf{Napake} & \textbf{$L$} \\
    \midrule
    1  & \res{x} & \res{x} & \res{x} & \res{x} & \res{x} & \res{x} \\
    5  & \res{x} & \res{x} & \res{x} & \res{x} & \res{x} & \res{x} \\
    10 & \res{x} & \res{x} & \res{x} & \res{x} & \res{x} & \res{x} \\
    20 & \res{x} & \res{x} & \res{x} & \res{x} & \res{x} & \res{x} \\
    40 & \res{x} & \res{x} & \res{x} & \res{x} & \res{x} & \res{x} \\
    \bottomrule
  \end{tabular}
\end{table}

\paragraph{Rezultat.}
\res{Prepustnost narašča približno linearno do sočasnosti X, nato se ustali pri
približno X zahtevkih na sekundo. Vse latence ostanejo pod nefunkcionalno
zahtevo 2~s.}

\paragraph{Skladnost meritve (Littlov zakon).}
\res{Stolpec $L$ (prepustnost krat povprečna latenca) ostaja blizu nastavljene
sočasnosti, kar potrjuje notranjo skladnost meritve.}

\paragraph{Dokaz ozkega grla.}
Teoretični strop postavitve je $1000 / S \approx 324$ vrednotenj na sekundo,
kjer je $S = 3{,}1$~ms povprečni čas ene evalvacije. Dejanska meja je nižja,
ker ves promet serializira \emph{ena sama deljena povezava do baze}
(\texttt{app/database.py}) in en sam strežniški proces.
\res{To potrjuje vzorčenje baze med testom: število aktivnih poizvedb kljub
naraščajoči sočasnosti odjemalcev ne preseže X.}

% =====================================================================
\section{Preverjanje funkcionalnih zahtev}
\label{sec:funkcionalne}

Tabela~\ref{tab:sledljivost} povezuje vsako funkcionalno zahtevo iz
razdelka~\ref{sec:zahteve} z načinom preverjanja in izidom.

\begin{table}[htbp]
  \centering
  \caption{Sledljivost funkcionalnih zahtev do preverjanja.}
  \label{tab:sledljivost}
  % vir: ročni scenariji + imena Jest testov v car-stock-market-be
  \begin{tabularx}{\textwidth}{@{}X X c@{}}
    \toprule
    \textbf{Zahteva} & \textbf{Način preverjanja} & \textbf{Izid} \\
    \midrule
    Registracija uporabnika              & \res{test / scenarij} & \res{OK} \\
    Prijava in odjava                    & \res{test / scenarij} & \res{OK} \\
    Objava oglasa (le registrirani)      & \res{test / scenarij} & \res{OK} \\
    Vrednotenje vozila po atributih      & \res{test / scenarij} & \res{OK} \\
    Všečkanje oglasov                    & \res{test / scenarij} & \res{OK} \\
    Ažuriranje lastnih/všečkanih oglasov & \res{test / scenarij} & \res{OK} \\
    Pregled oglasov (anonimni)           & \res{test / scenarij} & \res{OK} \\
    \bottomrule
  \end{tabularx}
\end{table}

\paragraph{Pokritost testov.}
Avtomatizirani testi zalednega sistema (Jest + Supertest,
razdelek~\ref{sec:testiranje-be}) obsegajo \res{X} testov v \res{X} zbirkah in
pokrivajo \res{X}\,\% stavkov ter \res{X}\,\% vej. Testi se izvedejo ob vsakem
zahtevku za združitev na cevovodu neprekinjene integracije
(razdelek~\ref{subsec:ci}).

% =====================================================================
\section{Preverjanje nefunkcionalnih zahtev}
\label{sec:nefunkcionalne}

Tabela~\ref{tab:nfr} za vsako nefunkcionalno zahtevo iz
razdelka~\ref{sec:zahteve} navaja merilo, izmerjeno vrednost in sodbo.

\begin{table}[htbp]
  \centering
  \caption{Preverjanje nefunkcionalnih zahtev.}
  \label{tab:nfr}
  \setlength{\tabcolsep}{4pt}
  \begin{tabularx}{\textwidth}{@{}>{\raggedright\arraybackslash}X
                                 >{\raggedright\arraybackslash}X
                                 >{\raggedright\arraybackslash}X c@{}}
    \toprule
    \textbf{Zahteva} & \textbf{Merilo} & \textbf{Izmerjeno} & \textbf{Sodba} \\
    \midrule
    Odziv vmesnika na dejanje uporabnika & $< 100$~ms & \res{x}~ms  & \res{OK} \\
    Zahtevki HTTP (razen vrednotenja)    & $< 300$~ms & \res{x}~ms ($p_{95}$) & \res{OK} \\
    Vrednotenje vozila                   & $< 2$~s    & $13{,}5$~ms ($p_{99}$, brez omrežja) & OK \\
    Vizualna privlačnost / dostopnost    & Lighthouse & \res{x} / \res{x} & \res{OK} \\
    Prilagodljiva postavitev             & preizkus na več velikostih zaslona & \res{da} & \res{OK} \\
    Varovanje uporabniških podatkov      & bcrypt, JWT z rokom & \res{opis} & \res{OK} \\
    Odzivnost pri zahtevnejših opravilih & asinhrona obdelava & \res{opis} & \res{OK} \\
    \bottomrule
  \end{tabularx}
\end{table}

\paragraph{Meritev odzivnosti vmesnika in endpointov.}
\res{Odzivnost uporabniškega vmesnika smo izmerili z orodjem Lighthouse na treh
ključnih straneh; ocene za performanco so X, za dostopnost pa X. Latenco
glavnih endpointov zalednega sistema smo izmerili z orodjem autocannon:
$p_{95}$ je v vseh primerih pod X~ms, torej pod zahtevo 300~ms.}

\paragraph{Varnost.}
Gesla so pred shranjevanjem zgoščena z algoritmom bcrypt
(razdelek~\ref{sec:avtentikacija}), žetoni JWT imajo rok veljavnosti $7$~dni,
poti vmesnika API pa ščiti sredinski preverjevalnik žetonov. \res{Orodje
\texttt{npm audit} ob zadnjem pregledu ni javilo kritičnih ranljivosti v
odvisnostih.}

% =====================================================================
\section{Uporabniški preizkus}
\label{sec:uporabniki}

Izvedli smo formativni uporabniški preizkus~\cite{nielsen1994} z namenom
odkriti težave uporabnosti, ne pa statistično ovrednotiti izdelka.

\paragraph{Udeleženci in postopek.}
V preizkus smo vključili \res{N} udeležencev iz kroga družine in bližnjih
prijateljev, katerih starost in računalniška pismenost sta se med seboj
razlikovali (tabela~\ref{tab:udelezenci}). Vsak udeleženec je prejel enak
seznam nalog: registracija, objava oglasa, všečkanje oglasa, izbira
slovenskega jezika v nastavitvah in vrednotenje poljubnega vozila. Med
izvajanjem je udeleženec razmišljal na glas, opazovalec pa je beležil, ali je
nalogo dokončal samostojno, in vsako opaženo težavo.

\begin{table}[htbp]
  \centering
  \caption{Udeleženci uporabniškega preizkusa.}
  \label{tab:udelezenci}
  \begin{tabular}{c c c l}
    \toprule
    \textbf{\#} & \textbf{Starost (razpon)} & \textbf{Rač.\ pismenost (1--5)} & \textbf{Naprava} \\
    \midrule
    1 & \res{x} & \res{x} & \res{x} \\
    2 & \res{x} & \res{x} & \res{x} \\
    3 & \res{x} & \res{x} & \res{x} \\
    \res{\dots} & & & \\
    \bottomrule
  \end{tabular}
\end{table}

\paragraph{Rezultati.}
Vsi udeleženci so vse naloge dokončali brez pomoči, kar potrjuje, da so glavne
poti v vmesniku dovolj jasne (RV4). Kljub temu so odkrili nekaj
pomanjkljivosti; tabela~\ref{tab:tezave} jih razvršča po resnosti (lestvica
$0$--$4$ po Nielsenu~\cite{nielsen1994}) in navaja, ali so bile odpravljene.

\begin{table}[htbp]
  \centering
  \caption{Težave, odkrite med uporabniškim preizkusom
  (resnost po lestvici $0$--$4$).}
  \label{tab:tezave}
  \setlength{\tabcolsep}{5pt}
  \begin{tabularx}{\textwidth}{@{}>{\raggedright\arraybackslash}X c c@{}}
    \toprule
    \textbf{Težava} & \textbf{Resnost} & \textbf{Status} \\
    \midrule
    Premajhni naslovi podstrani v glavi glavne strani       & \res{x} & odpravljeno \\
    Manjkajoče polje za ponovni vnos gesla pri registraciji & \res{x} & odpravljeno \\
    Velika števila brez ločil za tisočice                   & \res{x} & odpravljeno \\
    Manjkajoč opis pod slikami oblik vozil v obrazcu        & \res{x} & odpravljeno \\
    Premalo vozil, dostopnih za vrednotenje                 & \res{x} & \res{v načrtu} \\
    \bottomrule
  \end{tabularx}
\end{table}

Omejitve preizkusa (majhen in priložnostni vzorec, brez standardiziranega
vprašalnika, brez merjenja časa po nalogah) obravnavamo v
razdelku~\ref{sec:groznje}.

% =====================================================================
\section{Simulacija uporabe}
\label{sec:simulacija}

Za ponazoritev delovanja rešitve povzemamo tipičen potek uporabe. Uporabnik je
ob vstopu usmerjen na glavno stran (slika~\ref{fig:glavna-stran}), kjer sta
izpostavljeni prijava in registracija ter pregled novih in vseh objavljenih
oglasov. Registriran uporabnik lahko objavi oglas prek vodene forme
(slika~\ref{fig:objava-oglasa}), ovrednoti poljubno vozilo, všečka oglase in si
ogleda svoj profil z všečkanimi vozili (slika~\ref{fig:liked-cars}).

\begin{figure}[H]
  \centering
  \includegraphics[width=0.9\textwidth]{glavna-stran.png}
  \caption{Glavna stran prototipa}
  \label{fig:glavna-stran}
\end{figure}
\begin{figure}[H]
  \centering
  \includegraphics[width=0.9\textwidth]{objava-oglasa.png}
  \caption{Vodena forma za objavo oglasa}
  \label{fig:objava-oglasa}
\end{figure}
\begin{figure}[H]
  \centering
  \includegraphics[width=0.9\textwidth]{liked-cars.png}
  \caption{Stran uporabnika s prikazanimi všečkanimi vozili}
  \label{fig:liked-cars}
\end{figure}

\paragraph{Primerjava ocen z obstoječimi portali.}
Za ilustracijo smo delovanje sistema primerjali s ponudbo na dveh spletnih
portalih na treh vozilih (tabela~\ref{tab:vrednotenje}). Sistematično in
kvantitativno oceno natančnosti podaja razdelek~\ref{sec:natancnost}; ta
primerjava je zgolj kvalitativna ponazoritev.

\begin{table}[htbp]
  \centering
  \caption{Primerjava ocenjenih cen s ponudbo na spletnih portalih (v evrih).}
  \label{tab:vrednotenje}
  \setlength{\tabcolsep}{5pt}
  \begin{tabular}{lccc}
    \toprule
    & \textbf{Audi A3} & \textbf{BMW ser.\ 2} & \textbf{Citroën Berlingo} \\
    \midrule
    slov.\ portal (razpon)     & 4.000--6.000  & 9.000--12.000  & 7.500--10.000 \\
    nem.\ portal (razpon)      & 4.500--9.000  & pribl.\ 11.100 & 8.000--12.000 \\
    \midrule
    Sistem --- dobra cena      & 4.500  & 7.600   & 7.600  \\
    Sistem --- sredinska cena  & 6.350  & 9.300   & 9.941  \\
    Sistem --- slaba cena      & 7.600  & 10.500  & 12.100 \\
    Raznolikost                & $49\,\%$ & $31\,\%$ & $45\,\%$ \\
    \bottomrule
  \end{tabular}
\end{table}

Pri vozilih z nizko raznolikostjo ponudbe (Audi~A3, BMW serije~2) se sredinska
ocena dobro ujema s cenami na obeh portalih. Pri modelu Citroën Berlingo je
raznolikost bistveno višja, zato je razpon med dobro in slabo oceno širši,
sredinska ocena pa manj zanesljiva --- kar je skladno z ugotovitvijo iz
razdelka~\ref{sec:natancnost}, da napaka strmo narašča z raznolikostjo ponudbe.
Pri interpretaciji je treba upoštevati, da sistem ocenjuje pričakovano prodajno
in ne oglaševano ceno.

% =====================================================================
\section{Grožnje veljavnosti}
\label{sec:groznje}

\paragraph{Notranja veljavnost.}
Kot merilo resnice pri merjenju natančnosti uporabljamo \emph{oglaševano} ceno,
ne dejanske prodajne cene, ki ni javno dostopna. Namerni popust, s katerim
sistem cilja pričakovano prodajno ceno, je zato ocenjen in ne umerjen na
dejanskih prodajah. Isti fizični oglas se lahko pojavi pri več virih; prečno
preverjanje izloči zapise z istim izvornim oglasom, ne more pa zaznati
podvojitev med viri.

\paragraph{Zunanja veljavnost.}
Podatki pokrivajo $10$ znamk in pretežno slovenski trg (z manjšim nemškim
vzorcem); znamka z največ oglasi na trgu (okoli $8000$ oglasov) v šifrantu ni
zajeta in je iz analize izpuščena. Uporabniški preizkus je temeljil na majhnem
in priložnostnem vzorcu.

\paragraph{Konstruktna veljavnost.}
Pojem \enquote{tržna vrednost} operacionaliziramo kot uteženo povprečje
oglaševanih cen primerljivih vozil, znižano za pričakovani pogajalski popust;
\enquote{raznolikost ponudbe} pa kot relativno širino interkvartilnega razpona
teh cen. Drugačna operacionalizacija bi lahko dala drugačne rezultate.

\paragraph{Zanesljivost.}
Meritve latence in obremenitve so bile izvedene na deljenem produkcijskem
strežniku, kjer lahko na rezultate vpliva ozadje. Meritve natančnosti so
ponovljive nad fiksnim posnetkom baze.

% =====================================================================
\section{Povzetek ugotovitev}
\label{sec:povzetek-analize}

Tabela~\ref{tab:rv-povzetek} povzema odgovore na raziskovalna vprašanja glede
na vnaprej postavljena merila.

\begin{table}[htbp]
  \centering
  \caption{Povzetek odgovorov na raziskovalna vprašanja.}
  \label{tab:rv-povzetek}
  \begin{tabularx}{\textwidth}{@{}l X c@{}}
    \toprule
    \textbf{RV} & \textbf{Glavni rezultat} & \textbf{Sodba} \\
    \midrule
    RV1 natančnost & MdAPE $13{,}5\,\%$ ($9$--$10\,\%$ za dobro pokrita vozila),
                     bistveno bolje od naivne ocene ($37{,}8\,\%$), primerljivo z
                     drugimi metodami & izpolnjeno \\
    RV2 odvisnost  & napaka pada z velikostjo vzorca (ustali se pri $n\approx20$),
                     narašča s starostjo, kilometrino in raznolikostjo ponudbe
                     ($7{,}5\,\%$ do $29{,}7\,\%$) & pojasnjeno \\
    RV3 hitrost    & vrednotenje $2$~ms (p50), $13{,}5$~ms (p99); indeks odpravi
                     odvisnost od velikosti baze; \res{meja X~z/s} & \res{izpolnjeno} \\
    RV4 uporabnost & vse naloge dokončane brez pomoči; \res{X} težav odpravljenih & izpolnjeno \\
    \bottomrule
  \end{tabularx}
\end{table}

\paragraph{Sodba o zastavljenih izhodiščih.}
Rezultati \emph{potrjujejo} izhodišča, ki smo si jih zastavili na začetku. Sistem na podlagi štirih
atributov vozila vrne oceno vrednosti z mediano odstotne napake $13{,}5\,\%$
(brez namernega popusta), pri dobro pokritih vozilih pa $9$--$10\,\%$, kar je
znotraj zastavljenega praga $15\,\%$ in bistveno bolje od naivne metode.
Predlagani pristop se med primerljivo natančnimi metodami ne uvršča na prvo
mesto --- preprost $k$-najbližjih sosedov doseže $11{,}2\,\%$ --- je pa edini,
ki poleg točkovne ocene vrne tudi razložljiv cenovni razpon in indeks
zaupanja, ta pa se izkaže za dober pokazatelj dejanske zanesljivosti
posamezne ocene. Glavne omejitve so odvisnost od enega prevladujočega
podatkovnega vira, nižja natančnost pri starejših vozilih in modelih z
razpršeno ponudbo ter merilo resnice, ki temelji na oglaševanih in ne
prodajnih cenah. Konkretne smeri za izboljšave --- vključno s prehodom na
metode strojnega učenja, ki jih ta analiza dodatno utemeljuje --- navajamo v
poglavju~\ref{ch:sklep}.

%
%  Številke so vpisane iz izhodov skript v CarStockMarketEvaluation/analysis/.
%  Preostale oznake \res{...} (rdeče) je treba še izpolniti:
%    - obremenitveni test (benchmarks/loadtest.py)
%    - sledljivost testov, Jest coverage, Lighthouse
%    - tabeli udeležencev in najdenih težav
%  Slike (PDF) iz analysis/out/ prekopiraj v mapo slike/.
% =====================================================================

\newcommand{\res}[1]{\textcolor{red}{\textbf{[#1]}}}

\chapter{Analiza in vrednotenje rešitve}
\label{ch:analiza}

V tem poglavju sistematično pregleadamo in analiziramo razvito rešitev. Najprej iz zahtev in
izhodišč, ki smo si jih zastavili na začetku, izpeljemo raziskovalna vprašanja in za vsako določimo merilo
uspešnosti (razdelek~\ref{sec:rq}). Nato opišemo merilno okolje in postopke
(razdelek~\ref{sec:metodologija}) ter natančno definiramo vse uporabljene mere
(razdelek~\ref{sec:mere}). Sledijo rezultati: značilnosti zbranih
podatkov~(\ref{sec:podatki}), natančnost ocene
vrednosti~(\ref{sec:natancnost}), občutljivost algoritma na
parametre~(\ref{sec:obcutljivost}), latenca posameznega
vrednotenja~(\ref{sec:latenca}) in obnašanje pod
obremenitvijo~(\ref{sec:obremenitev}). Zatem preverimo funkcionalne in
nefunkcionalne zahteve (\ref{sec:funkcionalne}, \ref{sec:nefunkcionalne}) ter
predstavimo uporabniški preizkus~(\ref{sec:uporabniki}) in simulacijo
uporabe~(\ref{sec:simulacija}). Poglavje sklenemo z obravnavo groženj
veljavnosti~(\ref{sec:groznje}) in povzetkom ugotovitev glede na
zastavljena izhodišča~(\ref{sec:povzetek-analize}).

% =====================================================================
\section{Raziskovalna vprašanja in merila uspešnosti}
\label{sec:rq}

Na začetku diplomskega dela (poglavje~\ref{ch:uvod}) smo si zastavili, da je
mogoče uporabniku na podlagi zbranih tržnih podatkov in primerjave s podobnimi
vozili ponuditi \emph{uporabno} in \emph{dovolj natančno} oceno tržne vrednosti
vozila in da vključitev takega sistema izboljša uporabniško izkušnjo. Iz tega
izhodišča in iz nefunkcionalnih zahtev (razdelek~\ref{sec:zahteve}) izpeljemo
štiri raziskovalna vprašanja (RV1--RV4) in za vsako določimo merilo uspešnosti.

\emph{Mediana absolutne odstotne napake} (MdAPE) je osrednja mera natančnosti v
tem poglavju; nanjo se sklicujejo merila uspešnosti in vsi nadaljnji razdelki.
Izračunamo jo tako, da za vsak testni primer določimo absolutno odstotno napako,
torej za koliko odstotkov se ocenjena cena razlikuje od dejanske, nato pa čez
vse primere vzamemo mediano teh vrednosti. MdAPE $=15\,\%$ tako pomeni, da je pri
polovici vozil ocena znotraj $15\,\%$ od dejanske cene, pri drugi polovici pa je
odstopanje večje; manjša kot je vrednost, natančnejše so ocene. Mediano
uporabljamo namesto povprečja, ker je manj občutljiva na posamezne skrajne
napake. Natančne definicije MdAPE in preostalih mer so v
razdelku~\ref{sec:mere}.

Raziskovalna vprašanja so:

\begin{itemize}
  \item \textbf{RV1 --- natančnost ocene.} Kako natančno sistem oceni tržno
        vrednost vozila v primerjavi z dejanskimi oglasi na trgu?
        \emph{Merilo:} MdAPE je pri vozilih z vsaj $10$ primerljivimi
        vozili manjša od $15\,\%$, ocena pa je bistveno natančnejša od
        preproste (naivne) referenčne metode.
  \item \textbf{RV2 --- odvisnost natančnosti od podatkov.} Od česa je
        natančnost ocene najbolj odvisna (velikost vzorca primerljivih vozil,
        razpršenost ponudbe, segment vozila)?
  \item \textbf{RV3 --- hitrost in razširljivost.} Ali sistem zadošča
        nefunkcionalnim zahtevam glede odzivnosti tudi pod sočasno
        obremenitvijo, in kje je zgornja meja trenutne postavitve?
  \item \textbf{RV4 --- uporabnost.} Ali ciljni uporabniki znajo brez navodil
        opraviti glavne naloge (registracija, objava oglasa, vrednotenje)?
\end{itemize}
% =====================================================================
\section{Metodologija in merilno okolje}
\label{sec:metodologija}

\paragraph{Merilno okolje.}
Meritve hitrosti (razdelka~\ref{sec:latenca} in~\ref{sec:obremenitev}) so
izvedene v \emph{produkcijskem} okolju, torej na navideznem zasebnem strežniku
z dvema procesorskima jedroma in $4$~GB pomnilnika ter s tehnologijami in
njihovimi različicami, opisanimi v poglavju~\ref{ch:razvoj}
(razdelek~\ref{sec:namestitev}). Analitične in merilne skripte tečejo nad isto
podatkovno bazo kot je na produkciji.

\paragraph{Posnetek podatkov.}
Meritve v tem poglavju temeljijo na posnetku podatkovne baze iz septembra
$2026$, ki obsega približno $86.000$ normaliziranih oglasov
(razdelek~\ref{sec:podatki}). Ker se posnetek med analizo ne spreminja, so
rezultati ponovljivi.

\paragraph{Orodja.}
Za meritve in analizo v tem poglavju smo napisali namenske skripte, ločene od
produkcijske kode; ena od njih je dokumentirana kopija ocenjevalnega
algoritma, parametrizirana za eksperimente. Statistično obdelavo in izris
grafov smo izvedli s knjižnicama \texttt{NumPy}~\cite{harris2020numpy} in
\texttt{Matplotlib}~\cite{hunter2007matplotlib}.

\paragraph{Natančnost.}
Natančnost ocene (razdelek~\ref{sec:natancnost}) merimo z naključno izbranim
normaliziranim oglasom: njegove atribute (znamka, model, letnik, kilometrina)
uporabimo kot vhod v algoritem, sam oglas pa izločimo iz primerjalne množice. Nato izračunamo oceno in jo
primerjamo z \emph{dejansko oglaševano} ceno tega oglasa. Postopek ponovimo za
$3000$ naključnih oglasov.

% =====================================================================
\section{Uporabljene mere}
\label{sec:mere}

\paragraph{Latenca}
 je čas, ki preteče od trenutka tik pred oddajo zahtevka do
trenutka, ko odjemalec prejme celoten odgovor. Glede na to, kolikšen del te
poti zajamemo, ločimo tri vrste:
\begin{itemize}
  \item[(a)] \emph{latenca obdelave na strežniku} --- čas izvajanja funkcije
        za vrednotenje, od začetka do konca, brez omrežja;
  \item[(b)] \emph{latenca poizvedbe SQL} --- čas izvajanja poizvedbe v podatkovni bazi,
   od začetka do konca, brez omrežja;
  \item[(c)] \emph{latenca od konca do konca} (angl.\ \emph{end-to-end}) ---
        (a) skupaj z omrežnim obhodom med odjemalcem in strežnikom ter
        obdelavo zahtevka HTTP; to je čas, ki ga zazna uporabnik.
\end{itemize}
V razdelku~\ref{sec:latenca} merimo vrsti (a) in (b), v
razdelku~\ref{sec:obremenitev} pa vrsto (c).

\paragraph{Percentili.}
\textbf{Percentil} $p_k$ je vrednost, pod katero pade $k\,\%$ meritev. Poročamo
$p_{50}$ (mediana), $p_{95}$ in $p_{99}$.

\paragraph{Prepustnost, sočasnost, nasičenje.}
Obnašanje sistema pod obremenitvijo opišemo s tremi količinami:
\begin{itemize}
  \item \textbf{prepustnost} --- število uspešno obdelanih zahtevkov na
        sekundo;
  \item \textbf{sočasnost} --- število odjemalcev, ki hkrati držijo odprt
        zahtevek;
  \item \textbf{nasičenje} (\emph{koleno}) --- točka, v kateri dodatna
        sočasnost ne poveča več prepustnosti, latenca pa začne naraščati
        približno linearno.
\end{itemize}
Notranjo skladnost meritev preverimo z \textbf{Littlovim zakonom}~\cite{little1961},
\begin{equation}
  \label{eq:little}
  L = \lambda \cdot W,
\end{equation}
kjer je $L$ povprečno število zahtevkov v sistemu, $\lambda$ prepustnost in
$W$ povprečna latenca.

\paragraph{Napaka ocene.}
Za oglas $i$ z dejansko ceno $p_i$ in ocenjeno ceno $\hat p_i$ definiramo
absolutno odstotno napako
\begin{equation}
  \mathrm{APE}_i = \frac{\lvert \hat p_i - p_i \rvert}{p_i}.
\end{equation}
\textbf{MAPE} je povprečje $\mathrm{APE}_i$, \textbf{MdAPE} pa njihova mediana;
MdAPE je robustnejša na osamelce in je naša glavna mera~\cite{hyndman2006}.
Poročamo tudi \textbf{MAE} (povprečna absolutna napaka) in \textbf{RMSE}
(koren povprečne kvadratne napake) v evrih ter \emph{predznačeno} odstotno
napako, ki pokaže sistematičen odmik ocene navzgor ali navzdol.

\paragraph{Pokritost.}
\textbf{Pokritost} je delež poizvedb, za katere sistem vrne oceno, torej za
katere najde vsaj eno (oziroma vsaj $5$ oz.\ $10$) primerljivo vozilo.

\paragraph{Raznolikost ponudbe.}
\textbf{Raznolikost} (v izvorni kodi \texttt{ratio}) je relativna širina
cenovnega razpona primerljivih vozil,
\begin{equation}
  \label{eq:raznolikost}
  \mathrm{raznolikost} = \frac{P_{75} - P_{25}}{\hat{p}},
\end{equation}
kjer sta $P_{25}$ in $P_{75}$ (v kodi \texttt{goodPrice} in \texttt{badPrice})
uteženi $25.$ oziroma $75.$ percentil cen primerljivih vozil, $\hat{p}$
(v kodi \texttt{estimatedPrice}) pa ocenjena cena vozila. Vrednost izražamo v
odstotkih; enoten zapis uporabljamo v vseh tabelah tega poglavja.

% =====================================================================
\section{Značilnosti zbranih podatkov}
\label{sec:podatki}
\paragraph{Viri in obseg.}
Velika večina oglasov je bila zajeta z največjega slovenskega spletnega portala
za prodajo rabljenih vozil (razdelek~\ref{sec:zajem-podatkov}); manjši del
predstavlja vzorec z večjega nemškega oglasnega portala, ki smo ga uporabili za
preverjanje delovanja cevovoda na drugačni strukturi vira. Zajem je zajel
javno dostopne oglase za nekomercialno raziskovalno rabo, z omejeno hitrostjo
zahtevkov in brez shranjevanja kontaktnih ali osebnih podatkov prodajalcev.
Po normalizaciji baza obsega $85.846$ oglasov ($83.493$ s slovenskega in
$2.353$ z nemškega portala) za $10$ znamk vozil; razčlenitev po znamki je v
tabeli~\ref{tab:podatki-znamke}.

\begin{table}[htbp]
  \centering
  \caption{Število normaliziranih oglasov po znamki.}
  \label{tab:podatki-znamke}
  % vir: analysis/out/dataset_per_brand.csv
  \begin{tabular}{lr lr}
    \toprule
    \textbf{Znamka} & \textbf{Oglasi} & \textbf{Znamka} & \textbf{Oglasi} \\
    \midrule
    Volkswagen    & 21.191 & Peugeot     & 7.513 \\
    Audi          & 11.842 & Citroën     & 6.729 \\
    BMW           & 10.724 & Ford        & 5.783 \\
    Renault       & 9.615  & Toyota      & 2.648 \\
    Mercedes-Benz & 9.355  & Alfa Romeo  & 446 \\
    \bottomrule
  \end{tabular}
\end{table}

\paragraph{Porazdelitve ključnih atributov.}
Slika~\ref{fig:porazdelitve} prikazuje porazdelitve oglaševanih cen, prevožene
kilometrine in letnika. Porazdelitev cen je izrazito desno asimetrična (dolg
rep dragih vozil), zato jo prikazujemo na logaritemski osi. Mediana cene je
$14.890$~EUR (povprečje $18.137$~EUR), mediana kilometrine $159.950$~km,
mediana letnika pa $2017$; podrobnosti so v tabeli~\ref{tab:podatki-porazdelitve}.
Zaradi berljivosti so na grafih odrezane skrajne vrednosti (nekaj oglasov ima
očitno napačno ceno, kilometrino ali letnik); njihovo število je označeno v
kotu vsakega grafa, celoten razpon pa je razviden iz stolpca \emph{maks.}\ v
tabeli~\ref{tab:podatki-porazdelitve} in iz obravnave osamelcev v
tabeli~\ref{tab:osamelci}.

\begin{table}[htbp]
  \centering
  \caption{Povzetek porazdelitev ključnih numeričnih atributov.}
  \label{tab:podatki-porazdelitve}
  % vir: analysis/out/dataset_distribution_summary.json
  \setlength{\tabcolsep}{5pt}
  \begin{tabular}{lrrrrr}
    \toprule
    & \textbf{$p_{5}$} & \textbf{mediana} & \textbf{povprečje} & \textbf{$p_{95}$} & \textbf{maks.} \\
    \midrule
    cena [EUR]        & 1.550  & 14.890  & 18.137  & 49.014  & 374.000 \\
    kilometrina [km]  & 26.000 & 159.950 & 169.464 & 334.000 & 1.974.001 \\
    letnik            & 2005   & 2017    & 2016    & 2024    & 2029 \\
    moč [kW]          & 56     & 110     & 117     & 221     & 920 \\
    \bottomrule
  \end{tabular}
\end{table}

\begin{figure}[htbp]
  \centering
  \begin{subfigure}{0.48\textwidth}
    \includegraphics[width=\textwidth]{dataset_price_hist.pdf}
    \caption{oglaševana cena (log)}
  \end{subfigure}\hfill
  \begin{subfigure}{0.48\textwidth}
    \includegraphics[width=\textwidth]{dataset_km_hist.pdf}
    \caption{prevoženi kilometri}
  \end{subfigure}
  \\[1ex]
  \begin{subfigure}{0.48\textwidth}
    \includegraphics[width=\textwidth]{dataset_year_hist.pdf}
    \caption{letnik prve registracije}
  \end{subfigure}\hfill
  \begin{subfigure}{0.48\textwidth}
    \includegraphics[width=\textwidth]{dataset_brand_bar.pdf}
    \caption{oglasi po znamki}
  \end{subfigure}
  \caption{Porazdelitve ključnih atributov v normalizirani bazi.}
  \label{fig:porazdelitve}
  % vir: analysis/out/dataset_*_hist.pdf, dataset_brand_bar.pdf
\end{figure}

\paragraph{Pokritost po segmentih.}
Za pošteno oceno mora imeti algoritem dovolj primerljivih vozil v ožjem
segmentu (ista znamka in model, podoben letnik in kilometrina).
Slika~\ref{fig:pokritost} prikazuje toplotno karto za znamki Renault in
Volkswagen: vrstice so modeli, stolpci letniki, barva pa
število oglasov v tej kombinaciji. Vidno je, da so pogosti modeli srednjih
letnikov (npr.\ Volkswagen Golf ali Renault Clio) gosto pokriti, redki modeli
in skrajni letniki pa le redko, kar neposredno pojasni razlike v zanesljivosti
ocen (razdelek~\ref{sec:natancnost}).

\begin{figure}[htbp]
  \centering
  \includegraphics[width=\textwidth]{dataset_coverage_heatmap.pdf}
  \caption{Pokritost: število oglasov po modelu in letniku za znamki
  Renault in Volkswagen (po ena znamka na vrstico).}
  \label{fig:pokritost}
  % vir: analysis/out/dataset_coverage_heatmap.pdf
\end{figure}

\paragraph{Osamelci in kakovost podatkov.}
Postopek normalizacije trenutno ne izvaja filtriranja osamelcev, zato baza
vsebuje tudi vozila z očitno napačnimi vrednostmi.
Tabela~\ref{tab:osamelci} prikazuje njihov delež. Kar $10{,}7\,\%$ oglasov ima
vsaj eno vrednost zunaj smiselnih mej, največ pri kilometrini ($9{,}9\,\%$;
najvišja zabeležena vrednost je skoraj dva milijona kilometrov) in letniku
($8{,}1\,\%$; nekaj oglasov ima celo letnik v prihodnosti). Preprost filter po
interkvartilnem razponu (IQR) na ceni znotraj skupine
$(\text{znamka},\text{model})$ bi odstranil dodatnih $2{,}2\,\%$ oglasov.
Vpliv doslednega filtriranja osamelcev na natančnost obravnavamo v
razdelku~\ref{sec:obcutljivost}, uvedbo takega filtra pa v
poglavju~\ref{ch:sklep}.

\begin{table}[H]
  \centering
  \caption{Osamelci v normalizirani bazi ($85.846$ oglasov).}
  \label{tab:osamelci}
  % vir: analysis/out/dataset_outliers.csv
  \begin{tabularx}{\textwidth}{@{}X r r@{}}
    \toprule
    \textbf{Merilo} & \textbf{Št.\ vrstic} & \textbf{Delež} \\
    \midrule
    cena zunaj $[500,\ 200.000]$~EUR         & 427   & $0{,}5\,\%$ \\
    kilometrina zunaj $[1000,\ 500.000]$~km  & 8.495 & $9{,}9\,\%$ \\
    letnik zunaj $[1990,\ 2027]$             & 6.979 & $8{,}1\,\%$ \\
    vsaj ena vrednost zunaj mej              & 9.169 & $10{,}7\,\%$ \\
    dodatno po filtru IQR na ceni            & 1.878 & $2{,}2\,\%$ \\
    \bottomrule
  \end{tabularx}
\end{table}

% =====================================================================
\section{Natančnost ocene vrednosti}
\label{sec:natancnost}

To je osrednja meritev diplomskega dela: neposredno preverja zastavljeno
izhodišče o \emph{dovolj natančni} oceni (RV1). Protokol je opisan v
razdelku~\ref{sec:metodologija}, rezultate pa izdela skripta
\texttt{analysis/accuracy\_cv.py} nad vzorcem $3000$ oglasov (od tega jih je
$2986$ dobilo oceno).

\paragraph{Merilo resnice in namerni odmik ocene.}
Kot merilo resnice uporabljamo \emph{oglaševano} ceno, saj dejanske prodajne
cene niso javno dostopne. Sistem pa namenoma ne ocenjuje oglaševane cene,
temveč \emph{pričakovano prodajno ceno}: ker se rabljena vozila v pogajanjih
praviloma prodajo pod oglaševano ceno, oceno zniža s stopenjskim popustom
($9$--$15\,\%$ glede na cenovni razred). Zato je pri primerjavi z oglaševano
ceno pričakovan rahel \emph{negativni} odmik ocene, kar je zaželeno in ne
pomeni slabše ocene. To lastnost upoštevamo pri branju vseh rezultatov v
nadaljevanju: napako merimo dvakrat --- proti oglaševani ceni z vključenim
popustom (kar dejansko vrne sistem) in brez popusta (čista primerjalna
ocena).

\paragraph{Skupni rezultat.}
Brez popusta znaša MdAPE predlaganega pristopa $13{,}5\,\%$, MAPE $23{,}8\,\%$,
MAE $2798$~EUR in RMSE $5151$~EUR; z vključenim namernim popustom je MdAPE
proti oglaševani ceni $14{,}5\,\%$ (mediana predznačene napake $-6{,}8\,\%$).
Pokritost je $99{,}5\,\%$ pri $n \ge 1$ in $95{,}7\,\%$ pri $n \ge 10$.
Slika~\ref{fig:pred-vs-actual} prikazuje napovedano ceno proti dejanski
($R^2 = 0{,}87$), slika~\ref{fig:ape-hist} pa porazdelitev odstotne napake.

\begin{figure}[htbp]
  \centering
  \begin{subfigure}{0.48\textwidth}
    \includegraphics[width=\textwidth]{accuracy_pred_vs_actual.pdf}
    \caption{napovedano proti dejanskemu}
    \label{fig:pred-vs-actual}
  \end{subfigure}\hfill
  \begin{subfigure}{0.48\textwidth}
    \includegraphics[width=\textwidth]{accuracy_ape_hist.pdf}
    \caption{porazdelitev napake}
    \label{fig:ape-hist}
  \end{subfigure}
  \caption{Natančnost ocene vrednosti (leave-one-out, brez popusta).}
  % vir: analysis/out/accuracy_pred_vs_actual.pdf, accuracy_ape_hist.pdf
\end{figure}

\paragraph{Primerjava z osnovnimi modeli.}
Da presodimo, ali je predlagani pristop boljši od trivialnega, ga primerjamo s
štirimi referenčnimi modeli na istem protokolu in isti množici kandidatov
(tabela~\ref{tab:baselines}, slika~\ref{fig:baselines}): B0~globalna mediana
cene za dani model, B1~navadno povprečje cen v filtrirnem oknu, B1m~mediana cen
v oknu, B2~linearna regresija $\text{cena} \sim \text{letnik} + \text{km}$,
B3~metoda $k$~najbližjih sosedov ($k = 10$) po standardiziranih letniku in
kilometrini, ter B4~naš pristop (uteženo povprečje s podobnostnim jedrom, brez
popusta).

\begin{table}[htbp]
  \centering
  \caption{Primerjava natančnosti z osnovnimi modeli (leave-one-out).
  MdAPE in MAPE sta v odstotkih, MAE in RMSE v evrih.}
  \label{tab:baselines}
  % vir: analysis/out/accuracy_overall.csv
  \setlength{\tabcolsep}{5pt}
  \begin{tabular}{l r r r r}
    \toprule
    \textbf{Model} & \textbf{MdAPE} & \textbf{MAPE} & \textbf{MAE} & \textbf{RMSE} \\
    \midrule
    B0 globalna mediana       & 37{,}8 & 95{,}7 & 7395 & 10883 \\
    B1 povprečje v oknu       & 14{,}2 & 24{,}9 & 2930 & 5333 \\
    B1m mediana v oknu        & 13{,}1 & 22{,}5 & 2877 & 5561 \\
    B2 linearna regresija     & 12{,}9 & 22{,}4 & 2677 & 5032 \\
    B3 $k$-najbližjih sosedov & \textbf{11{,}2} & 21{,}5 & 2497 & 4883 \\
    B4 naš pristop            & 13{,}5 & 23{,}8 & 2798 & 5151 \\
    \bottomrule
  \end{tabular}
\end{table}

\begin{figure}[htbp]
  \centering
  \includegraphics[width=0.7\textwidth]{accuracy_baselines.pdf}
  \caption{MdAPE osnovnih modelov in predlaganega pristopa.}
  \label{fig:baselines}
  % vir: analysis/out/accuracy_baselines.pdf
\end{figure}

Predlagani pristop je \emph{bistveno} boljši od naivne globalne mediane
($13{,}5\,\%$ proti $37{,}8\,\%$) in primerljiv s preostalimi metodami, ki se
gibljejo med $11$ in $14\,\%$. Ni pa najboljši: metoda $k$~najbližjih sosedov
doseže $11{,}2\,\%$, tudi navadna mediana istega okna ($13{,}1\,\%$) je
malenkost natančnejša. To je pričakovano za zasnovo, katere cilj je
\emph{razložljivost}: za razliko od naštetih metod predlagani pristop poleg
točkovne ocene vrne tudi cenovni razpon in indeks zaupanja, uteži pa so
neposredno interpretabilne. Rezultat hkrati potrjuje smiselnost prehoda na
metode strojnega učenja, ki jo predlagamo v poglavju~\ref{ch:sklep}: če že
preprost $k$-NN preseže sedanji pristop, je od zmogljivejših modelov mogoče
pričakovati še boljše rezultate.

\paragraph{Odvisnost od velikosti vzorca in segmenta (RV2).}
Slika~\ref{fig:mdape-n} prikazuje MdAPE glede na število primerljivih vozil
$n$: napaka je pri $n < 10$ okoli $20\,\%$, pri $n \ge 20$ pa se ustali pri
približno $13\,\%$. To empirično podpira izbiro konstante $20$ v indeksu
zaupanja (enačba~\ref{eq:confidence}) --- pri toliko primerljivih vozilih je
ocena že blizu svoje najboljše natančnosti. Slika~\ref{fig:mdape-brand} prikazuje MdAPE po znamki: najbolje se
odrežejo pogosti, gosto pokriti modeli (Audi $9{,}8\,\%$, Toyota $10{,}8\,\%$,
Volkswagen $12{,}1\,\%$), najslabše pa Mercedes-Benz ($24{,}9\,\%$), pri
katerem se pri normalizaciji več različnih modelov zlije v isto kanonično
ime, kar poslabša izbor primerljivih vozil.
Tabela~\ref{tab:natancnost-segmenti} kaže še razčlenitev po starosti,
kilometrini in raznolikosti ponudbe.

\begin{figure}[htbp]
  \centering
  \begin{subfigure}{0.48\textwidth}
    \includegraphics[width=\textwidth]{accuracy_mdape_vs_n.pdf}
    \caption{glede na velikost vzorca $n$}
    \label{fig:mdape-n}
  \end{subfigure}\hfill
  \begin{subfigure}{0.48\textwidth}
    \includegraphics[width=\textwidth]{accuracy_mdape_by_brand.pdf}
    \caption{po znamki}
    \label{fig:mdape-brand}
  \end{subfigure}
  \caption{Odvisnost natančnosti ocene od podatkov.}
  % vir: analysis/out/accuracy_mdape_vs_n.pdf, accuracy_mdape_by_brand.pdf
\end{figure}

\begin{table}[htbp]
  \centering
  \caption{Natančnost predlaganega pristopa po segmentih (MdAPE).}
  \label{tab:natancnost-segmenti}
  % vir: analysis/out/accuracy_by_age.csv, accuracy_by_km.csv, accuracy_by_ratio.csv
  \begin{tabularx}{\textwidth}{@{}l X r r@{}}
    \toprule
    \textbf{Razčlenitev} & \textbf{Koš} & \textbf{Primeri} & \textbf{MdAPE [\%]} \\
    \midrule
    starost [leta] & 0--3   & 367 & 9{,}0 \\
                   & 4--7   & 921 & 10{,}6 \\
                   & 8--12  & 932 & 14{,}5 \\
                   & $13+$  & 766 & 22{,}9 \\
    \midrule
    kilometrina    & $<50.000$        & 280  & 9{,}6 \\
                   & $50.000$--$100.000$  & 449  & 11{,}9 \\
                   & $100.000$--$150.000$ & 634  & 12{,}4 \\
                   & $150.000$--$200.000$ & 629  & 13{,}7 \\
                   & $>200.000$       & 1021 & 16{,}7 \\
    \midrule
    raznolikost    & $0$--$0{,}10$    & 88   & 7{,}5 \\
                   & $0{,}10$--$0{,}25$ & 1002 & 8{,}4 \\
                   & $0{,}25$--$0{,}50$ & 1613 & 15{,}5 \\
                   & $>0{,}50$        & 447  & 29{,}7 \\
    \bottomrule
  \end{tabularx}
\end{table}

Vzorec pokaže jasne zakonitosti: napaka narašča s starostjo vozila ($9\,\%$ pri
najnovejših, $23\,\%$ pri vozilih, starejših od $13$~let), s prevoženo
kilometrino in --- najmočneje --- z \emph{raznolikostjo ponudbe}. Pri modelih z
enotno ponudbo (raznolikost pod $0{,}1$) je MdAPE le $7{,}5\,\%$, pri razpršeni
ponudbi (nad $0{,}5$) pa $29{,}7\,\%$. To potrjuje, da je raznolikost, ki jo
sistem uporabniku tudi prikaže, dober vgrajen pokazatelj zanesljivosti
posamezne ocene.

\paragraph{Vpliv popusta.}
Slika~\ref{fig:popust} prikazuje predznačeno napako z in brez popusta. Brez
popusta je ocena rahlo optimistična (mediana $+3{,}8\,\%$ nad oglaševano ceno),
z namernim popustom pa se premakne na $-6{,}8\,\%$ pod oglaševano ceno. Ta
premik je skladen z zasnovo sistema, ki cilja pričakovano prodajno in ne
oglaševane cene; dokončno umerjanje popusta bi bilo mogoče šele s podatki o
dejanskih prodajah (razdelek~\ref{sec:groznje}).

\begin{figure}[htbp]
  \centering
  \includegraphics[width=0.7\textwidth]{accuracy_discount_bias.pdf}
  \caption{Predznačena napaka proti oglaševani ceni, z in brez namernega popusta.}
  \label{fig:popust}
  % vir: analysis/out/accuracy_discount_bias.pdf
\end{figure}

\paragraph{Časovni razrez.}
Pri napovedovanju najnovejših $10\,\%$ oglasov zgolj iz starejših $90\,\%$
znaša MdAPE predlaganega pristopa $13{,}0\,\%$ (tabela~\ref{tab:temporal}), kar
je celo malenkost bolje od rezultata leave-one-out in potrjuje, da ocena ne
izhaja iz uhajanja informacije.

\begin{table}[htbp]
  \centering
  \caption{Natančnost pri časovnem razrezu (napoved najnovejših $10\,\%$).
  MdAPE in MAPE sta v odstotkih, MAE in RMSE v evrih.}
  \label{tab:temporal}
  % vir: analysis/out/accuracy_temporal.csv
  \setlength{\tabcolsep}{5pt}
  \begin{tabular}{l r r r r}
    \toprule
    \textbf{Model} & \textbf{MdAPE} & \textbf{MAPE} & \textbf{MAE} & \textbf{RMSE} \\
    \midrule
    B0 globalna mediana       & 35{,}9 & 74{,}5 & 7596 & 11222 \\
    B1m mediana v oknu        & 12{,}8 & 19{,}7 & 3088 & 5805 \\
    B3 $k$-najbližjih sosedov & 10{,}5 & 18{,}2 & 2667 & 4890 \\
    B4 naš pristop            & 13{,}0 & 20{,}6 & 2994 & 5347 \\
    \bottomrule
  \end{tabular}
\end{table}

% =====================================================================
\section{Občutljivost algoritma na parametre}
\label{sec:obcutljivost}

Ocenjevalni algoritem vsebuje več hevristično določenih konstant: dovoljeno
odstopanje kilometrine ($\pm 40.000$~km) in letnika ($\pm 3$~leta), uteži
delnih ocen ($0{,}6$ in $0{,}4$), oblika padajoče funkcije in stopenjski
popust. V tem razdelku vsako spreminjamo ob nespremenjenih ostalih in merimo
MdAPE ter pokritost (vzorec $2000$ oglasov, leave-one-out). Rezultate izdela
skripta \texttt{analysis/sensitivity.py}.

\paragraph{Okno primerljivosti.}
Tabela~\ref{tab:sens-okno} in slika~\ref{fig:sens-km} kažeta, da je algoritem
na širino \emph{kilometrinskega} okna skoraj neobčutljiv: MdAPE se med
$\pm 10.000$~km in $\pm 80.000$~km giblje le med $13{,}2$ in $13{,}6\,\%$.
Nasprotno je odvisen od \emph{letnikovega} okna: pri $\pm 1$~letu je MdAPE
$12{,}6\,\%$, pri $\pm 5$~letih pa $15{,}2\,\%$. Ožje okno da bolj primerljiva
vozila, a manjši vzorec (nižja pokritost); gre za znano razmerje med
pristranskostjo in varianco~\cite{bias-variance}. Toplotna karta
(slika~\ref{fig:sens-heatmap}) potrjuje, da natančnost določa skoraj izključno
letnikovo okno. Sedanja izbira $\pm 3$~leta je razumen kompromis med
natančnostjo in pokritostjo; morebitna zožitev na $\pm 2$~leti bi MdAPE
znižala za približno $0{,}3$~odstotne točke ob še vedno visoki pokritosti.

\begin{table}[htbp]
  \centering
  \caption{Občutljivost na širino okna primerljivosti. Pokritost je delež
  poizvedb z vsaj $5$ primerljivimi vozili.}
  \label{tab:sens-okno}
  % vir: analysis/out/sensitivity_km.csv, sensitivity_year.csv
  \setlength{\tabcolsep}{5pt}
  \begin{tabular}{l r r r}
    \toprule
    \textbf{Parameter} & \textbf{MdAPE [\%]} & \textbf{Pokritost [\%]} & \textbf{Mediana $n$} \\
    \midrule
    kilometrina $\pm 10.000$~km & 13{,}2 & 92{,}3 & 53 \\
    kilometrina $\pm 40.000$~km & 13{,}4 & 98{,}1 & 196 \\
    kilometrina $\pm 80.000$~km & 13{,}6 & 99{,}0 & 362 \\
    \midrule
    letnik $\pm 1$~leto  & 12{,}6 & 95{,}6 & 97 \\
    letnik $\pm 3$~leta  & 13{,}4 & 98{,}1 & 196 \\
    letnik $\pm 5$~let   & 15{,}2 & 98{,}5 & 262 \\
    \bottomrule
  \end{tabular}
\end{table}

\begin{figure}[htbp]
  \centering
  \begin{subfigure}{0.48\textwidth}
    \includegraphics[width=\textwidth]{sensitivity_km.pdf}
    \caption{okno kilometrine}
  \end{subfigure}\hfill
  \begin{subfigure}{0.48\textwidth}
    \includegraphics[width=\textwidth]{sensitivity_year.pdf}
    \caption{okno letnika}
  \end{subfigure}
  \caption{Občutljivost na širino okna primerljivosti (MdAPE in pokritost).}
  \label{fig:sens-km}
  % vir: analysis/out/sensitivity_km.pdf, sensitivity_year.pdf
\end{figure}

\begin{figure}[htbp]
  \centering
  \includegraphics[width=0.7\textwidth]{sensitivity_km_year_heatmap.pdf}
  \caption{MdAPE [\%] po mreži dovoljenih odstopanj kilometrine in letnika.}
  \label{fig:sens-heatmap}
  % vir: analysis/out/sensitivity_km_year_heatmap.pdf
\end{figure}

\paragraph{Uteži, oblika jedra in popust.}
Tabela~\ref{tab:sens-ostalo} povzema preostale parametre. Razmerje uteži med
kilometrino in letnikom na natančnost praktično ne vpliva ($13{,}4$ do
$13{,}6\,\%$ za vsa preizkušena razmerja), prav tako ne zamenjava linearne
padajoče funkcije z eksponentno ($13{,}4$ proti $13{,}3\,\%$). Edina parametra,
ki opazno premakneta rezultat, sta torej širina letnikovega okna in namerni
popust: ta predznačeno napako proti oglaševani ceni premakne iz $+3{,}8\,\%$ v
$-6{,}8\,\%$, MdAPE pa zviša z $13{,}4$ na $14{,}5\,\%$. Ta premik je posledica
zavestne odločitve, da sistem ocenjuje pričakovano prodajno in ne oglaševane
cene, in ne slabše ocene.

\begin{table}[htbp]
  \centering
  \caption{Občutljivost na uteži, obliko jedra in popust.}
  \label{tab:sens-ostalo}
  % vir: analysis/out/sensitivity_weights.csv, sensitivity_kernel.csv, sensitivity_discount.csv
  \begin{tabularx}{\textwidth}{@{}l X r >{\raggedleft\arraybackslash}p{2.3cm}@{}}
    \toprule
    \textbf{Parameter} & \textbf{Nastavitev} & \textbf{MdAPE [\%]} & \textbf{Predznačena napaka [\%]} \\
    \midrule
    uteži (km / letnik) & $0{,}5 / 0{,}5$              & 13{,}4 & $+3{,}7$ \\
    uteži (km / letnik) & $0{,}6 / 0{,}4$ (privzeto)   & 13{,}4 & $+3{,}8$ \\
    uteži (km / letnik) & $0{,}8 / 0{,}2$              & 13{,}6 & $+3{,}8$ \\
    \midrule
    jedro               & linearno (privzeto)         & 13{,}4 & $+3{,}8$ \\
    jedro               & eksponentno                 & 13{,}3 & $+3{,}8$ \\
    \midrule
    popust              & brez                        & 13{,}4 & $+3{,}8$ \\
    popust              & stopenjski (privzeto)        & 14{,}5 & $-6{,}8$ \\
    \bottomrule
  \end{tabularx}
\end{table}

\paragraph{Sklep o robustnosti.}
Algoritem je robusten: na večino svojih konstant se skoraj ne odziva, kar
pomeni, da njihova hevristična izbira ni sporna. Njegova natančnost je odvisna
predvsem od kakovosti in gostote podatkov (razdelek~\ref{sec:natancnost}), ne
od finega nastavljanja parametrov.

% =====================================================================
\section{Latenca posameznega vrednotenja}
\label{sec:latenca}

Merimo latenco vrste (a) in (b) iz razdelka~\ref{sec:mere}. Meritev izvede
\texttt{benchmarks/bench\_algo.py}: $2000$ neposrednih klicev funkcije nad
naključnimi realnimi vhodi (brez HTTP), po $100$ ogrevalnih klicih, ki se ne
merijo. Ločeno merimo čas poizvedbe SQL, čas točkovanja v Pythonu in čas
serializacije odgovora.

\paragraph{Ozko grlo: pregled celotne tabele.}
Prvotna izvedba poizvedbe ni imela ustreznega indeksa.
Izsek~\ref{lst:explain-before} prikazuje izpis \texttt{EXPLAIN (ANALYZE,
BUFFERS)}: načrt je \texttt{Seq Scan}, baza pri vsakem vrednotenju prebere
celotno tabelo in zavrže $85.273$ vrstic, ki ne ustrezajo pogojem
(\texttt{Rows Removed by Filter}). Ker so podatki v celoti v predpomnilniku
(\texttt{shared hit}, brez \texttt{read}), je čas stabilen, a raste linearno z
velikostjo tabele.

\begin{center}
\captionof{listing}{\texttt{EXPLAIN (ANALYZE, BUFFERS)} pred dodajanjem indeksa.}
\label{lst:explain-before}
\begin{minted}{text}
Seq Scan on normalized_market_listing
    (cost=0.00..3314.15 rows=12 width=97)
    (actual time=0.555..6.847 rows=573 loops=1)
  Filter: ((kilometers >= 18500) AND (kilometers <= 98500)
       AND (registration_year >= 2019) AND (registration_year <= 2025)
       AND (brand_name = 'Citroën') AND (model_name = 'C3'))
  Rows Removed by Filter: 85273
  Buffers: shared hit=1168
Execution Time: 6.880 ms
\end{minted}
\end{center}

\paragraph{Učinek sestavljenega indeksa.}
Po dodajanju sestavljenega indeksa na
$(\text{brand\_name},\ \text{model\_name},\ \text{registration\_year},\
\text{kilometers})$ se načrt spremeni v \texttt{Bitmap Index Scan}
(izsek~\ref{lst:explain-after}): baza prebere le nekaj strani indeksa in
pripadajoče vrstice, čas izvajanja pa pade z $6{,}88$~ms na $0{,}69$~ms.
Latenca poizvedbe SQL se v mediani zmanjša z $9{,}7$~ms na $2{,}0$~ms
(tabela~\ref{tab:indeks-pred-po}, slika~\ref{fig:indeks-pred-po}). Pomembneje
od trenutnega pohitrenja je, da postane čas poizvedbe neodvisen od velikosti
tabele: namesto $O(\text{tabela})$ postane $O(\text{zadetki})$.

\begin{table}[htbp]
  \centering
  \caption{Latenca poizvedbe SQL pred in po dodajanju indeksa, v milisekundah.}
  \label{tab:indeks-pred-po}
  % vir: benchmarks/report.md (results_algo_before/after.json)
  \begin{tabular}{lrrrr}
    \toprule
    & \textbf{p50} & \textbf{p95} & \textbf{p99} & \textbf{Povprečje} \\
    \midrule
    pred (Seq Scan)  & 9{,}7 & 21{,}0 & 24{,}8 & 11{,}3 \\
    po (Index Scan)  & 2{,}0 & 8{,}6  & 12{,}5 & 3{,}0 \\
    \bottomrule
  \end{tabular}
\end{table}

\begin{figure}[htbp]
  \centering
  \includegraphics[width=0.6\textwidth]{latency_index_before_after.pdf}
  \caption{Latenca poizvedbe SQL pred in po dodajanju sestavljenega indeksa.}
  \label{fig:indeks-pred-po}
  % vir: analysis/out/latency_index_before_after.pdf
\end{figure}

\begin{center}
\captionof{listing}{\texttt{EXPLAIN (ANALYZE, BUFFERS)} po dodajanju indeksa.}
\label{lst:explain-after}
\begin{minted}{text}
Bitmap Heap Scan on normalized_market_listing
    (cost=5.84..50.49 rows=12 width=97)
    (actual time=0.202..0.615 rows=573 loops=1)
  Recheck Cond: (brand_name = 'Citroën' AND model_name = 'C3'
       AND registration_year BETWEEN 2019 AND 2025
       AND kilometers BETWEEN 18500 AND 98500)
  Buffers: shared hit=195
  ->  Bitmap Index Scan on idx_nml_valuation
        (actual time=0.166..0.166 rows=573 loops=1)
        Buffers: shared hit=6
Execution Time: 0.688 ms
\end{minted}
\end{center}

\paragraph{Razčlenitev latence po komponentah.}
Tabela~\ref{tab:latenca-vrednotenja} in slika~\ref{fig:latency-breakdown}
prikazujeta latenco celotnega vrednotenja po dodajanju indeksa. Tipično
vrednotenje traja $2{,}0$~ms, v najslabšem odstotku primerov (p99) pa
$13{,}5$~ms. Skoraj celotna latenca ($\approx 88\,\%$ povprečja) odpade na
poizvedbo SQL; točkovanje v Pythonu ($0{,}4$~ms povprečno) in serializacija
odgovora ($\approx 0{,}01$~ms) sta zanemarljiva. Poizvedba v povprečju vrne
$287$ primerljivih vozil (mediana $175$, največ $1495$).

\begin{table}[htbp]
  \centering
  \caption{Latenca posameznega vrednotenja po dodajanju indeksa
  ($N = 2000$), v milisekundah.}
  \label{tab:latenca-vrednotenja}
  % vir: benchmarks/report.md (results_algo.json)
  \begin{tabular}{lrrrr}
    \toprule
    \textbf{Komponenta} & \textbf{p50} & \textbf{p95} & \textbf{p99} & \textbf{Povprečje} \\
    \midrule
    Skupaj                & 2{,}0  & 9{,}7 & 13{,}5 & 3{,}1 \\
    Poizvedba SQL         & 1{,}8  & 8{,}4 & 11{,}5 & 2{,}7 \\
    Točkovanje (Python)   & 0{,}2  & 1{,}2 & 1{,}8  & 0{,}4 \\
    Serializacija (JSON)  & 0{,}01 & 0{,}02 & 0{,}03 & 0{,}01 \\
    \bottomrule
  \end{tabular}
\end{table}

\begin{figure}[htbp]
  \centering
  \includegraphics[width=0.6\textwidth]{latency_breakdown.pdf}
  \caption{Razčlenitev latence posameznega vrednotenja po komponentah.}
  \label{fig:latency-breakdown}
  % vir: analysis/out/latency_breakdown.pdf
\end{figure}

Vrednotenje se torej izvede daleč pod nefunkcionalno zahtevo $2$~s
(razdelek~\ref{sec:zahteve}); prostora za rast podatkovne baze je ob
prisotnosti indeksa dovolj.

% =====================================================================
\section{Obnašanje pod obremenitvijo}
\label{sec:obremenitev}

Merimo end-to-end latenco (vrsta (c)) in prepustnost končne točke
\texttt{POST /valuate} pri naraščajoči sočasnosti. Obremenitev ustvarja
skripta \texttt{benchmarks/loadtest.py} z več vzporednimi odjemalci, ki v
fiksnem $20$-sekundnem oknu neprekinjeno pošiljajo zahtevke z naključnimi
realnimi vhodi. Test se samodejno prekine ob deležu napak nad $5\,\%$ ali
$p_{95}$ nad $5$~s. Med testom vzorčimo tudi zasedenost baze
(\texttt{pg\_stat\_activity}).

\begin{table}[htbp]
  \centering
  \caption{Prepustnost in latenca pri naraščajoči sočasnosti.
  Latence so v milisekundah, $L$ je ocena po Littlovem zakonu.}
  \label{tab:obremenitev}
  % vir: benchmarks/results_load.json (loadtest.py --db-sample)
  \begin{tabular}{r r r r r r r}
    \toprule
    \textbf{Sočasnost} & \textbf{Prep.\ [z/s]} & \textbf{p50} & \textbf{p95} & \textbf{p99}
      & \textbf{Napake} & \textbf{$L$} \\
    \midrule
    1  & \res{x} & \res{x} & \res{x} & \res{x} & \res{x} & \res{x} \\
    5  & \res{x} & \res{x} & \res{x} & \res{x} & \res{x} & \res{x} \\
    10 & \res{x} & \res{x} & \res{x} & \res{x} & \res{x} & \res{x} \\
    20 & \res{x} & \res{x} & \res{x} & \res{x} & \res{x} & \res{x} \\
    40 & \res{x} & \res{x} & \res{x} & \res{x} & \res{x} & \res{x} \\
    \bottomrule
  \end{tabular}
\end{table}

\paragraph{Rezultat.}
\res{Prepustnost narašča približno linearno do sočasnosti X, nato se ustali pri
približno X zahtevkih na sekundo. Vse latence ostanejo pod nefunkcionalno
zahtevo 2~s.}

\paragraph{Skladnost meritve (Littlov zakon).}
\res{Stolpec $L$ (prepustnost krat povprečna latenca) ostaja blizu nastavljene
sočasnosti, kar potrjuje notranjo skladnost meritve.}

\paragraph{Dokaz ozkega grla.}
Teoretični strop postavitve je $1000 / S \approx 324$ vrednotenj na sekundo,
kjer je $S = 3{,}1$~ms povprečni čas ene evalvacije. Dejanska meja je nižja,
ker ves promet serializira \emph{ena sama deljena povezava do baze}
(\texttt{app/database.py}) in en sam strežniški proces.
\res{To potrjuje vzorčenje baze med testom: število aktivnih poizvedb kljub
naraščajoči sočasnosti odjemalcev ne preseže X.}

% =====================================================================
\section{Preverjanje funkcionalnih zahtev}
\label{sec:funkcionalne}

Tabela~\ref{tab:sledljivost} povezuje vsako funkcionalno zahtevo iz
razdelka~\ref{sec:zahteve} z načinom preverjanja in izidom.

\begin{table}[htbp]
  \centering
  \caption{Sledljivost funkcionalnih zahtev do preverjanja.}
  \label{tab:sledljivost}
  % vir: ročni scenariji + imena Jest testov v car-stock-market-be
  \begin{tabularx}{\textwidth}{@{}X X c@{}}
    \toprule
    \textbf{Zahteva} & \textbf{Način preverjanja} & \textbf{Izid} \\
    \midrule
    Registracija uporabnika              & \res{test / scenarij} & \res{OK} \\
    Prijava in odjava                    & \res{test / scenarij} & \res{OK} \\
    Objava oglasa (le registrirani)      & \res{test / scenarij} & \res{OK} \\
    Vrednotenje vozila po atributih      & \res{test / scenarij} & \res{OK} \\
    Všečkanje oglasov                    & \res{test / scenarij} & \res{OK} \\
    Ažuriranje lastnih/všečkanih oglasov & \res{test / scenarij} & \res{OK} \\
    Pregled oglasov (anonimni)           & \res{test / scenarij} & \res{OK} \\
    \bottomrule
  \end{tabularx}
\end{table}

\paragraph{Pokritost testov.}
Avtomatizirani testi zalednega sistema (Jest + Supertest,
razdelek~\ref{sec:testiranje-be}) obsegajo \res{X} testov v \res{X} zbirkah in
pokrivajo \res{X}\,\% stavkov ter \res{X}\,\% vej. Testi se izvedejo ob vsakem
zahtevku za združitev na cevovodu neprekinjene integracije
(razdelek~\ref{subsec:ci}).

% =====================================================================
\section{Preverjanje nefunkcionalnih zahtev}
\label{sec:nefunkcionalne}

Tabela~\ref{tab:nfr} za vsako nefunkcionalno zahtevo iz
razdelka~\ref{sec:zahteve} navaja merilo, izmerjeno vrednost in sodbo.

\begin{table}[htbp]
  \centering
  \caption{Preverjanje nefunkcionalnih zahtev.}
  \label{tab:nfr}
  \setlength{\tabcolsep}{4pt}
  \begin{tabularx}{\textwidth}{@{}>{\raggedright\arraybackslash}X
                                 >{\raggedright\arraybackslash}X
                                 >{\raggedright\arraybackslash}X c@{}}
    \toprule
    \textbf{Zahteva} & \textbf{Merilo} & \textbf{Izmerjeno} & \textbf{Sodba} \\
    \midrule
    Odziv vmesnika na dejanje uporabnika & $< 100$~ms & \res{x}~ms  & \res{OK} \\
    Zahtevki HTTP (razen vrednotenja)    & $< 300$~ms & \res{x}~ms ($p_{95}$) & \res{OK} \\
    Vrednotenje vozila                   & $< 2$~s    & $13{,}5$~ms ($p_{99}$, brez omrežja) & OK \\
    Vizualna privlačnost / dostopnost    & Lighthouse & \res{x} / \res{x} & \res{OK} \\
    Prilagodljiva postavitev             & preizkus na več velikostih zaslona & \res{da} & \res{OK} \\
    Varovanje uporabniških podatkov      & bcrypt, JWT z rokom & \res{opis} & \res{OK} \\
    Odzivnost pri zahtevnejših opravilih & asinhrona obdelava & \res{opis} & \res{OK} \\
    \bottomrule
  \end{tabularx}
\end{table}

\paragraph{Meritev odzivnosti vmesnika in endpointov.}
\res{Odzivnost uporabniškega vmesnika smo izmerili z orodjem Lighthouse na treh
ključnih straneh; ocene za performanco so X, za dostopnost pa X. Latenco
glavnih endpointov zalednega sistema smo izmerili z orodjem autocannon:
$p_{95}$ je v vseh primerih pod X~ms, torej pod zahtevo 300~ms.}

\paragraph{Varnost.}
Gesla so pred shranjevanjem zgoščena z algoritmom bcrypt
(razdelek~\ref{sec:avtentikacija}), žetoni JWT imajo rok veljavnosti $7$~dni,
poti vmesnika API pa ščiti sredinski preverjevalnik žetonov. \res{Orodje
\texttt{npm audit} ob zadnjem pregledu ni javilo kritičnih ranljivosti v
odvisnostih.}

% =====================================================================
\section{Uporabniški preizkus}
\label{sec:uporabniki}

Izvedli smo formativni uporabniški preizkus~\cite{nielsen1994} z namenom
odkriti težave uporabnosti, ne pa statistično ovrednotiti izdelka.

\paragraph{Udeleženci in postopek.}
V preizkus smo vključili \res{N} udeležencev iz kroga družine in bližnjih
prijateljev, katerih starost in računalniška pismenost sta se med seboj
razlikovali (tabela~\ref{tab:udelezenci}). Vsak udeleženec je prejel enak
seznam nalog: registracija, objava oglasa, všečkanje oglasa, izbira
slovenskega jezika v nastavitvah in vrednotenje poljubnega vozila. Med
izvajanjem je udeleženec razmišljal na glas, opazovalec pa je beležil, ali je
nalogo dokončal samostojno, in vsako opaženo težavo.

\begin{table}[htbp]
  \centering
  \caption{Udeleženci uporabniškega preizkusa.}
  \label{tab:udelezenci}
  \begin{tabular}{c c c l}
    \toprule
    \textbf{\#} & \textbf{Starost (razpon)} & \textbf{Rač.\ pismenost (1--5)} & \textbf{Naprava} \\
    \midrule
    1 & \res{x} & \res{x} & \res{x} \\
    2 & \res{x} & \res{x} & \res{x} \\
    3 & \res{x} & \res{x} & \res{x} \\
    \res{\dots} & & & \\
    \bottomrule
  \end{tabular}
\end{table}

\paragraph{Rezultati.}
Vsi udeleženci so vse naloge dokončali brez pomoči, kar potrjuje, da so glavne
poti v vmesniku dovolj jasne (RV4). Kljub temu so odkrili nekaj
pomanjkljivosti; tabela~\ref{tab:tezave} jih razvršča po resnosti (lestvica
$0$--$4$ po Nielsenu~\cite{nielsen1994}) in navaja, ali so bile odpravljene.

\begin{table}[htbp]
  \centering
  \caption{Težave, odkrite med uporabniškim preizkusom
  (resnost po lestvici $0$--$4$).}
  \label{tab:tezave}
  \setlength{\tabcolsep}{5pt}
  \begin{tabularx}{\textwidth}{@{}>{\raggedright\arraybackslash}X c c@{}}
    \toprule
    \textbf{Težava} & \textbf{Resnost} & \textbf{Status} \\
    \midrule
    Premajhni naslovi podstrani v glavi glavne strani       & \res{x} & odpravljeno \\
    Manjkajoče polje za ponovni vnos gesla pri registraciji & \res{x} & odpravljeno \\
    Velika števila brez ločil za tisočice                   & \res{x} & odpravljeno \\
    Manjkajoč opis pod slikami oblik vozil v obrazcu        & \res{x} & odpravljeno \\
    Premalo vozil, dostopnih za vrednotenje                 & \res{x} & \res{v načrtu} \\
    \bottomrule
  \end{tabularx}
\end{table}

Omejitve preizkusa (majhen in priložnostni vzorec, brez standardiziranega
vprašalnika, brez merjenja časa po nalogah) obravnavamo v
razdelku~\ref{sec:groznje}.

% =====================================================================
\section{Simulacija uporabe}
\label{sec:simulacija}

Za ponazoritev delovanja rešitve povzemamo tipičen potek uporabe. Uporabnik je
ob vstopu usmerjen na glavno stran (slika~\ref{fig:glavna-stran}), kjer sta
izpostavljeni prijava in registracija ter pregled novih in vseh objavljenih
oglasov. Registriran uporabnik lahko objavi oglas prek vodene forme
(slika~\ref{fig:objava-oglasa}), ovrednoti poljubno vozilo, všečka oglase in si
ogleda svoj profil z všečkanimi vozili (slika~\ref{fig:liked-cars}).

\begin{figure}[H]
  \centering
  \includegraphics[width=0.9\textwidth]{glavna-stran.png}
  \caption{Glavna stran prototipa}
  \label{fig:glavna-stran}
\end{figure}
\begin{figure}[H]
  \centering
  \includegraphics[width=0.9\textwidth]{objava-oglasa.png}
  \caption{Vodena forma za objavo oglasa}
  \label{fig:objava-oglasa}
\end{figure}
\begin{figure}[H]
  \centering
  \includegraphics[width=0.9\textwidth]{liked-cars.png}
  \caption{Stran uporabnika s prikazanimi všečkanimi vozili}
  \label{fig:liked-cars}
\end{figure}

\paragraph{Primerjava ocen z obstoječimi portali.}
Za ilustracijo smo delovanje sistema primerjali s ponudbo na dveh spletnih
portalih na treh vozilih (tabela~\ref{tab:vrednotenje}). Sistematično in
kvantitativno oceno natančnosti podaja razdelek~\ref{sec:natancnost}; ta
primerjava je zgolj kvalitativna ponazoritev.

\begin{table}[htbp]
  \centering
  \caption{Primerjava ocenjenih cen s ponudbo na spletnih portalih (v evrih).}
  \label{tab:vrednotenje}
  \setlength{\tabcolsep}{5pt}
  \begin{tabular}{lccc}
    \toprule
    & \textbf{Audi A3} & \textbf{BMW ser.\ 2} & \textbf{Citroën Berlingo} \\
    \midrule
    slov.\ portal (razpon)     & 4.000--6.000  & 9.000--12.000  & 7.500--10.000 \\
    nem.\ portal (razpon)      & 4.500--9.000  & pribl.\ 11.100 & 8.000--12.000 \\
    \midrule
    Sistem --- dobra cena      & 4.500  & 7.600   & 7.600  \\
    Sistem --- sredinska cena  & 6.350  & 9.300   & 9.941  \\
    Sistem --- slaba cena      & 7.600  & 10.500  & 12.100 \\
    Raznolikost                & $49\,\%$ & $31\,\%$ & $45\,\%$ \\
    \bottomrule
  \end{tabular}
\end{table}

Pri vozilih z nizko raznolikostjo ponudbe (Audi~A3, BMW serije~2) se sredinska
ocena dobro ujema s cenami na obeh portalih. Pri modelu Citroën Berlingo je
raznolikost bistveno višja, zato je razpon med dobro in slabo oceno širši,
sredinska ocena pa manj zanesljiva --- kar je skladno z ugotovitvijo iz
razdelka~\ref{sec:natancnost}, da napaka strmo narašča z raznolikostjo ponudbe.
Pri interpretaciji je treba upoštevati, da sistem ocenjuje pričakovano prodajno
in ne oglaševano ceno.

% =====================================================================
\section{Grožnje veljavnosti}
\label{sec:groznje}

\paragraph{Notranja veljavnost.}
Kot merilo resnice pri merjenju natančnosti uporabljamo \emph{oglaševano} ceno,
ne dejanske prodajne cene, ki ni javno dostopna. Namerni popust, s katerim
sistem cilja pričakovano prodajno ceno, je zato ocenjen in ne umerjen na
dejanskih prodajah. Isti fizični oglas se lahko pojavi pri več virih; prečno
preverjanje izloči zapise z istim izvornim oglasom, ne more pa zaznati
podvojitev med viri.

\paragraph{Zunanja veljavnost.}
Podatki pokrivajo $10$ znamk in pretežno slovenski trg (z manjšim nemškim
vzorcem); znamka z največ oglasi na trgu (okoli $8000$ oglasov) v šifrantu ni
zajeta in je iz analize izpuščena. Uporabniški preizkus je temeljil na majhnem
in priložnostnem vzorcu.

\paragraph{Konstruktna veljavnost.}
Pojem \enquote{tržna vrednost} operacionaliziramo kot uteženo povprečje
oglaševanih cen primerljivih vozil, znižano za pričakovani pogajalski popust;
\enquote{raznolikost ponudbe} pa kot relativno širino interkvartilnega razpona
teh cen. Drugačna operacionalizacija bi lahko dala drugačne rezultate.

\paragraph{Zanesljivost.}
Meritve latence in obremenitve so bile izvedene na deljenem produkcijskem
strežniku, kjer lahko na rezultate vpliva ozadje. Meritve natančnosti so
ponovljive nad fiksnim posnetkom baze.

% =====================================================================
\section{Povzetek ugotovitev}
\label{sec:povzetek-analize}

Tabela~\ref{tab:rv-povzetek} povzema odgovore na raziskovalna vprašanja glede
na vnaprej postavljena merila.

\begin{table}[htbp]
  \centering
  \caption{Povzetek odgovorov na raziskovalna vprašanja.}
  \label{tab:rv-povzetek}
  \begin{tabularx}{\textwidth}{@{}l X c@{}}
    \toprule
    \textbf{RV} & \textbf{Glavni rezultat} & \textbf{Sodba} \\
    \midrule
    RV1 natančnost & MdAPE $13{,}5\,\%$ ($9$--$10\,\%$ za dobro pokrita vozila),
                     bistveno bolje od naivne ocene ($37{,}8\,\%$), primerljivo z
                     drugimi metodami & izpolnjeno \\
    RV2 odvisnost  & napaka pada z velikostjo vzorca (ustali se pri $n\approx20$),
                     narašča s starostjo, kilometrino in raznolikostjo ponudbe
                     ($7{,}5\,\%$ do $29{,}7\,\%$) & pojasnjeno \\
    RV3 hitrost    & vrednotenje $2$~ms (p50), $13{,}5$~ms (p99); indeks odpravi
                     odvisnost od velikosti baze; \res{meja X~z/s} & \res{izpolnjeno} \\
    RV4 uporabnost & vse naloge dokončane brez pomoči; \res{X} težav odpravljenih & izpolnjeno \\
    \bottomrule
  \end{tabularx}
\end{table}

\paragraph{Sodba o zastavljenih izhodiščih.}
Rezultati \emph{potrjujejo} izhodišča, ki smo si jih zastavili na začetku. Sistem na podlagi štirih
atributov vozila vrne oceno vrednosti z mediano odstotne napake $13{,}5\,\%$
(brez namernega popusta), pri dobro pokritih vozilih pa $9$--$10\,\%$, kar je
znotraj zastavljenega praga $15\,\%$ in bistveno bolje od naivne metode.
Predlagani pristop se med primerljivo natančnimi metodami ne uvršča na prvo
mesto --- preprost $k$-najbližjih sosedov doseže $11{,}2\,\%$ --- je pa edini,
ki poleg točkovne ocene vrne tudi razložljiv cenovni razpon in indeks
zaupanja, ta pa se izkaže za dober pokazatelj dejanske zanesljivosti
posamezne ocene. Glavne omejitve so odvisnost od enega prevladujočega
podatkovnega vira, nižja natančnost pri starejših vozilih in modelih z
razpršeno ponudbo ter merilo resnice, ki temelji na oglaševanih in ne
prodajnih cenah. Konkretne smeri za izboljšave --- vključno s prehodom na
metode strojnega učenja, ki jih ta analiza dodatno utemeljuje --- navajamo v
poglavju~\ref{ch:sklep}.

%
%  Številke so vpisane iz izhodov skript v CarStockMarketEvaluation/analysis/.
%  Preostale oznake \res{...} (rdeče) je treba še izpolniti:
%    - obremenitveni test (benchmarks/loadtest.py)
%    - sledljivost testov, Jest coverage, Lighthouse
%    - tabeli udeležencev in najdenih težav
%  Slike (PDF) iz analysis/out/ prekopiraj v mapo slike/.
% =====================================================================

\newcommand{\res}[1]{\textcolor{red}{\textbf{[#1]}}}

\chapter{Analiza in vrednotenje rešitve}
\label{ch:analiza}

V tem poglavju sistematično pregleadamo in analiziramo razvito rešitev. Najprej iz zahtev in
izhodišč, ki smo si jih zastavili na začetku, izpeljemo raziskovalna vprašanja in za vsako določimo merilo
uspešnosti (razdelek~\ref{sec:rq}). Nato opišemo merilno okolje in postopke
(razdelek~\ref{sec:metodologija}) ter natančno definiramo vse uporabljene mere
(razdelek~\ref{sec:mere}). Sledijo rezultati: značilnosti zbranih
podatkov~(\ref{sec:podatki}), natančnost ocene
vrednosti~(\ref{sec:natancnost}), občutljivost algoritma na
parametre~(\ref{sec:obcutljivost}), latenca posameznega
vrednotenja~(\ref{sec:latenca}) in obnašanje pod
obremenitvijo~(\ref{sec:obremenitev}). Zatem preverimo funkcionalne in
nefunkcionalne zahteve (\ref{sec:funkcionalne}, \ref{sec:nefunkcionalne}) ter
predstavimo uporabniški preizkus~(\ref{sec:uporabniki}) in simulacijo
uporabe~(\ref{sec:simulacija}). Poglavje sklenemo z obravnavo groženj
veljavnosti~(\ref{sec:groznje}) in povzetkom ugotovitev glede na
zastavljena izhodišča~(\ref{sec:povzetek-analize}).

% =====================================================================
\section{Raziskovalna vprašanja in merila uspešnosti}
\label{sec:rq}

Na začetku diplomskega dela (poglavje~\ref{ch:uvod}) smo si zastavili, da je
mogoče uporabniku na podlagi zbranih tržnih podatkov in primerjave s podobnimi
vozili ponuditi \emph{uporabno} in \emph{dovolj natančno} oceno tržne vrednosti
vozila in da vključitev takega sistema izboljša uporabniško izkušnjo. Iz tega
izhodišča in iz nefunkcionalnih zahtev (razdelek~\ref{sec:zahteve}) izpeljemo
štiri raziskovalna vprašanja (RV1--RV4) in za vsako določimo merilo uspešnosti.

\emph{Mediana absolutne odstotne napake} (MdAPE) je osrednja mera natančnosti v
tem poglavju; nanjo se sklicujejo merila uspešnosti in vsi nadaljnji razdelki.
Izračunamo jo tako, da za vsak testni primer določimo absolutno odstotno napako,
torej za koliko odstotkov se ocenjena cena razlikuje od dejanske, nato pa čez
vse primere vzamemo mediano teh vrednosti. MdAPE $=15\,\%$ tako pomeni, da je pri
polovici vozil ocena znotraj $15\,\%$ od dejanske cene, pri drugi polovici pa je
odstopanje večje; manjša kot je vrednost, natančnejše so ocene. Mediano
uporabljamo namesto povprečja, ker je manj občutljiva na posamezne skrajne
napake. Natančne definicije MdAPE in preostalih mer so v
razdelku~\ref{sec:mere}.

Raziskovalna vprašanja so:

\begin{itemize}
  \item \textbf{RV1 --- natančnost ocene.} Kako natančno sistem oceni tržno
        vrednost vozila v primerjavi z dejanskimi oglasi na trgu?
        \emph{Merilo:} MdAPE je pri vozilih z vsaj $10$ primerljivimi
        vozili manjša od $15\,\%$, ocena pa je bistveno natančnejša od
        preproste (naivne) referenčne metode.
  \item \textbf{RV2 --- odvisnost natančnosti od podatkov.} Od česa je
        natančnost ocene najbolj odvisna (velikost vzorca primerljivih vozil,
        razpršenost ponudbe, segment vozila)?
  \item \textbf{RV3 --- hitrost in razširljivost.} Ali sistem zadošča
        nefunkcionalnim zahtevam glede odzivnosti tudi pod sočasno
        obremenitvijo, in kje je zgornja meja trenutne postavitve?
  \item \textbf{RV4 --- uporabnost.} Ali ciljni uporabniki znajo brez navodil
        opraviti glavne naloge (registracija, objava oglasa, vrednotenje)?
\end{itemize}
% =====================================================================
\section{Metodologija in merilno okolje}
\label{sec:metodologija}

\paragraph{Merilno okolje.}
Meritve hitrosti (razdelka~\ref{sec:latenca} in~\ref{sec:obremenitev}) so
izvedene v \emph{produkcijskem} okolju, torej na navideznem zasebnem strežniku
z dvema procesorskima jedroma in $4$~GB pomnilnika ter s tehnologijami in
njihovimi različicami, opisanimi v poglavju~\ref{ch:razvoj}
(razdelek~\ref{sec:namestitev}). Analitične in merilne skripte tečejo nad isto
podatkovno bazo kot je na produkciji.

\paragraph{Posnetek podatkov.}
Meritve v tem poglavju temeljijo na posnetku podatkovne baze iz septembra
$2026$, ki obsega približno $86.000$ normaliziranih oglasov
(razdelek~\ref{sec:podatki}). Ker se posnetek med analizo ne spreminja, so
rezultati ponovljivi.

\paragraph{Orodja.}
Za meritve in analizo v tem poglavju smo napisali namenske skripte, ločene od
produkcijske kode; ena od njih je dokumentirana kopija ocenjevalnega
algoritma, parametrizirana za eksperimente. Statistično obdelavo in izris
grafov smo izvedli s knjižnicama \texttt{NumPy}~\cite{harris2020numpy} in
\texttt{Matplotlib}~\cite{hunter2007matplotlib}.

\paragraph{Natančnost.}
Natančnost ocene (razdelek~\ref{sec:natancnost}) merimo z naključno izbranim
normaliziranim oglasom: njegove atribute (znamka, model, letnik, kilometrina)
uporabimo kot vhod v algoritem, sam oglas pa izločimo iz primerjalne množice. Nato izračunamo oceno in jo
primerjamo z \emph{dejansko oglaševano} ceno tega oglasa. Postopek ponovimo za
$3000$ naključnih oglasov.

% =====================================================================
\section{Uporabljene mere}
\label{sec:mere}

\paragraph{Latenca}
 je čas, ki preteče od trenutka tik pred oddajo zahtevka do
trenutka, ko odjemalec prejme celoten odgovor. Glede na to, kolikšen del te
poti zajamemo, ločimo tri vrste:
\begin{itemize}
  \item[(a)] \emph{latenca obdelave na strežniku} --- čas izvajanja funkcije
        za vrednotenje, od začetka do konca, brez omrežja;
  \item[(b)] \emph{latenca poizvedbe SQL} --- čas izvajanja poizvedbe v podatkovni bazi,
   od začetka do konca, brez omrežja;
  \item[(c)] \emph{latenca od konca do konca} (angl.\ \emph{end-to-end}) ---
        (a) skupaj z omrežnim obhodom med odjemalcem in strežnikom ter
        obdelavo zahtevka HTTP; to je čas, ki ga zazna uporabnik.
\end{itemize}
V razdelku~\ref{sec:latenca} merimo vrsti (a) in (b), v
razdelku~\ref{sec:obremenitev} pa vrsto (c).

\paragraph{Percentili.}
\textbf{Percentil} $p_k$ je vrednost, pod katero pade $k\,\%$ meritev. Poročamo
$p_{50}$ (mediana), $p_{95}$ in $p_{99}$.

\paragraph{Prepustnost, sočasnost, nasičenje.}
Obnašanje sistema pod obremenitvijo opišemo s tremi količinami:
\begin{itemize}
  \item \textbf{prepustnost} --- število uspešno obdelanih zahtevkov na
        sekundo;
  \item \textbf{sočasnost} --- število odjemalcev, ki hkrati držijo odprt
        zahtevek;
  \item \textbf{nasičenje} (\emph{koleno}) --- točka, v kateri dodatna
        sočasnost ne poveča več prepustnosti, latenca pa začne naraščati
        približno linearno.
\end{itemize}
Notranjo skladnost meritev preverimo z \textbf{Littlovim zakonom}~\cite{little1961},
\begin{equation}
  \label{eq:little}
  L = \lambda \cdot W,
\end{equation}
kjer je $L$ povprečno število zahtevkov v sistemu, $\lambda$ prepustnost in
$W$ povprečna latenca.

\paragraph{Napaka ocene.}
Za oglas $i$ z dejansko ceno $p_i$ in ocenjeno ceno $\hat p_i$ definiramo
absolutno odstotno napako
\begin{equation}
  \mathrm{APE}_i = \frac{\lvert \hat p_i - p_i \rvert}{p_i}.
\end{equation}
\textbf{MAPE} je povprečje $\mathrm{APE}_i$, \textbf{MdAPE} pa njihova mediana;
MdAPE je robustnejša na osamelce in je naša glavna mera~\cite{hyndman2006}.
Poročamo tudi \textbf{MAE} (povprečna absolutna napaka) in \textbf{RMSE}
(koren povprečne kvadratne napake) v evrih ter \emph{predznačeno} odstotno
napako, ki pokaže sistematičen odmik ocene navzgor ali navzdol.

\paragraph{Pokritost.}
\textbf{Pokritost} je delež poizvedb, za katere sistem vrne oceno, torej za
katere najde vsaj eno (oziroma vsaj $5$ oz.\ $10$) primerljivo vozilo.

\paragraph{Raznolikost ponudbe.}
\textbf{Raznolikost} (v izvorni kodi \texttt{ratio}) je relativna širina
cenovnega razpona primerljivih vozil,
\begin{equation}
  \label{eq:raznolikost}
  \mathrm{raznolikost} = \frac{P_{75} - P_{25}}{\hat{p}},
\end{equation}
kjer sta $P_{25}$ in $P_{75}$ (v kodi \texttt{goodPrice} in \texttt{badPrice})
uteženi $25.$ oziroma $75.$ percentil cen primerljivih vozil, $\hat{p}$
(v kodi \texttt{estimatedPrice}) pa ocenjena cena vozila. Vrednost izražamo v
odstotkih; enoten zapis uporabljamo v vseh tabelah tega poglavja.

% =====================================================================
\section{Značilnosti zbranih podatkov}
\label{sec:podatki}
\paragraph{Viri in obseg.}
Velika večina oglasov je bila zajeta z največjega slovenskega spletnega portala
za prodajo rabljenih vozil (razdelek~\ref{sec:zajem-podatkov}); manjši del
predstavlja vzorec z večjega nemškega oglasnega portala, ki smo ga uporabili za
preverjanje delovanja cevovoda na drugačni strukturi vira. Zajem je zajel
javno dostopne oglase za nekomercialno raziskovalno rabo, z omejeno hitrostjo
zahtevkov in brez shranjevanja kontaktnih ali osebnih podatkov prodajalcev.
Po normalizaciji baza obsega $85.846$ oglasov ($83.493$ s slovenskega in
$2.353$ z nemškega portala) za $10$ znamk vozil; razčlenitev po znamki je v
tabeli~\ref{tab:podatki-znamke}.

\begin{table}[htbp]
  \centering
  \caption{Število normaliziranih oglasov po znamki.}
  \label{tab:podatki-znamke}
  % vir: analysis/out/dataset_per_brand.csv
  \begin{tabular}{lr lr}
    \toprule
    \textbf{Znamka} & \textbf{Oglasi} & \textbf{Znamka} & \textbf{Oglasi} \\
    \midrule
    Volkswagen    & 21.191 & Peugeot     & 7.513 \\
    Audi          & 11.842 & Citroën     & 6.729 \\
    BMW           & 10.724 & Ford        & 5.783 \\
    Renault       & 9.615  & Toyota      & 2.648 \\
    Mercedes-Benz & 9.355  & Alfa Romeo  & 446 \\
    \bottomrule
  \end{tabular}
\end{table}

\paragraph{Porazdelitve ključnih atributov.}
Slika~\ref{fig:porazdelitve} prikazuje porazdelitve oglaševanih cen, prevožene
kilometrine in letnika. Porazdelitev cen je izrazito desno asimetrična (dolg
rep dragih vozil), zato jo prikazujemo na logaritemski osi. Mediana cene je
$14.890$~EUR (povprečje $18.137$~EUR), mediana kilometrine $159.950$~km,
mediana letnika pa $2017$; podrobnosti so v tabeli~\ref{tab:podatki-porazdelitve}.
Zaradi berljivosti so na grafih odrezane skrajne vrednosti (nekaj oglasov ima
očitno napačno ceno, kilometrino ali letnik); njihovo število je označeno v
kotu vsakega grafa, celoten razpon pa je razviden iz stolpca \emph{maks.}\ v
tabeli~\ref{tab:podatki-porazdelitve} in iz obravnave osamelcev v
tabeli~\ref{tab:osamelci}.

\begin{table}[htbp]
  \centering
  \caption{Povzetek porazdelitev ključnih numeričnih atributov.}
  \label{tab:podatki-porazdelitve}
  % vir: analysis/out/dataset_distribution_summary.json
  \setlength{\tabcolsep}{5pt}
  \begin{tabular}{lrrrrr}
    \toprule
    & \textbf{$p_{5}$} & \textbf{mediana} & \textbf{povprečje} & \textbf{$p_{95}$} & \textbf{maks.} \\
    \midrule
    cena [EUR]        & 1.550  & 14.890  & 18.137  & 49.014  & 374.000 \\
    kilometrina [km]  & 26.000 & 159.950 & 169.464 & 334.000 & 1.974.001 \\
    letnik            & 2005   & 2017    & 2016    & 2024    & 2029 \\
    moč [kW]          & 56     & 110     & 117     & 221     & 920 \\
    \bottomrule
  \end{tabular}
\end{table}

\begin{figure}[htbp]
  \centering
  \begin{subfigure}{0.48\textwidth}
    \includegraphics[width=\textwidth]{dataset_price_hist.pdf}
    \caption{oglaševana cena (log)}
  \end{subfigure}\hfill
  \begin{subfigure}{0.48\textwidth}
    \includegraphics[width=\textwidth]{dataset_km_hist.pdf}
    \caption{prevoženi kilometri}
  \end{subfigure}
  \\[1ex]
  \begin{subfigure}{0.48\textwidth}
    \includegraphics[width=\textwidth]{dataset_year_hist.pdf}
    \caption{letnik prve registracije}
  \end{subfigure}\hfill
  \begin{subfigure}{0.48\textwidth}
    \includegraphics[width=\textwidth]{dataset_brand_bar.pdf}
    \caption{oglasi po znamki}
  \end{subfigure}
  \caption{Porazdelitve ključnih atributov v normalizirani bazi.}
  \label{fig:porazdelitve}
  % vir: analysis/out/dataset_*_hist.pdf, dataset_brand_bar.pdf
\end{figure}

\paragraph{Pokritost po segmentih.}
Za pošteno oceno mora imeti algoritem dovolj primerljivih vozil v ožjem
segmentu (ista znamka in model, podoben letnik in kilometrina).
Slika~\ref{fig:pokritost} prikazuje toplotno karto za znamki Renault in
Volkswagen: vrstice so modeli, stolpci letniki, barva pa
število oglasov v tej kombinaciji. Vidno je, da so pogosti modeli srednjih
letnikov (npr.\ Volkswagen Golf ali Renault Clio) gosto pokriti, redki modeli
in skrajni letniki pa le redko, kar neposredno pojasni razlike v zanesljivosti
ocen (razdelek~\ref{sec:natancnost}).

\begin{figure}[htbp]
  \centering
  \includegraphics[width=\textwidth]{dataset_coverage_heatmap.pdf}
  \caption{Pokritost: število oglasov po modelu in letniku za znamki
  Renault in Volkswagen (po ena znamka na vrstico).}
  \label{fig:pokritost}
  % vir: analysis/out/dataset_coverage_heatmap.pdf
\end{figure}

\paragraph{Osamelci in kakovost podatkov.}
Postopek normalizacije trenutno ne izvaja filtriranja osamelcev, zato baza
vsebuje tudi vozila z očitno napačnimi vrednostmi.
Tabela~\ref{tab:osamelci} prikazuje njihov delež. Kar $10{,}7\,\%$ oglasov ima
vsaj eno vrednost zunaj smiselnih mej, največ pri kilometrini ($9{,}9\,\%$;
najvišja zabeležena vrednost je skoraj dva milijona kilometrov) in letniku
($8{,}1\,\%$; nekaj oglasov ima celo letnik v prihodnosti). Preprost filter po
interkvartilnem razponu (IQR) na ceni znotraj skupine
$(\text{znamka},\text{model})$ bi odstranil dodatnih $2{,}2\,\%$ oglasov.
Vpliv doslednega filtriranja osamelcev na natančnost obravnavamo v
razdelku~\ref{sec:obcutljivost}, uvedbo takega filtra pa v
poglavju~\ref{ch:sklep}.

\begin{table}[H]
  \centering
  \caption{Osamelci v normalizirani bazi ($85.846$ oglasov).}
  \label{tab:osamelci}
  % vir: analysis/out/dataset_outliers.csv
  \begin{tabularx}{\textwidth}{@{}X r r@{}}
    \toprule
    \textbf{Merilo} & \textbf{Št.\ vrstic} & \textbf{Delež} \\
    \midrule
    cena zunaj $[500,\ 200.000]$~EUR         & 427   & $0{,}5\,\%$ \\
    kilometrina zunaj $[1000,\ 500.000]$~km  & 8.495 & $9{,}9\,\%$ \\
    letnik zunaj $[1990,\ 2027]$             & 6.979 & $8{,}1\,\%$ \\
    vsaj ena vrednost zunaj mej              & 9.169 & $10{,}7\,\%$ \\
    dodatno po filtru IQR na ceni            & 1.878 & $2{,}2\,\%$ \\
    \bottomrule
  \end{tabularx}
\end{table}

% =====================================================================
\section{Natančnost ocene vrednosti}
\label{sec:natancnost}

To je osrednja meritev diplomskega dela: neposredno preverja zastavljeno
izhodišče o \emph{dovolj natančni} oceni (RV1). Protokol je opisan v
razdelku~\ref{sec:metodologija}, rezultate pa izdela skripta
\texttt{analysis/accuracy\_cv.py} nad vzorcem $3000$ oglasov (od tega jih je
$2986$ dobilo oceno).

\paragraph{Merilo resnice in namerni odmik ocene.}
Kot merilo resnice uporabljamo \emph{oglaševano} ceno, saj dejanske prodajne
cene niso javno dostopne. Sistem pa namenoma ne ocenjuje oglaševane cene,
temveč \emph{pričakovano prodajno ceno}: ker se rabljena vozila v pogajanjih
praviloma prodajo pod oglaševano ceno, oceno zniža s stopenjskim popustom
($9$--$15\,\%$ glede na cenovni razred). Zato je pri primerjavi z oglaševano
ceno pričakovan rahel \emph{negativni} odmik ocene, kar je zaželeno in ne
pomeni slabše ocene. To lastnost upoštevamo pri branju vseh rezultatov v
nadaljevanju: napako merimo dvakrat --- proti oglaševani ceni z vključenim
popustom (kar dejansko vrne sistem) in brez popusta (čista primerjalna
ocena).

\paragraph{Skupni rezultat.}
Brez popusta znaša MdAPE predlaganega pristopa $13{,}5\,\%$, MAPE $23{,}8\,\%$,
MAE $2798$~EUR in RMSE $5151$~EUR; z vključenim namernim popustom je MdAPE
proti oglaševani ceni $14{,}5\,\%$ (mediana predznačene napake $-6{,}8\,\%$).
Pokritost je $99{,}5\,\%$ pri $n \ge 1$ in $95{,}7\,\%$ pri $n \ge 10$.
Slika~\ref{fig:pred-vs-actual} prikazuje napovedano ceno proti dejanski
($R^2 = 0{,}87$), slika~\ref{fig:ape-hist} pa porazdelitev odstotne napake.

\begin{figure}[htbp]
  \centering
  \begin{subfigure}{0.48\textwidth}
    \includegraphics[width=\textwidth]{accuracy_pred_vs_actual.pdf}
    \caption{napovedano proti dejanskemu}
    \label{fig:pred-vs-actual}
  \end{subfigure}\hfill
  \begin{subfigure}{0.48\textwidth}
    \includegraphics[width=\textwidth]{accuracy_ape_hist.pdf}
    \caption{porazdelitev napake}
    \label{fig:ape-hist}
  \end{subfigure}
  \caption{Natančnost ocene vrednosti (leave-one-out, brez popusta).}
  % vir: analysis/out/accuracy_pred_vs_actual.pdf, accuracy_ape_hist.pdf
\end{figure}

\paragraph{Primerjava z osnovnimi modeli.}
Da presodimo, ali je predlagani pristop boljši od trivialnega, ga primerjamo s
štirimi referenčnimi modeli na istem protokolu in isti množici kandidatov
(tabela~\ref{tab:baselines}, slika~\ref{fig:baselines}): B0~globalna mediana
cene za dani model, B1~navadno povprečje cen v filtrirnem oknu, B1m~mediana cen
v oknu, B2~linearna regresija $\text{cena} \sim \text{letnik} + \text{km}$,
B3~metoda $k$~najbližjih sosedov ($k = 10$) po standardiziranih letniku in
kilometrini, ter B4~naš pristop (uteženo povprečje s podobnostnim jedrom, brez
popusta).

\begin{table}[htbp]
  \centering
  \caption{Primerjava natančnosti z osnovnimi modeli (leave-one-out).
  MdAPE in MAPE sta v odstotkih, MAE in RMSE v evrih.}
  \label{tab:baselines}
  % vir: analysis/out/accuracy_overall.csv
  \setlength{\tabcolsep}{5pt}
  \begin{tabular}{l r r r r}
    \toprule
    \textbf{Model} & \textbf{MdAPE} & \textbf{MAPE} & \textbf{MAE} & \textbf{RMSE} \\
    \midrule
    B0 globalna mediana       & 37{,}8 & 95{,}7 & 7395 & 10883 \\
    B1 povprečje v oknu       & 14{,}2 & 24{,}9 & 2930 & 5333 \\
    B1m mediana v oknu        & 13{,}1 & 22{,}5 & 2877 & 5561 \\
    B2 linearna regresija     & 12{,}9 & 22{,}4 & 2677 & 5032 \\
    B3 $k$-najbližjih sosedov & \textbf{11{,}2} & 21{,}5 & 2497 & 4883 \\
    B4 naš pristop            & 13{,}5 & 23{,}8 & 2798 & 5151 \\
    \bottomrule
  \end{tabular}
\end{table}

\begin{figure}[htbp]
  \centering
  \includegraphics[width=0.7\textwidth]{accuracy_baselines.pdf}
  \caption{MdAPE osnovnih modelov in predlaganega pristopa.}
  \label{fig:baselines}
  % vir: analysis/out/accuracy_baselines.pdf
\end{figure}

Predlagani pristop je \emph{bistveno} boljši od naivne globalne mediane
($13{,}5\,\%$ proti $37{,}8\,\%$) in primerljiv s preostalimi metodami, ki se
gibljejo med $11$ in $14\,\%$. Ni pa najboljši: metoda $k$~najbližjih sosedov
doseže $11{,}2\,\%$, tudi navadna mediana istega okna ($13{,}1\,\%$) je
malenkost natančnejša. To je pričakovano za zasnovo, katere cilj je
\emph{razložljivost}: za razliko od naštetih metod predlagani pristop poleg
točkovne ocene vrne tudi cenovni razpon in indeks zaupanja, uteži pa so
neposredno interpretabilne. Rezultat hkrati potrjuje smiselnost prehoda na
metode strojnega učenja, ki jo predlagamo v poglavju~\ref{ch:sklep}: če že
preprost $k$-NN preseže sedanji pristop, je od zmogljivejših modelov mogoče
pričakovati še boljše rezultate.

\paragraph{Odvisnost od velikosti vzorca in segmenta (RV2).}
Slika~\ref{fig:mdape-n} prikazuje MdAPE glede na število primerljivih vozil
$n$: napaka je pri $n < 10$ okoli $20\,\%$, pri $n \ge 20$ pa se ustali pri
približno $13\,\%$. To empirično podpira izbiro konstante $20$ v indeksu
zaupanja (enačba~\ref{eq:confidence}) --- pri toliko primerljivih vozilih je
ocena že blizu svoje najboljše natančnosti. Slika~\ref{fig:mdape-brand} prikazuje MdAPE po znamki: najbolje se
odrežejo pogosti, gosto pokriti modeli (Audi $9{,}8\,\%$, Toyota $10{,}8\,\%$,
Volkswagen $12{,}1\,\%$), najslabše pa Mercedes-Benz ($24{,}9\,\%$), pri
katerem se pri normalizaciji več različnih modelov zlije v isto kanonično
ime, kar poslabša izbor primerljivih vozil.
Tabela~\ref{tab:natancnost-segmenti} kaže še razčlenitev po starosti,
kilometrini in raznolikosti ponudbe.

\begin{figure}[htbp]
  \centering
  \begin{subfigure}{0.48\textwidth}
    \includegraphics[width=\textwidth]{accuracy_mdape_vs_n.pdf}
    \caption{glede na velikost vzorca $n$}
    \label{fig:mdape-n}
  \end{subfigure}\hfill
  \begin{subfigure}{0.48\textwidth}
    \includegraphics[width=\textwidth]{accuracy_mdape_by_brand.pdf}
    \caption{po znamki}
    \label{fig:mdape-brand}
  \end{subfigure}
  \caption{Odvisnost natančnosti ocene od podatkov.}
  % vir: analysis/out/accuracy_mdape_vs_n.pdf, accuracy_mdape_by_brand.pdf
\end{figure}

\begin{table}[htbp]
  \centering
  \caption{Natančnost predlaganega pristopa po segmentih (MdAPE).}
  \label{tab:natancnost-segmenti}
  % vir: analysis/out/accuracy_by_age.csv, accuracy_by_km.csv, accuracy_by_ratio.csv
  \begin{tabularx}{\textwidth}{@{}l X r r@{}}
    \toprule
    \textbf{Razčlenitev} & \textbf{Koš} & \textbf{Primeri} & \textbf{MdAPE [\%]} \\
    \midrule
    starost [leta] & 0--3   & 367 & 9{,}0 \\
                   & 4--7   & 921 & 10{,}6 \\
                   & 8--12  & 932 & 14{,}5 \\
                   & $13+$  & 766 & 22{,}9 \\
    \midrule
    kilometrina    & $<50.000$        & 280  & 9{,}6 \\
                   & $50.000$--$100.000$  & 449  & 11{,}9 \\
                   & $100.000$--$150.000$ & 634  & 12{,}4 \\
                   & $150.000$--$200.000$ & 629  & 13{,}7 \\
                   & $>200.000$       & 1021 & 16{,}7 \\
    \midrule
    raznolikost    & $0$--$0{,}10$    & 88   & 7{,}5 \\
                   & $0{,}10$--$0{,}25$ & 1002 & 8{,}4 \\
                   & $0{,}25$--$0{,}50$ & 1613 & 15{,}5 \\
                   & $>0{,}50$        & 447  & 29{,}7 \\
    \bottomrule
  \end{tabularx}
\end{table}

Vzorec pokaže jasne zakonitosti: napaka narašča s starostjo vozila ($9\,\%$ pri
najnovejših, $23\,\%$ pri vozilih, starejših od $13$~let), s prevoženo
kilometrino in --- najmočneje --- z \emph{raznolikostjo ponudbe}. Pri modelih z
enotno ponudbo (raznolikost pod $0{,}1$) je MdAPE le $7{,}5\,\%$, pri razpršeni
ponudbi (nad $0{,}5$) pa $29{,}7\,\%$. To potrjuje, da je raznolikost, ki jo
sistem uporabniku tudi prikaže, dober vgrajen pokazatelj zanesljivosti
posamezne ocene.

\paragraph{Vpliv popusta.}
Slika~\ref{fig:popust} prikazuje predznačeno napako z in brez popusta. Brez
popusta je ocena rahlo optimistična (mediana $+3{,}8\,\%$ nad oglaševano ceno),
z namernim popustom pa se premakne na $-6{,}8\,\%$ pod oglaševano ceno. Ta
premik je skladen z zasnovo sistema, ki cilja pričakovano prodajno in ne
oglaševane cene; dokončno umerjanje popusta bi bilo mogoče šele s podatki o
dejanskih prodajah (razdelek~\ref{sec:groznje}).

\begin{figure}[htbp]
  \centering
  \includegraphics[width=0.7\textwidth]{accuracy_discount_bias.pdf}
  \caption{Predznačena napaka proti oglaševani ceni, z in brez namernega popusta.}
  \label{fig:popust}
  % vir: analysis/out/accuracy_discount_bias.pdf
\end{figure}

\paragraph{Časovni razrez.}
Pri napovedovanju najnovejših $10\,\%$ oglasov zgolj iz starejših $90\,\%$
znaša MdAPE predlaganega pristopa $13{,}0\,\%$ (tabela~\ref{tab:temporal}), kar
je celo malenkost bolje od rezultata leave-one-out in potrjuje, da ocena ne
izhaja iz uhajanja informacije.

\begin{table}[htbp]
  \centering
  \caption{Natančnost pri časovnem razrezu (napoved najnovejših $10\,\%$).
  MdAPE in MAPE sta v odstotkih, MAE in RMSE v evrih.}
  \label{tab:temporal}
  % vir: analysis/out/accuracy_temporal.csv
  \setlength{\tabcolsep}{5pt}
  \begin{tabular}{l r r r r}
    \toprule
    \textbf{Model} & \textbf{MdAPE} & \textbf{MAPE} & \textbf{MAE} & \textbf{RMSE} \\
    \midrule
    B0 globalna mediana       & 35{,}9 & 74{,}5 & 7596 & 11222 \\
    B1m mediana v oknu        & 12{,}8 & 19{,}7 & 3088 & 5805 \\
    B3 $k$-najbližjih sosedov & 10{,}5 & 18{,}2 & 2667 & 4890 \\
    B4 naš pristop            & 13{,}0 & 20{,}6 & 2994 & 5347 \\
    \bottomrule
  \end{tabular}
\end{table}

% =====================================================================
\section{Občutljivost algoritma na parametre}
\label{sec:obcutljivost}

Ocenjevalni algoritem vsebuje več hevristično določenih konstant: dovoljeno
odstopanje kilometrine ($\pm 40.000$~km) in letnika ($\pm 3$~leta), uteži
delnih ocen ($0{,}6$ in $0{,}4$), oblika padajoče funkcije in stopenjski
popust. V tem razdelku vsako spreminjamo ob nespremenjenih ostalih in merimo
MdAPE ter pokritost (vzorec $2000$ oglasov, leave-one-out). Rezultate izdela
skripta \texttt{analysis/sensitivity.py}.

\paragraph{Okno primerljivosti.}
Tabela~\ref{tab:sens-okno} in slika~\ref{fig:sens-km} kažeta, da je algoritem
na širino \emph{kilometrinskega} okna skoraj neobčutljiv: MdAPE se med
$\pm 10.000$~km in $\pm 80.000$~km giblje le med $13{,}2$ in $13{,}6\,\%$.
Nasprotno je odvisen od \emph{letnikovega} okna: pri $\pm 1$~letu je MdAPE
$12{,}6\,\%$, pri $\pm 5$~letih pa $15{,}2\,\%$. Ožje okno da bolj primerljiva
vozila, a manjši vzorec (nižja pokritost); gre za znano razmerje med
pristranskostjo in varianco~\cite{bias-variance}. Toplotna karta
(slika~\ref{fig:sens-heatmap}) potrjuje, da natančnost določa skoraj izključno
letnikovo okno. Sedanja izbira $\pm 3$~leta je razumen kompromis med
natančnostjo in pokritostjo; morebitna zožitev na $\pm 2$~leti bi MdAPE
znižala za približno $0{,}3$~odstotne točke ob še vedno visoki pokritosti.

\begin{table}[htbp]
  \centering
  \caption{Občutljivost na širino okna primerljivosti. Pokritost je delež
  poizvedb z vsaj $5$ primerljivimi vozili.}
  \label{tab:sens-okno}
  % vir: analysis/out/sensitivity_km.csv, sensitivity_year.csv
  \setlength{\tabcolsep}{5pt}
  \begin{tabular}{l r r r}
    \toprule
    \textbf{Parameter} & \textbf{MdAPE [\%]} & \textbf{Pokritost [\%]} & \textbf{Mediana $n$} \\
    \midrule
    kilometrina $\pm 10.000$~km & 13{,}2 & 92{,}3 & 53 \\
    kilometrina $\pm 40.000$~km & 13{,}4 & 98{,}1 & 196 \\
    kilometrina $\pm 80.000$~km & 13{,}6 & 99{,}0 & 362 \\
    \midrule
    letnik $\pm 1$~leto  & 12{,}6 & 95{,}6 & 97 \\
    letnik $\pm 3$~leta  & 13{,}4 & 98{,}1 & 196 \\
    letnik $\pm 5$~let   & 15{,}2 & 98{,}5 & 262 \\
    \bottomrule
  \end{tabular}
\end{table}

\begin{figure}[htbp]
  \centering
  \begin{subfigure}{0.48\textwidth}
    \includegraphics[width=\textwidth]{sensitivity_km.pdf}
    \caption{okno kilometrine}
  \end{subfigure}\hfill
  \begin{subfigure}{0.48\textwidth}
    \includegraphics[width=\textwidth]{sensitivity_year.pdf}
    \caption{okno letnika}
  \end{subfigure}
  \caption{Občutljivost na širino okna primerljivosti (MdAPE in pokritost).}
  \label{fig:sens-km}
  % vir: analysis/out/sensitivity_km.pdf, sensitivity_year.pdf
\end{figure}

\begin{figure}[htbp]
  \centering
  \includegraphics[width=0.7\textwidth]{sensitivity_km_year_heatmap.pdf}
  \caption{MdAPE [\%] po mreži dovoljenih odstopanj kilometrine in letnika.}
  \label{fig:sens-heatmap}
  % vir: analysis/out/sensitivity_km_year_heatmap.pdf
\end{figure}

\paragraph{Uteži, oblika jedra in popust.}
Tabela~\ref{tab:sens-ostalo} povzema preostale parametre. Razmerje uteži med
kilometrino in letnikom na natančnost praktično ne vpliva ($13{,}4$ do
$13{,}6\,\%$ za vsa preizkušena razmerja), prav tako ne zamenjava linearne
padajoče funkcije z eksponentno ($13{,}4$ proti $13{,}3\,\%$). Edina parametra,
ki opazno premakneta rezultat, sta torej širina letnikovega okna in namerni
popust: ta predznačeno napako proti oglaševani ceni premakne iz $+3{,}8\,\%$ v
$-6{,}8\,\%$, MdAPE pa zviša z $13{,}4$ na $14{,}5\,\%$. Ta premik je posledica
zavestne odločitve, da sistem ocenjuje pričakovano prodajno in ne oglaševane
cene, in ne slabše ocene.

\begin{table}[htbp]
  \centering
  \caption{Občutljivost na uteži, obliko jedra in popust.}
  \label{tab:sens-ostalo}
  % vir: analysis/out/sensitivity_weights.csv, sensitivity_kernel.csv, sensitivity_discount.csv
  \begin{tabularx}{\textwidth}{@{}l X r >{\raggedleft\arraybackslash}p{2.3cm}@{}}
    \toprule
    \textbf{Parameter} & \textbf{Nastavitev} & \textbf{MdAPE [\%]} & \textbf{Predznačena napaka [\%]} \\
    \midrule
    uteži (km / letnik) & $0{,}5 / 0{,}5$              & 13{,}4 & $+3{,}7$ \\
    uteži (km / letnik) & $0{,}6 / 0{,}4$ (privzeto)   & 13{,}4 & $+3{,}8$ \\
    uteži (km / letnik) & $0{,}8 / 0{,}2$              & 13{,}6 & $+3{,}8$ \\
    \midrule
    jedro               & linearno (privzeto)         & 13{,}4 & $+3{,}8$ \\
    jedro               & eksponentno                 & 13{,}3 & $+3{,}8$ \\
    \midrule
    popust              & brez                        & 13{,}4 & $+3{,}8$ \\
    popust              & stopenjski (privzeto)        & 14{,}5 & $-6{,}8$ \\
    \bottomrule
  \end{tabularx}
\end{table}

\paragraph{Sklep o robustnosti.}
Algoritem je robusten: na večino svojih konstant se skoraj ne odziva, kar
pomeni, da njihova hevristična izbira ni sporna. Njegova natančnost je odvisna
predvsem od kakovosti in gostote podatkov (razdelek~\ref{sec:natancnost}), ne
od finega nastavljanja parametrov.

% =====================================================================
\section{Latenca posameznega vrednotenja}
\label{sec:latenca}

Merimo latenco vrste (a) in (b) iz razdelka~\ref{sec:mere}. Meritev izvede
\texttt{benchmarks/bench\_algo.py}: $2000$ neposrednih klicev funkcije nad
naključnimi realnimi vhodi (brez HTTP), po $100$ ogrevalnih klicih, ki se ne
merijo. Ločeno merimo čas poizvedbe SQL, čas točkovanja v Pythonu in čas
serializacije odgovora.

\paragraph{Ozko grlo: pregled celotne tabele.}
Prvotna izvedba poizvedbe ni imela ustreznega indeksa.
Izsek~\ref{lst:explain-before} prikazuje izpis \texttt{EXPLAIN (ANALYZE,
BUFFERS)}: načrt je \texttt{Seq Scan}, baza pri vsakem vrednotenju prebere
celotno tabelo in zavrže $85.273$ vrstic, ki ne ustrezajo pogojem
(\texttt{Rows Removed by Filter}). Ker so podatki v celoti v predpomnilniku
(\texttt{shared hit}, brez \texttt{read}), je čas stabilen, a raste linearno z
velikostjo tabele.

\begin{center}
\captionof{listing}{\texttt{EXPLAIN (ANALYZE, BUFFERS)} pred dodajanjem indeksa.}
\label{lst:explain-before}
\begin{minted}{text}
Seq Scan on normalized_market_listing
    (cost=0.00..3314.15 rows=12 width=97)
    (actual time=0.555..6.847 rows=573 loops=1)
  Filter: ((kilometers >= 18500) AND (kilometers <= 98500)
       AND (registration_year >= 2019) AND (registration_year <= 2025)
       AND (brand_name = 'Citroën') AND (model_name = 'C3'))
  Rows Removed by Filter: 85273
  Buffers: shared hit=1168
Execution Time: 6.880 ms
\end{minted}
\end{center}

\paragraph{Učinek sestavljenega indeksa.}
Po dodajanju sestavljenega indeksa na
$(\text{brand\_name},\ \text{model\_name},\ \text{registration\_year},\
\text{kilometers})$ se načrt spremeni v \texttt{Bitmap Index Scan}
(izsek~\ref{lst:explain-after}): baza prebere le nekaj strani indeksa in
pripadajoče vrstice, čas izvajanja pa pade z $6{,}88$~ms na $0{,}69$~ms.
Latenca poizvedbe SQL se v mediani zmanjša z $9{,}7$~ms na $2{,}0$~ms
(tabela~\ref{tab:indeks-pred-po}, slika~\ref{fig:indeks-pred-po}). Pomembneje
od trenutnega pohitrenja je, da postane čas poizvedbe neodvisen od velikosti
tabele: namesto $O(\text{tabela})$ postane $O(\text{zadetki})$.

\begin{table}[htbp]
  \centering
  \caption{Latenca poizvedbe SQL pred in po dodajanju indeksa, v milisekundah.}
  \label{tab:indeks-pred-po}
  % vir: benchmarks/report.md (results_algo_before/after.json)
  \begin{tabular}{lrrrr}
    \toprule
    & \textbf{p50} & \textbf{p95} & \textbf{p99} & \textbf{Povprečje} \\
    \midrule
    pred (Seq Scan)  & 9{,}7 & 21{,}0 & 24{,}8 & 11{,}3 \\
    po (Index Scan)  & 2{,}0 & 8{,}6  & 12{,}5 & 3{,}0 \\
    \bottomrule
  \end{tabular}
\end{table}

\begin{figure}[htbp]
  \centering
  \includegraphics[width=0.6\textwidth]{latency_index_before_after.pdf}
  \caption{Latenca poizvedbe SQL pred in po dodajanju sestavljenega indeksa.}
  \label{fig:indeks-pred-po}
  % vir: analysis/out/latency_index_before_after.pdf
\end{figure}

\begin{center}
\captionof{listing}{\texttt{EXPLAIN (ANALYZE, BUFFERS)} po dodajanju indeksa.}
\label{lst:explain-after}
\begin{minted}{text}
Bitmap Heap Scan on normalized_market_listing
    (cost=5.84..50.49 rows=12 width=97)
    (actual time=0.202..0.615 rows=573 loops=1)
  Recheck Cond: (brand_name = 'Citroën' AND model_name = 'C3'
       AND registration_year BETWEEN 2019 AND 2025
       AND kilometers BETWEEN 18500 AND 98500)
  Buffers: shared hit=195
  ->  Bitmap Index Scan on idx_nml_valuation
        (actual time=0.166..0.166 rows=573 loops=1)
        Buffers: shared hit=6
Execution Time: 0.688 ms
\end{minted}
\end{center}

\paragraph{Razčlenitev latence po komponentah.}
Tabela~\ref{tab:latenca-vrednotenja} in slika~\ref{fig:latency-breakdown}
prikazujeta latenco celotnega vrednotenja po dodajanju indeksa. Tipično
vrednotenje traja $2{,}0$~ms, v najslabšem odstotku primerov (p99) pa
$13{,}5$~ms. Skoraj celotna latenca ($\approx 88\,\%$ povprečja) odpade na
poizvedbo SQL; točkovanje v Pythonu ($0{,}4$~ms povprečno) in serializacija
odgovora ($\approx 0{,}01$~ms) sta zanemarljiva. Poizvedba v povprečju vrne
$287$ primerljivih vozil (mediana $175$, največ $1495$).

\begin{table}[htbp]
  \centering
  \caption{Latenca posameznega vrednotenja po dodajanju indeksa
  ($N = 2000$), v milisekundah.}
  \label{tab:latenca-vrednotenja}
  % vir: benchmarks/report.md (results_algo.json)
  \begin{tabular}{lrrrr}
    \toprule
    \textbf{Komponenta} & \textbf{p50} & \textbf{p95} & \textbf{p99} & \textbf{Povprečje} \\
    \midrule
    Skupaj                & 2{,}0  & 9{,}7 & 13{,}5 & 3{,}1 \\
    Poizvedba SQL         & 1{,}8  & 8{,}4 & 11{,}5 & 2{,}7 \\
    Točkovanje (Python)   & 0{,}2  & 1{,}2 & 1{,}8  & 0{,}4 \\
    Serializacija (JSON)  & 0{,}01 & 0{,}02 & 0{,}03 & 0{,}01 \\
    \bottomrule
  \end{tabular}
\end{table}

\begin{figure}[htbp]
  \centering
  \includegraphics[width=0.6\textwidth]{latency_breakdown.pdf}
  \caption{Razčlenitev latence posameznega vrednotenja po komponentah.}
  \label{fig:latency-breakdown}
  % vir: analysis/out/latency_breakdown.pdf
\end{figure}

Vrednotenje se torej izvede daleč pod nefunkcionalno zahtevo $2$~s
(razdelek~\ref{sec:zahteve}); prostora za rast podatkovne baze je ob
prisotnosti indeksa dovolj.

% =====================================================================
\section{Obnašanje pod obremenitvijo}
\label{sec:obremenitev}

Merimo end-to-end latenco (vrsta (c)) in prepustnost končne točke
\texttt{POST /valuate} pri naraščajoči sočasnosti. Obremenitev ustvarja
skripta \texttt{benchmarks/loadtest.py} z več vzporednimi odjemalci, ki v
fiksnem $20$-sekundnem oknu neprekinjeno pošiljajo zahtevke z naključnimi
realnimi vhodi. Test se samodejno prekine ob deležu napak nad $5\,\%$ ali
$p_{95}$ nad $5$~s. Med testom vzorčimo tudi zasedenost baze
(\texttt{pg\_stat\_activity}).

\begin{table}[htbp]
  \centering
  \caption{Prepustnost in latenca pri naraščajoči sočasnosti.
  Latence so v milisekundah, $L$ je ocena po Littlovem zakonu.}
  \label{tab:obremenitev}
  % vir: benchmarks/results_load.json (loadtest.py --db-sample)
  \begin{tabular}{r r r r r r r}
    \toprule
    \textbf{Sočasnost} & \textbf{Prep.\ [z/s]} & \textbf{p50} & \textbf{p95} & \textbf{p99}
      & \textbf{Napake} & \textbf{$L$} \\
    \midrule
    1  & \res{x} & \res{x} & \res{x} & \res{x} & \res{x} & \res{x} \\
    5  & \res{x} & \res{x} & \res{x} & \res{x} & \res{x} & \res{x} \\
    10 & \res{x} & \res{x} & \res{x} & \res{x} & \res{x} & \res{x} \\
    20 & \res{x} & \res{x} & \res{x} & \res{x} & \res{x} & \res{x} \\
    40 & \res{x} & \res{x} & \res{x} & \res{x} & \res{x} & \res{x} \\
    \bottomrule
  \end{tabular}
\end{table}

\paragraph{Rezultat.}
\res{Prepustnost narašča približno linearno do sočasnosti X, nato se ustali pri
približno X zahtevkih na sekundo. Vse latence ostanejo pod nefunkcionalno
zahtevo 2~s.}

\paragraph{Skladnost meritve (Littlov zakon).}
\res{Stolpec $L$ (prepustnost krat povprečna latenca) ostaja blizu nastavljene
sočasnosti, kar potrjuje notranjo skladnost meritve.}

\paragraph{Dokaz ozkega grla.}
Teoretični strop postavitve je $1000 / S \approx 324$ vrednotenj na sekundo,
kjer je $S = 3{,}1$~ms povprečni čas ene evalvacije. Dejanska meja je nižja,
ker ves promet serializira \emph{ena sama deljena povezava do baze}
(\texttt{app/database.py}) in en sam strežniški proces.
\res{To potrjuje vzorčenje baze med testom: število aktivnih poizvedb kljub
naraščajoči sočasnosti odjemalcev ne preseže X.}

% =====================================================================
\section{Preverjanje funkcionalnih zahtev}
\label{sec:funkcionalne}

Tabela~\ref{tab:sledljivost} povezuje vsako funkcionalno zahtevo iz
razdelka~\ref{sec:zahteve} z načinom preverjanja in izidom.

\begin{table}[htbp]
  \centering
  \caption{Sledljivost funkcionalnih zahtev do preverjanja.}
  \label{tab:sledljivost}
  % vir: ročni scenariji + imena Jest testov v car-stock-market-be
  \begin{tabularx}{\textwidth}{@{}X X c@{}}
    \toprule
    \textbf{Zahteva} & \textbf{Način preverjanja} & \textbf{Izid} \\
    \midrule
    Registracija uporabnika              & \res{test / scenarij} & \res{OK} \\
    Prijava in odjava                    & \res{test / scenarij} & \res{OK} \\
    Objava oglasa (le registrirani)      & \res{test / scenarij} & \res{OK} \\
    Vrednotenje vozila po atributih      & \res{test / scenarij} & \res{OK} \\
    Všečkanje oglasov                    & \res{test / scenarij} & \res{OK} \\
    Ažuriranje lastnih/všečkanih oglasov & \res{test / scenarij} & \res{OK} \\
    Pregled oglasov (anonimni)           & \res{test / scenarij} & \res{OK} \\
    \bottomrule
  \end{tabularx}
\end{table}

\paragraph{Pokritost testov.}
Avtomatizirani testi zalednega sistema (Jest + Supertest,
razdelek~\ref{sec:testiranje-be}) obsegajo \res{X} testov v \res{X} zbirkah in
pokrivajo \res{X}\,\% stavkov ter \res{X}\,\% vej. Testi se izvedejo ob vsakem
zahtevku za združitev na cevovodu neprekinjene integracije
(razdelek~\ref{subsec:ci}).

% =====================================================================
\section{Preverjanje nefunkcionalnih zahtev}
\label{sec:nefunkcionalne}

Tabela~\ref{tab:nfr} za vsako nefunkcionalno zahtevo iz
razdelka~\ref{sec:zahteve} navaja merilo, izmerjeno vrednost in sodbo.

\begin{table}[htbp]
  \centering
  \caption{Preverjanje nefunkcionalnih zahtev.}
  \label{tab:nfr}
  \setlength{\tabcolsep}{4pt}
  \begin{tabularx}{\textwidth}{@{}>{\raggedright\arraybackslash}X
                                 >{\raggedright\arraybackslash}X
                                 >{\raggedright\arraybackslash}X c@{}}
    \toprule
    \textbf{Zahteva} & \textbf{Merilo} & \textbf{Izmerjeno} & \textbf{Sodba} \\
    \midrule
    Odziv vmesnika na dejanje uporabnika & $< 100$~ms & \res{x}~ms  & \res{OK} \\
    Zahtevki HTTP (razen vrednotenja)    & $< 300$~ms & \res{x}~ms ($p_{95}$) & \res{OK} \\
    Vrednotenje vozila                   & $< 2$~s    & $13{,}5$~ms ($p_{99}$, brez omrežja) & OK \\
    Vizualna privlačnost / dostopnost    & Lighthouse & \res{x} / \res{x} & \res{OK} \\
    Prilagodljiva postavitev             & preizkus na več velikostih zaslona & \res{da} & \res{OK} \\
    Varovanje uporabniških podatkov      & bcrypt, JWT z rokom & \res{opis} & \res{OK} \\
    Odzivnost pri zahtevnejših opravilih & asinhrona obdelava & \res{opis} & \res{OK} \\
    \bottomrule
  \end{tabularx}
\end{table}

\paragraph{Meritev odzivnosti vmesnika in endpointov.}
\res{Odzivnost uporabniškega vmesnika smo izmerili z orodjem Lighthouse na treh
ključnih straneh; ocene za performanco so X, za dostopnost pa X. Latenco
glavnih endpointov zalednega sistema smo izmerili z orodjem autocannon:
$p_{95}$ je v vseh primerih pod X~ms, torej pod zahtevo 300~ms.}

\paragraph{Varnost.}
Gesla so pred shranjevanjem zgoščena z algoritmom bcrypt
(razdelek~\ref{sec:avtentikacija}), žetoni JWT imajo rok veljavnosti $7$~dni,
poti vmesnika API pa ščiti sredinski preverjevalnik žetonov. \res{Orodje
\texttt{npm audit} ob zadnjem pregledu ni javilo kritičnih ranljivosti v
odvisnostih.}

% =====================================================================
\section{Uporabniški preizkus}
\label{sec:uporabniki}

Izvedli smo formativni uporabniški preizkus~\cite{nielsen1994} z namenom
odkriti težave uporabnosti, ne pa statistično ovrednotiti izdelka.

\paragraph{Udeleženci in postopek.}
V preizkus smo vključili \res{N} udeležencev iz kroga družine in bližnjih
prijateljev, katerih starost in računalniška pismenost sta se med seboj
razlikovali (tabela~\ref{tab:udelezenci}). Vsak udeleženec je prejel enak
seznam nalog: registracija, objava oglasa, všečkanje oglasa, izbira
slovenskega jezika v nastavitvah in vrednotenje poljubnega vozila. Med
izvajanjem je udeleženec razmišljal na glas, opazovalec pa je beležil, ali je
nalogo dokončal samostojno, in vsako opaženo težavo.

\begin{table}[htbp]
  \centering
  \caption{Udeleženci uporabniškega preizkusa.}
  \label{tab:udelezenci}
  \begin{tabular}{c c c l}
    \toprule
    \textbf{\#} & \textbf{Starost (razpon)} & \textbf{Rač.\ pismenost (1--5)} & \textbf{Naprava} \\
    \midrule
    1 & \res{x} & \res{x} & \res{x} \\
    2 & \res{x} & \res{x} & \res{x} \\
    3 & \res{x} & \res{x} & \res{x} \\
    \res{\dots} & & & \\
    \bottomrule
  \end{tabular}
\end{table}

\paragraph{Rezultati.}
Vsi udeleženci so vse naloge dokončali brez pomoči, kar potrjuje, da so glavne
poti v vmesniku dovolj jasne (RV4). Kljub temu so odkrili nekaj
pomanjkljivosti; tabela~\ref{tab:tezave} jih razvršča po resnosti (lestvica
$0$--$4$ po Nielsenu~\cite{nielsen1994}) in navaja, ali so bile odpravljene.

\begin{table}[htbp]
  \centering
  \caption{Težave, odkrite med uporabniškim preizkusom
  (resnost po lestvici $0$--$4$).}
  \label{tab:tezave}
  \setlength{\tabcolsep}{5pt}
  \begin{tabularx}{\textwidth}{@{}>{\raggedright\arraybackslash}X c c@{}}
    \toprule
    \textbf{Težava} & \textbf{Resnost} & \textbf{Status} \\
    \midrule
    Premajhni naslovi podstrani v glavi glavne strani       & \res{x} & odpravljeno \\
    Manjkajoče polje za ponovni vnos gesla pri registraciji & \res{x} & odpravljeno \\
    Velika števila brez ločil za tisočice                   & \res{x} & odpravljeno \\
    Manjkajoč opis pod slikami oblik vozil v obrazcu        & \res{x} & odpravljeno \\
    Premalo vozil, dostopnih za vrednotenje                 & \res{x} & \res{v načrtu} \\
    \bottomrule
  \end{tabularx}
\end{table}

Omejitve preizkusa (majhen in priložnostni vzorec, brez standardiziranega
vprašalnika, brez merjenja časa po nalogah) obravnavamo v
razdelku~\ref{sec:groznje}.

% =====================================================================
\section{Simulacija uporabe}
\label{sec:simulacija}

Za ponazoritev delovanja rešitve povzemamo tipičen potek uporabe. Uporabnik je
ob vstopu usmerjen na glavno stran (slika~\ref{fig:glavna-stran}), kjer sta
izpostavljeni prijava in registracija ter pregled novih in vseh objavljenih
oglasov. Registriran uporabnik lahko objavi oglas prek vodene forme
(slika~\ref{fig:objava-oglasa}), ovrednoti poljubno vozilo, všečka oglase in si
ogleda svoj profil z všečkanimi vozili (slika~\ref{fig:liked-cars}).

\begin{figure}[H]
  \centering
  \includegraphics[width=0.9\textwidth]{glavna-stran.png}
  \caption{Glavna stran prototipa}
  \label{fig:glavna-stran}
\end{figure}
\begin{figure}[H]
  \centering
  \includegraphics[width=0.9\textwidth]{objava-oglasa.png}
  \caption{Vodena forma za objavo oglasa}
  \label{fig:objava-oglasa}
\end{figure}
\begin{figure}[H]
  \centering
  \includegraphics[width=0.9\textwidth]{liked-cars.png}
  \caption{Stran uporabnika s prikazanimi všečkanimi vozili}
  \label{fig:liked-cars}
\end{figure}

\paragraph{Primerjava ocen z obstoječimi portali.}
Za ilustracijo smo delovanje sistema primerjali s ponudbo na dveh spletnih
portalih na treh vozilih (tabela~\ref{tab:vrednotenje}). Sistematično in
kvantitativno oceno natančnosti podaja razdelek~\ref{sec:natancnost}; ta
primerjava je zgolj kvalitativna ponazoritev.

\begin{table}[htbp]
  \centering
  \caption{Primerjava ocenjenih cen s ponudbo na spletnih portalih (v evrih).}
  \label{tab:vrednotenje}
  \setlength{\tabcolsep}{5pt}
  \begin{tabular}{lccc}
    \toprule
    & \textbf{Audi A3} & \textbf{BMW ser.\ 2} & \textbf{Citroën Berlingo} \\
    \midrule
    slov.\ portal (razpon)     & 4.000--6.000  & 9.000--12.000  & 7.500--10.000 \\
    nem.\ portal (razpon)      & 4.500--9.000  & pribl.\ 11.100 & 8.000--12.000 \\
    \midrule
    Sistem --- dobra cena      & 4.500  & 7.600   & 7.600  \\
    Sistem --- sredinska cena  & 6.350  & 9.300   & 9.941  \\
    Sistem --- slaba cena      & 7.600  & 10.500  & 12.100 \\
    Raznolikost                & $49\,\%$ & $31\,\%$ & $45\,\%$ \\
    \bottomrule
  \end{tabular}
\end{table}

Pri vozilih z nizko raznolikostjo ponudbe (Audi~A3, BMW serije~2) se sredinska
ocena dobro ujema s cenami na obeh portalih. Pri modelu Citroën Berlingo je
raznolikost bistveno višja, zato je razpon med dobro in slabo oceno širši,
sredinska ocena pa manj zanesljiva --- kar je skladno z ugotovitvijo iz
razdelka~\ref{sec:natancnost}, da napaka strmo narašča z raznolikostjo ponudbe.
Pri interpretaciji je treba upoštevati, da sistem ocenjuje pričakovano prodajno
in ne oglaševano ceno.

% =====================================================================
\section{Grožnje veljavnosti}
\label{sec:groznje}

\paragraph{Notranja veljavnost.}
Kot merilo resnice pri merjenju natančnosti uporabljamo \emph{oglaševano} ceno,
ne dejanske prodajne cene, ki ni javno dostopna. Namerni popust, s katerim
sistem cilja pričakovano prodajno ceno, je zato ocenjen in ne umerjen na
dejanskih prodajah. Isti fizični oglas se lahko pojavi pri več virih; prečno
preverjanje izloči zapise z istim izvornim oglasom, ne more pa zaznati
podvojitev med viri.

\paragraph{Zunanja veljavnost.}
Podatki pokrivajo $10$ znamk in pretežno slovenski trg (z manjšim nemškim
vzorcem); znamka z največ oglasi na trgu (okoli $8000$ oglasov) v šifrantu ni
zajeta in je iz analize izpuščena. Uporabniški preizkus je temeljil na majhnem
in priložnostnem vzorcu.

\paragraph{Konstruktna veljavnost.}
Pojem \enquote{tržna vrednost} operacionaliziramo kot uteženo povprečje
oglaševanih cen primerljivih vozil, znižano za pričakovani pogajalski popust;
\enquote{raznolikost ponudbe} pa kot relativno širino interkvartilnega razpona
teh cen. Drugačna operacionalizacija bi lahko dala drugačne rezultate.

\paragraph{Zanesljivost.}
Meritve latence in obremenitve so bile izvedene na deljenem produkcijskem
strežniku, kjer lahko na rezultate vpliva ozadje. Meritve natančnosti so
ponovljive nad fiksnim posnetkom baze.

% =====================================================================
\section{Povzetek ugotovitev}
\label{sec:povzetek-analize}

Tabela~\ref{tab:rv-povzetek} povzema odgovore na raziskovalna vprašanja glede
na vnaprej postavljena merila.

\begin{table}[htbp]
  \centering
  \caption{Povzetek odgovorov na raziskovalna vprašanja.}
  \label{tab:rv-povzetek}
  \begin{tabularx}{\textwidth}{@{}l X c@{}}
    \toprule
    \textbf{RV} & \textbf{Glavni rezultat} & \textbf{Sodba} \\
    \midrule
    RV1 natančnost & MdAPE $13{,}5\,\%$ ($9$--$10\,\%$ za dobro pokrita vozila),
                     bistveno bolje od naivne ocene ($37{,}8\,\%$), primerljivo z
                     drugimi metodami & izpolnjeno \\
    RV2 odvisnost  & napaka pada z velikostjo vzorca (ustali se pri $n\approx20$),
                     narašča s starostjo, kilometrino in raznolikostjo ponudbe
                     ($7{,}5\,\%$ do $29{,}7\,\%$) & pojasnjeno \\
    RV3 hitrost    & vrednotenje $2$~ms (p50), $13{,}5$~ms (p99); indeks odpravi
                     odvisnost od velikosti baze; \res{meja X~z/s} & \res{izpolnjeno} \\
    RV4 uporabnost & vse naloge dokončane brez pomoči; \res{X} težav odpravljenih & izpolnjeno \\
    \bottomrule
  \end{tabularx}
\end{table}

\paragraph{Sodba o zastavljenih izhodiščih.}
Rezultati \emph{potrjujejo} izhodišča, ki smo si jih zastavili na začetku. Sistem na podlagi štirih
atributov vozila vrne oceno vrednosti z mediano odstotne napake $13{,}5\,\%$
(brez namernega popusta), pri dobro pokritih vozilih pa $9$--$10\,\%$, kar je
znotraj zastavljenega praga $15\,\%$ in bistveno bolje od naivne metode.
Predlagani pristop se med primerljivo natančnimi metodami ne uvršča na prvo
mesto --- preprost $k$-najbližjih sosedov doseže $11{,}2\,\%$ --- je pa edini,
ki poleg točkovne ocene vrne tudi razložljiv cenovni razpon in indeks
zaupanja, ta pa se izkaže za dober pokazatelj dejanske zanesljivosti
posamezne ocene. Glavne omejitve so odvisnost od enega prevladujočega
podatkovnega vira, nižja natančnost pri starejših vozilih in modelih z
razpršeno ponudbo ter merilo resnice, ki temelji na oglaševanih in ne
prodajnih cenah. Konkretne smeri za izboljšave --- vključno s prehodom na
metode strojnega učenja, ki jih ta analiza dodatno utemeljuje --- navajamo v
poglavju~\ref{ch:sklep}.

%
%  Številke so vpisane iz izhodov skript v CarStockMarketEvaluation/analysis/.
%  Preostale oznake \res{...} (rdeče) je treba še izpolniti:
%    - obremenitveni test (benchmarks/loadtest.py)
%    - sledljivost testov, Jest coverage, Lighthouse
%    - tabeli udeležencev in najdenih težav
%  Slike (PDF) iz analysis/out/ prekopiraj v mapo slike/.
% =====================================================================

\newcommand{\res}[1]{\textcolor{red}{\textbf{[#1]}}}

\chapter{Analiza in vrednotenje rešitve}
\label{ch:analiza}

V tem poglavju sistematično pregleadamo in analiziramo razvito rešitev. Najprej iz zahtev in
izhodišč, ki smo si jih zastavili na začetku, izpeljemo raziskovalna vprašanja in za vsako določimo merilo
uspešnosti (razdelek~\ref{sec:rq}). Nato opišemo merilno okolje in postopke
(razdelek~\ref{sec:metodologija}) ter natančno definiramo vse uporabljene mere
(razdelek~\ref{sec:mere}). Sledijo rezultati: značilnosti zbranih
podatkov~(\ref{sec:podatki}), natančnost ocene
vrednosti~(\ref{sec:natancnost}), občutljivost algoritma na
parametre~(\ref{sec:obcutljivost}), latenca posameznega
vrednotenja~(\ref{sec:latenca}) in obnašanje pod
obremenitvijo~(\ref{sec:obremenitev}). Zatem preverimo funkcionalne in
nefunkcionalne zahteve (\ref{sec:funkcionalne}, \ref{sec:nefunkcionalne}) ter
predstavimo uporabniški preizkus~(\ref{sec:uporabniki}) in simulacijo
uporabe~(\ref{sec:simulacija}). Poglavje sklenemo z obravnavo groženj
veljavnosti~(\ref{sec:groznje}) in povzetkom ugotovitev glede na
zastavljena izhodišča~(\ref{sec:povzetek-analize}).

% =====================================================================
\section{Raziskovalna vprašanja in merila uspešnosti}
\label{sec:rq}

Na začetku diplomskega dela (poglavje~\ref{ch:uvod}) smo si zastavili, da je
mogoče uporabniku na podlagi zbranih tržnih podatkov in primerjave s podobnimi
vozili ponuditi \emph{uporabno} in \emph{dovolj natančno} oceno tržne vrednosti
vozila in da vključitev takega sistema izboljša uporabniško izkušnjo. Iz tega
izhodišča in iz nefunkcionalnih zahtev (razdelek~\ref{sec:zahteve}) izpeljemo
štiri raziskovalna vprašanja (RV1--RV4) in za vsako določimo merilo uspešnosti.

\emph{Mediana absolutne odstotne napake} (MdAPE) je osrednja mera natančnosti v
tem poglavju; nanjo se sklicujejo merila uspešnosti in vsi nadaljnji razdelki.
Izračunamo jo tako, da za vsak testni primer določimo absolutno odstotno napako,
torej za koliko odstotkov se ocenjena cena razlikuje od dejanske, nato pa čez
vse primere vzamemo mediano teh vrednosti. MdAPE $=15\,\%$ tako pomeni, da je pri
polovici vozil ocena znotraj $15\,\%$ od dejanske cene, pri drugi polovici pa je
odstopanje večje; manjša kot je vrednost, natančnejše so ocene. Mediano
uporabljamo namesto povprečja, ker je manj občutljiva na posamezne skrajne
napake. Natančne definicije MdAPE in preostalih mer so v
razdelku~\ref{sec:mere}.

Raziskovalna vprašanja so:

\begin{itemize}
  \item \textbf{RV1 --- natančnost ocene.} Kako natančno sistem oceni tržno
        vrednost vozila v primerjavi z dejanskimi oglasi na trgu?
        \emph{Merilo:} MdAPE je pri vozilih z vsaj $10$ primerljivimi
        vozili manjša od $15\,\%$, ocena pa je bistveno natančnejša od
        preproste (naivne) referenčne metode.
  \item \textbf{RV2 --- odvisnost natančnosti od podatkov.} Od česa je
        natančnost ocene najbolj odvisna (velikost vzorca primerljivih vozil,
        razpršenost ponudbe, segment vozila)?
  \item \textbf{RV3 --- hitrost in razširljivost.} Ali sistem zadošča
        nefunkcionalnim zahtevam glede odzivnosti tudi pod sočasno
        obremenitvijo, in kje je zgornja meja trenutne postavitve?
  \item \textbf{RV4 --- uporabnost.} Ali ciljni uporabniki znajo brez navodil
        opraviti glavne naloge (registracija, objava oglasa, vrednotenje)?
\end{itemize}
% =====================================================================
\section{Metodologija in merilno okolje}
\label{sec:metodologija}

\paragraph{Merilno okolje.}
Meritve hitrosti (razdelka~\ref{sec:latenca} in~\ref{sec:obremenitev}) so
izvedene v \emph{produkcijskem} okolju, torej na navideznem zasebnem strežniku
z dvema procesorskima jedroma in $4$~GB pomnilnika ter s tehnologijami in
njihovimi različicami, opisanimi v poglavju~\ref{ch:razvoj}
(razdelek~\ref{sec:namestitev}). Analitične in merilne skripte tečejo nad isto
podatkovno bazo kot je na produkciji.

\paragraph{Posnetek podatkov.}
Meritve v tem poglavju temeljijo na posnetku podatkovne baze iz septembra
$2026$, ki obsega približno $86.000$ normaliziranih oglasov
(razdelek~\ref{sec:podatki}). Ker se posnetek med analizo ne spreminja, so
rezultati ponovljivi.

\paragraph{Orodja.}
Za meritve in analizo v tem poglavju smo napisali namenske skripte, ločene od
produkcijske kode; ena od njih je dokumentirana kopija ocenjevalnega
algoritma, parametrizirana za eksperimente. Statistično obdelavo in izris
grafov smo izvedli s knjižnicama \texttt{NumPy}~\cite{harris2020numpy} in
\texttt{Matplotlib}~\cite{hunter2007matplotlib}.

\paragraph{Natančnost.}
Natančnost ocene (razdelek~\ref{sec:natancnost}) merimo z naključno izbranim
normaliziranim oglasom: njegove atribute (znamka, model, letnik, kilometrina)
uporabimo kot vhod v algoritem, sam oglas pa izločimo iz primerjalne množice. Nato izračunamo oceno in jo
primerjamo z \emph{dejansko oglaševano} ceno tega oglasa. Postopek ponovimo za
$3000$ naključnih oglasov.

% =====================================================================
\section{Uporabljene mere}
\label{sec:mere}

\paragraph{Latenca}
 je čas, ki preteče od trenutka tik pred oddajo zahtevka do
trenutka, ko odjemalec prejme celoten odgovor. Glede na to, kolikšen del te
poti zajamemo, ločimo tri vrste:
\begin{itemize}
  \item[(a)] \emph{latenca obdelave na strežniku} --- čas izvajanja funkcije
        za vrednotenje, od začetka do konca, brez omrežja;
  \item[(b)] \emph{latenca poizvedbe SQL} --- čas izvajanja poizvedbe v podatkovni bazi,
   od začetka do konca, brez omrežja;
  \item[(c)] \emph{latenca od konca do konca} (angl.\ \emph{end-to-end}) ---
        (a) skupaj z omrežnim obhodom med odjemalcem in strežnikom ter
        obdelavo zahtevka HTTP; to je čas, ki ga zazna uporabnik.
\end{itemize}
V razdelku~\ref{sec:latenca} merimo vrsti (a) in (b), v
razdelku~\ref{sec:obremenitev} pa vrsto (c).

\paragraph{Percentili.}
\textbf{Percentil} $p_k$ je vrednost, pod katero pade $k\,\%$ meritev. Poročamo
$p_{50}$ (mediana), $p_{95}$ in $p_{99}$.

\paragraph{Prepustnost, sočasnost, nasičenje.}
Obnašanje sistema pod obremenitvijo opišemo s tremi količinami:
\begin{itemize}
  \item \textbf{prepustnost} --- število uspešno obdelanih zahtevkov na
        sekundo;
  \item \textbf{sočasnost} --- število odjemalcev, ki hkrati držijo odprt
        zahtevek;
  \item \textbf{nasičenje} (\emph{koleno}) --- točka, v kateri dodatna
        sočasnost ne poveča več prepustnosti, latenca pa začne naraščati
        približno linearno.
\end{itemize}
Notranjo skladnost meritev preverimo z \textbf{Littlovim zakonom}~\cite{little1961},
\begin{equation}
  \label{eq:little}
  L = \lambda \cdot W,
\end{equation}
kjer je $L$ povprečno število zahtevkov v sistemu, $\lambda$ prepustnost in
$W$ povprečna latenca.

\paragraph{Napaka ocene.}
Za oglas $i$ z dejansko ceno $p_i$ in ocenjeno ceno $\hat p_i$ definiramo
absolutno odstotno napako
\begin{equation}
  \mathrm{APE}_i = \frac{\lvert \hat p_i - p_i \rvert}{p_i}.
\end{equation}
\textbf{MAPE} je povprečje $\mathrm{APE}_i$, \textbf{MdAPE} pa njihova mediana;
MdAPE je robustnejša na osamelce in je naša glavna mera~\cite{hyndman2006}.
Poročamo tudi \textbf{MAE} (povprečna absolutna napaka) in \textbf{RMSE}
(koren povprečne kvadratne napake) v evrih ter \emph{predznačeno} odstotno
napako, ki pokaže sistematičen odmik ocene navzgor ali navzdol.

\paragraph{Pokritost.}
\textbf{Pokritost} je delež poizvedb, za katere sistem vrne oceno, torej za
katere najde vsaj eno (oziroma vsaj $5$ oz.\ $10$) primerljivo vozilo.

\paragraph{Raznolikost ponudbe.}
\textbf{Raznolikost} (v izvorni kodi \texttt{ratio}) je relativna širina
cenovnega razpona primerljivih vozil,
\begin{equation}
  \label{eq:raznolikost}
  \mathrm{raznolikost} = \frac{P_{75} - P_{25}}{\hat{p}},
\end{equation}
kjer sta $P_{25}$ in $P_{75}$ (v kodi \texttt{goodPrice} in \texttt{badPrice})
uteženi $25.$ oziroma $75.$ percentil cen primerljivih vozil, $\hat{p}$
(v kodi \texttt{estimatedPrice}) pa ocenjena cena vozila. Vrednost izražamo v
odstotkih; enoten zapis uporabljamo v vseh tabelah tega poglavja.

% =====================================================================
\section{Značilnosti zbranih podatkov}
\label{sec:podatki}
\paragraph{Viri in obseg.}
Velika večina oglasov je bila zajeta z največjega slovenskega spletnega portala
za prodajo rabljenih vozil (razdelek~\ref{sec:zajem-podatkov}); manjši del
predstavlja vzorec z večjega nemškega oglasnega portala, ki smo ga uporabili za
preverjanje delovanja cevovoda na drugačni strukturi vira. Zajem je zajel
javno dostopne oglase za nekomercialno raziskovalno rabo, z omejeno hitrostjo
zahtevkov in brez shranjevanja kontaktnih ali osebnih podatkov prodajalcev.
Po normalizaciji baza obsega $85.846$ oglasov ($83.493$ s slovenskega in
$2.353$ z nemškega portala) za $10$ znamk vozil; razčlenitev po znamki je v
tabeli~\ref{tab:podatki-znamke}.

\begin{table}[htbp]
  \centering
  \caption{Število normaliziranih oglasov po znamki.}
  \label{tab:podatki-znamke}
  % vir: analysis/out/dataset_per_brand.csv
  \begin{tabular}{lr lr}
    \toprule
    \textbf{Znamka} & \textbf{Oglasi} & \textbf{Znamka} & \textbf{Oglasi} \\
    \midrule
    Volkswagen    & 21.191 & Peugeot     & 7.513 \\
    Audi          & 11.842 & Citroën     & 6.729 \\
    BMW           & 10.724 & Ford        & 5.783 \\
    Renault       & 9.615  & Toyota      & 2.648 \\
    Mercedes-Benz & 9.355  & Alfa Romeo  & 446 \\
    \bottomrule
  \end{tabular}
\end{table}

\paragraph{Porazdelitve ključnih atributov.}
Slika~\ref{fig:porazdelitve} prikazuje porazdelitve oglaševanih cen, prevožene
kilometrine in letnika. Porazdelitev cen je izrazito desno asimetrična (dolg
rep dragih vozil), zato jo prikazujemo na logaritemski osi. Mediana cene je
$14.890$~EUR (povprečje $18.137$~EUR), mediana kilometrine $159.950$~km,
mediana letnika pa $2017$; podrobnosti so v tabeli~\ref{tab:podatki-porazdelitve}.
Zaradi berljivosti so na grafih odrezane skrajne vrednosti (nekaj oglasov ima
očitno napačno ceno, kilometrino ali letnik); njihovo število je označeno v
kotu vsakega grafa, celoten razpon pa je razviden iz stolpca \emph{maks.}\ v
tabeli~\ref{tab:podatki-porazdelitve} in iz obravnave osamelcev v
tabeli~\ref{tab:osamelci}.

\begin{table}[htbp]
  \centering
  \caption{Povzetek porazdelitev ključnih numeričnih atributov.}
  \label{tab:podatki-porazdelitve}
  % vir: analysis/out/dataset_distribution_summary.json
  \setlength{\tabcolsep}{5pt}
  \begin{tabular}{lrrrrr}
    \toprule
    & \textbf{$p_{5}$} & \textbf{mediana} & \textbf{povprečje} & \textbf{$p_{95}$} & \textbf{maks.} \\
    \midrule
    cena [EUR]        & 1.550  & 14.890  & 18.137  & 49.014  & 374.000 \\
    kilometrina [km]  & 26.000 & 159.950 & 169.464 & 334.000 & 1.974.001 \\
    letnik            & 2005   & 2017    & 2016    & 2024    & 2029 \\
    moč [kW]          & 56     & 110     & 117     & 221     & 920 \\
    \bottomrule
  \end{tabular}
\end{table}

\begin{figure}[htbp]
  \centering
  \begin{subfigure}{0.48\textwidth}
    \includegraphics[width=\textwidth]{dataset_price_hist.pdf}
    \caption{oglaševana cena (log)}
  \end{subfigure}\hfill
  \begin{subfigure}{0.48\textwidth}
    \includegraphics[width=\textwidth]{dataset_km_hist.pdf}
    \caption{prevoženi kilometri}
  \end{subfigure}
  \\[1ex]
  \begin{subfigure}{0.48\textwidth}
    \includegraphics[width=\textwidth]{dataset_year_hist.pdf}
    \caption{letnik prve registracije}
  \end{subfigure}\hfill
  \begin{subfigure}{0.48\textwidth}
    \includegraphics[width=\textwidth]{dataset_brand_bar.pdf}
    \caption{oglasi po znamki}
  \end{subfigure}
  \caption{Porazdelitve ključnih atributov v normalizirani bazi.}
  \label{fig:porazdelitve}
  % vir: analysis/out/dataset_*_hist.pdf, dataset_brand_bar.pdf
\end{figure}

\paragraph{Pokritost po segmentih.}
Za pošteno oceno mora imeti algoritem dovolj primerljivih vozil v ožjem
segmentu (ista znamka in model, podoben letnik in kilometrina).
Slika~\ref{fig:pokritost} prikazuje toplotno karto za znamki Renault in
Volkswagen: vrstice so modeli, stolpci letniki, barva pa
število oglasov v tej kombinaciji. Vidno je, da so pogosti modeli srednjih
letnikov (npr.\ Volkswagen Golf ali Renault Clio) gosto pokriti, redki modeli
in skrajni letniki pa le redko, kar neposredno pojasni razlike v zanesljivosti
ocen (razdelek~\ref{sec:natancnost}).

\begin{figure}[htbp]
  \centering
  \includegraphics[width=\textwidth]{dataset_coverage_heatmap.pdf}
  \caption{Pokritost: število oglasov po modelu in letniku za znamki
  Renault in Volkswagen (po ena znamka na vrstico).}
  \label{fig:pokritost}
  % vir: analysis/out/dataset_coverage_heatmap.pdf
\end{figure}

\paragraph{Osamelci in kakovost podatkov.}
Postopek normalizacije trenutno ne izvaja filtriranja osamelcev, zato baza
vsebuje tudi vozila z očitno napačnimi vrednostmi.
Tabela~\ref{tab:osamelci} prikazuje njihov delež. Kar $10{,}7\,\%$ oglasov ima
vsaj eno vrednost zunaj smiselnih mej, največ pri kilometrini ($9{,}9\,\%$;
najvišja zabeležena vrednost je skoraj dva milijona kilometrov) in letniku
($8{,}1\,\%$; nekaj oglasov ima celo letnik v prihodnosti). Preprost filter po
interkvartilnem razponu (IQR) na ceni znotraj skupine
$(\text{znamka},\text{model})$ bi odstranil dodatnih $2{,}2\,\%$ oglasov.
Vpliv doslednega filtriranja osamelcev na natančnost obravnavamo v
razdelku~\ref{sec:obcutljivost}, uvedbo takega filtra pa v
poglavju~\ref{ch:sklep}.

\begin{table}[H]
  \centering
  \caption{Osamelci v normalizirani bazi ($85.846$ oglasov).}
  \label{tab:osamelci}
  % vir: analysis/out/dataset_outliers.csv
  \begin{tabularx}{\textwidth}{@{}X r r@{}}
    \toprule
    \textbf{Merilo} & \textbf{Št.\ vrstic} & \textbf{Delež} \\
    \midrule
    cena zunaj $[500,\ 200.000]$~EUR         & 427   & $0{,}5\,\%$ \\
    kilometrina zunaj $[1000,\ 500.000]$~km  & 8.495 & $9{,}9\,\%$ \\
    letnik zunaj $[1990,\ 2027]$             & 6.979 & $8{,}1\,\%$ \\
    vsaj ena vrednost zunaj mej              & 9.169 & $10{,}7\,\%$ \\
    dodatno po filtru IQR na ceni            & 1.878 & $2{,}2\,\%$ \\
    \bottomrule
  \end{tabularx}
\end{table}

% =====================================================================
\section{Natančnost ocene vrednosti}
\label{sec:natancnost}

To je osrednja meritev diplomskega dela: neposredno preverja zastavljeno
izhodišče o \emph{dovolj natančni} oceni (RV1). Protokol je opisan v
razdelku~\ref{sec:metodologija}, rezultate pa izdela skripta
\texttt{analysis/accuracy\_cv.py} nad vzorcem $3000$ oglasov (od tega jih je
$2986$ dobilo oceno).

\paragraph{Merilo resnice in namerni odmik ocene.}
Kot merilo resnice uporabljamo \emph{oglaševano} ceno, saj dejanske prodajne
cene niso javno dostopne. Sistem pa namenoma ne ocenjuje oglaševane cene,
temveč \emph{pričakovano prodajno ceno}: ker se rabljena vozila v pogajanjih
praviloma prodajo pod oglaševano ceno, oceno zniža s stopenjskim popustom
($9$--$15\,\%$ glede na cenovni razred). Zato je pri primerjavi z oglaševano
ceno pričakovan rahel \emph{negativni} odmik ocene, kar je zaželeno in ne
pomeni slabše ocene. To lastnost upoštevamo pri branju vseh rezultatov v
nadaljevanju: napako merimo dvakrat --- proti oglaševani ceni z vključenim
popustom (kar dejansko vrne sistem) in brez popusta (čista primerjalna
ocena).

\paragraph{Skupni rezultat.}
Brez popusta znaša MdAPE predlaganega pristopa $13{,}5\,\%$, MAPE $23{,}8\,\%$,
MAE $2798$~EUR in RMSE $5151$~EUR; z vključenim namernim popustom je MdAPE
proti oglaševani ceni $14{,}5\,\%$ (mediana predznačene napake $-6{,}8\,\%$).
Pokritost je $99{,}5\,\%$ pri $n \ge 1$ in $95{,}7\,\%$ pri $n \ge 10$.
Slika~\ref{fig:pred-vs-actual} prikazuje napovedano ceno proti dejanski
($R^2 = 0{,}87$), slika~\ref{fig:ape-hist} pa porazdelitev odstotne napake.

\begin{figure}[htbp]
  \centering
  \begin{subfigure}{0.48\textwidth}
    \includegraphics[width=\textwidth]{accuracy_pred_vs_actual.pdf}
    \caption{napovedano proti dejanskemu}
    \label{fig:pred-vs-actual}
  \end{subfigure}\hfill
  \begin{subfigure}{0.48\textwidth}
    \includegraphics[width=\textwidth]{accuracy_ape_hist.pdf}
    \caption{porazdelitev napake}
    \label{fig:ape-hist}
  \end{subfigure}
  \caption{Natančnost ocene vrednosti (leave-one-out, brez popusta).}
  % vir: analysis/out/accuracy_pred_vs_actual.pdf, accuracy_ape_hist.pdf
\end{figure}

\paragraph{Primerjava z osnovnimi modeli.}
Da presodimo, ali je predlagani pristop boljši od trivialnega, ga primerjamo s
štirimi referenčnimi modeli na istem protokolu in isti množici kandidatov
(tabela~\ref{tab:baselines}, slika~\ref{fig:baselines}): B0~globalna mediana
cene za dani model, B1~navadno povprečje cen v filtrirnem oknu, B1m~mediana cen
v oknu, B2~linearna regresija $\text{cena} \sim \text{letnik} + \text{km}$,
B3~metoda $k$~najbližjih sosedov ($k = 10$) po standardiziranih letniku in
kilometrini, ter B4~naš pristop (uteženo povprečje s podobnostnim jedrom, brez
popusta).

\begin{table}[htbp]
  \centering
  \caption{Primerjava natančnosti z osnovnimi modeli (leave-one-out).
  MdAPE in MAPE sta v odstotkih, MAE in RMSE v evrih.}
  \label{tab:baselines}
  % vir: analysis/out/accuracy_overall.csv
  \setlength{\tabcolsep}{5pt}
  \begin{tabular}{l r r r r}
    \toprule
    \textbf{Model} & \textbf{MdAPE} & \textbf{MAPE} & \textbf{MAE} & \textbf{RMSE} \\
    \midrule
    B0 globalna mediana       & 37{,}8 & 95{,}7 & 7395 & 10883 \\
    B1 povprečje v oknu       & 14{,}2 & 24{,}9 & 2930 & 5333 \\
    B1m mediana v oknu        & 13{,}1 & 22{,}5 & 2877 & 5561 \\
    B2 linearna regresija     & 12{,}9 & 22{,}4 & 2677 & 5032 \\
    B3 $k$-najbližjih sosedov & \textbf{11{,}2} & 21{,}5 & 2497 & 4883 \\
    B4 naš pristop            & 13{,}5 & 23{,}8 & 2798 & 5151 \\
    \bottomrule
  \end{tabular}
\end{table}

\begin{figure}[htbp]
  \centering
  \includegraphics[width=0.7\textwidth]{accuracy_baselines.pdf}
  \caption{MdAPE osnovnih modelov in predlaganega pristopa.}
  \label{fig:baselines}
  % vir: analysis/out/accuracy_baselines.pdf
\end{figure}

Predlagani pristop je \emph{bistveno} boljši od naivne globalne mediane
($13{,}5\,\%$ proti $37{,}8\,\%$) in primerljiv s preostalimi metodami, ki se
gibljejo med $11$ in $14\,\%$. Ni pa najboljši: metoda $k$~najbližjih sosedov
doseže $11{,}2\,\%$, tudi navadna mediana istega okna ($13{,}1\,\%$) je
malenkost natančnejša. To je pričakovano za zasnovo, katere cilj je
\emph{razložljivost}: za razliko od naštetih metod predlagani pristop poleg
točkovne ocene vrne tudi cenovni razpon in indeks zaupanja, uteži pa so
neposredno interpretabilne. Rezultat hkrati potrjuje smiselnost prehoda na
metode strojnega učenja, ki jo predlagamo v poglavju~\ref{ch:sklep}: če že
preprost $k$-NN preseže sedanji pristop, je od zmogljivejših modelov mogoče
pričakovati še boljše rezultate.

\paragraph{Odvisnost od velikosti vzorca in segmenta (RV2).}
Slika~\ref{fig:mdape-n} prikazuje MdAPE glede na število primerljivih vozil
$n$: napaka je pri $n < 10$ okoli $20\,\%$, pri $n \ge 20$ pa se ustali pri
približno $13\,\%$. To empirično podpira izbiro konstante $20$ v indeksu
zaupanja (enačba~\ref{eq:confidence}) --- pri toliko primerljivih vozilih je
ocena že blizu svoje najboljše natančnosti. Slika~\ref{fig:mdape-brand} prikazuje MdAPE po znamki: najbolje se
odrežejo pogosti, gosto pokriti modeli (Audi $9{,}8\,\%$, Toyota $10{,}8\,\%$,
Volkswagen $12{,}1\,\%$), najslabše pa Mercedes-Benz ($24{,}9\,\%$), pri
katerem se pri normalizaciji več različnih modelov zlije v isto kanonično
ime, kar poslabša izbor primerljivih vozil.
Tabela~\ref{tab:natancnost-segmenti} kaže še razčlenitev po starosti,
kilometrini in raznolikosti ponudbe.

\begin{figure}[htbp]
  \centering
  \begin{subfigure}{0.48\textwidth}
    \includegraphics[width=\textwidth]{accuracy_mdape_vs_n.pdf}
    \caption{glede na velikost vzorca $n$}
    \label{fig:mdape-n}
  \end{subfigure}\hfill
  \begin{subfigure}{0.48\textwidth}
    \includegraphics[width=\textwidth]{accuracy_mdape_by_brand.pdf}
    \caption{po znamki}
    \label{fig:mdape-brand}
  \end{subfigure}
  \caption{Odvisnost natančnosti ocene od podatkov.}
  % vir: analysis/out/accuracy_mdape_vs_n.pdf, accuracy_mdape_by_brand.pdf
\end{figure}

\begin{table}[htbp]
  \centering
  \caption{Natančnost predlaganega pristopa po segmentih (MdAPE).}
  \label{tab:natancnost-segmenti}
  % vir: analysis/out/accuracy_by_age.csv, accuracy_by_km.csv, accuracy_by_ratio.csv
  \begin{tabularx}{\textwidth}{@{}l X r r@{}}
    \toprule
    \textbf{Razčlenitev} & \textbf{Koš} & \textbf{Primeri} & \textbf{MdAPE [\%]} \\
    \midrule
    starost [leta] & 0--3   & 367 & 9{,}0 \\
                   & 4--7   & 921 & 10{,}6 \\
                   & 8--12  & 932 & 14{,}5 \\
                   & $13+$  & 766 & 22{,}9 \\
    \midrule
    kilometrina    & $<50.000$        & 280  & 9{,}6 \\
                   & $50.000$--$100.000$  & 449  & 11{,}9 \\
                   & $100.000$--$150.000$ & 634  & 12{,}4 \\
                   & $150.000$--$200.000$ & 629  & 13{,}7 \\
                   & $>200.000$       & 1021 & 16{,}7 \\
    \midrule
    raznolikost    & $0$--$0{,}10$    & 88   & 7{,}5 \\
                   & $0{,}10$--$0{,}25$ & 1002 & 8{,}4 \\
                   & $0{,}25$--$0{,}50$ & 1613 & 15{,}5 \\
                   & $>0{,}50$        & 447  & 29{,}7 \\
    \bottomrule
  \end{tabularx}
\end{table}

Vzorec pokaže jasne zakonitosti: napaka narašča s starostjo vozila ($9\,\%$ pri
najnovejših, $23\,\%$ pri vozilih, starejših od $13$~let), s prevoženo
kilometrino in --- najmočneje --- z \emph{raznolikostjo ponudbe}. Pri modelih z
enotno ponudbo (raznolikost pod $0{,}1$) je MdAPE le $7{,}5\,\%$, pri razpršeni
ponudbi (nad $0{,}5$) pa $29{,}7\,\%$. To potrjuje, da je raznolikost, ki jo
sistem uporabniku tudi prikaže, dober vgrajen pokazatelj zanesljivosti
posamezne ocene.

\paragraph{Vpliv popusta.}
Slika~\ref{fig:popust} prikazuje predznačeno napako z in brez popusta. Brez
popusta je ocena rahlo optimistična (mediana $+3{,}8\,\%$ nad oglaševano ceno),
z namernim popustom pa se premakne na $-6{,}8\,\%$ pod oglaševano ceno. Ta
premik je skladen z zasnovo sistema, ki cilja pričakovano prodajno in ne
oglaševane cene; dokončno umerjanje popusta bi bilo mogoče šele s podatki o
dejanskih prodajah (razdelek~\ref{sec:groznje}).

\begin{figure}[htbp]
  \centering
  \includegraphics[width=0.7\textwidth]{accuracy_discount_bias.pdf}
  \caption{Predznačena napaka proti oglaševani ceni, z in brez namernega popusta.}
  \label{fig:popust}
  % vir: analysis/out/accuracy_discount_bias.pdf
\end{figure}

\paragraph{Časovni razrez.}
Pri napovedovanju najnovejših $10\,\%$ oglasov zgolj iz starejših $90\,\%$
znaša MdAPE predlaganega pristopa $13{,}0\,\%$ (tabela~\ref{tab:temporal}), kar
je celo malenkost bolje od rezultata leave-one-out in potrjuje, da ocena ne
izhaja iz uhajanja informacije.

\begin{table}[htbp]
  \centering
  \caption{Natančnost pri časovnem razrezu (napoved najnovejših $10\,\%$).
  MdAPE in MAPE sta v odstotkih, MAE in RMSE v evrih.}
  \label{tab:temporal}
  % vir: analysis/out/accuracy_temporal.csv
  \setlength{\tabcolsep}{5pt}
  \begin{tabular}{l r r r r}
    \toprule
    \textbf{Model} & \textbf{MdAPE} & \textbf{MAPE} & \textbf{MAE} & \textbf{RMSE} \\
    \midrule
    B0 globalna mediana       & 35{,}9 & 74{,}5 & 7596 & 11222 \\
    B1m mediana v oknu        & 12{,}8 & 19{,}7 & 3088 & 5805 \\
    B3 $k$-najbližjih sosedov & 10{,}5 & 18{,}2 & 2667 & 4890 \\
    B4 naš pristop            & 13{,}0 & 20{,}6 & 2994 & 5347 \\
    \bottomrule
  \end{tabular}
\end{table}

% =====================================================================
\section{Občutljivost algoritma na parametre}
\label{sec:obcutljivost}

Ocenjevalni algoritem vsebuje več hevristično določenih konstant: dovoljeno
odstopanje kilometrine ($\pm 40.000$~km) in letnika ($\pm 3$~leta), uteži
delnih ocen ($0{,}6$ in $0{,}4$), oblika padajoče funkcije in stopenjski
popust. V tem razdelku vsako spreminjamo ob nespremenjenih ostalih in merimo
MdAPE ter pokritost (vzorec $2000$ oglasov, leave-one-out). Rezultate izdela
skripta \texttt{analysis/sensitivity.py}.

\paragraph{Okno primerljivosti.}
Tabela~\ref{tab:sens-okno} in slika~\ref{fig:sens-km} kažeta, da je algoritem
na širino \emph{kilometrinskega} okna skoraj neobčutljiv: MdAPE se med
$\pm 10.000$~km in $\pm 80.000$~km giblje le med $13{,}2$ in $13{,}6\,\%$.
Nasprotno je odvisen od \emph{letnikovega} okna: pri $\pm 1$~letu je MdAPE
$12{,}6\,\%$, pri $\pm 5$~letih pa $15{,}2\,\%$. Ožje okno da bolj primerljiva
vozila, a manjši vzorec (nižja pokritost); gre za znano razmerje med
pristranskostjo in varianco~\cite{bias-variance}. Toplotna karta
(slika~\ref{fig:sens-heatmap}) potrjuje, da natančnost določa skoraj izključno
letnikovo okno. Sedanja izbira $\pm 3$~leta je razumen kompromis med
natančnostjo in pokritostjo; morebitna zožitev na $\pm 2$~leti bi MdAPE
znižala za približno $0{,}3$~odstotne točke ob še vedno visoki pokritosti.

\begin{table}[htbp]
  \centering
  \caption{Občutljivost na širino okna primerljivosti. Pokritost je delež
  poizvedb z vsaj $5$ primerljivimi vozili.}
  \label{tab:sens-okno}
  % vir: analysis/out/sensitivity_km.csv, sensitivity_year.csv
  \setlength{\tabcolsep}{5pt}
  \begin{tabular}{l r r r}
    \toprule
    \textbf{Parameter} & \textbf{MdAPE [\%]} & \textbf{Pokritost [\%]} & \textbf{Mediana $n$} \\
    \midrule
    kilometrina $\pm 10.000$~km & 13{,}2 & 92{,}3 & 53 \\
    kilometrina $\pm 40.000$~km & 13{,}4 & 98{,}1 & 196 \\
    kilometrina $\pm 80.000$~km & 13{,}6 & 99{,}0 & 362 \\
    \midrule
    letnik $\pm 1$~leto  & 12{,}6 & 95{,}6 & 97 \\
    letnik $\pm 3$~leta  & 13{,}4 & 98{,}1 & 196 \\
    letnik $\pm 5$~let   & 15{,}2 & 98{,}5 & 262 \\
    \bottomrule
  \end{tabular}
\end{table}

\begin{figure}[htbp]
  \centering
  \begin{subfigure}{0.48\textwidth}
    \includegraphics[width=\textwidth]{sensitivity_km.pdf}
    \caption{okno kilometrine}
  \end{subfigure}\hfill
  \begin{subfigure}{0.48\textwidth}
    \includegraphics[width=\textwidth]{sensitivity_year.pdf}
    \caption{okno letnika}
  \end{subfigure}
  \caption{Občutljivost na širino okna primerljivosti (MdAPE in pokritost).}
  \label{fig:sens-km}
  % vir: analysis/out/sensitivity_km.pdf, sensitivity_year.pdf
\end{figure}

\begin{figure}[htbp]
  \centering
  \includegraphics[width=0.7\textwidth]{sensitivity_km_year_heatmap.pdf}
  \caption{MdAPE [\%] po mreži dovoljenih odstopanj kilometrine in letnika.}
  \label{fig:sens-heatmap}
  % vir: analysis/out/sensitivity_km_year_heatmap.pdf
\end{figure}

\paragraph{Uteži, oblika jedra in popust.}
Tabela~\ref{tab:sens-ostalo} povzema preostale parametre. Razmerje uteži med
kilometrino in letnikom na natančnost praktično ne vpliva ($13{,}4$ do
$13{,}6\,\%$ za vsa preizkušena razmerja), prav tako ne zamenjava linearne
padajoče funkcije z eksponentno ($13{,}4$ proti $13{,}3\,\%$). Edina parametra,
ki opazno premakneta rezultat, sta torej širina letnikovega okna in namerni
popust: ta predznačeno napako proti oglaševani ceni premakne iz $+3{,}8\,\%$ v
$-6{,}8\,\%$, MdAPE pa zviša z $13{,}4$ na $14{,}5\,\%$. Ta premik je posledica
zavestne odločitve, da sistem ocenjuje pričakovano prodajno in ne oglaševane
cene, in ne slabše ocene.

\begin{table}[htbp]
  \centering
  \caption{Občutljivost na uteži, obliko jedra in popust.}
  \label{tab:sens-ostalo}
  % vir: analysis/out/sensitivity_weights.csv, sensitivity_kernel.csv, sensitivity_discount.csv
  \begin{tabularx}{\textwidth}{@{}l X r >{\raggedleft\arraybackslash}p{2.3cm}@{}}
    \toprule
    \textbf{Parameter} & \textbf{Nastavitev} & \textbf{MdAPE [\%]} & \textbf{Predznačena napaka [\%]} \\
    \midrule
    uteži (km / letnik) & $0{,}5 / 0{,}5$              & 13{,}4 & $+3{,}7$ \\
    uteži (km / letnik) & $0{,}6 / 0{,}4$ (privzeto)   & 13{,}4 & $+3{,}8$ \\
    uteži (km / letnik) & $0{,}8 / 0{,}2$              & 13{,}6 & $+3{,}8$ \\
    \midrule
    jedro               & linearno (privzeto)         & 13{,}4 & $+3{,}8$ \\
    jedro               & eksponentno                 & 13{,}3 & $+3{,}8$ \\
    \midrule
    popust              & brez                        & 13{,}4 & $+3{,}8$ \\
    popust              & stopenjski (privzeto)        & 14{,}5 & $-6{,}8$ \\
    \bottomrule
  \end{tabularx}
\end{table}

\paragraph{Sklep o robustnosti.}
Algoritem je robusten: na večino svojih konstant se skoraj ne odziva, kar
pomeni, da njihova hevristična izbira ni sporna. Njegova natančnost je odvisna
predvsem od kakovosti in gostote podatkov (razdelek~\ref{sec:natancnost}), ne
od finega nastavljanja parametrov.

% =====================================================================
\section{Latenca posameznega vrednotenja}
\label{sec:latenca}

Merimo latenco vrste (a) in (b) iz razdelka~\ref{sec:mere}. Meritev izvede
\texttt{benchmarks/bench\_algo.py}: $2000$ neposrednih klicev funkcije nad
naključnimi realnimi vhodi (brez HTTP), po $100$ ogrevalnih klicih, ki se ne
merijo. Ločeno merimo čas poizvedbe SQL, čas točkovanja v Pythonu in čas
serializacije odgovora.

\paragraph{Ozko grlo: pregled celotne tabele.}
Prvotna izvedba poizvedbe ni imela ustreznega indeksa.
Izsek~\ref{lst:explain-before} prikazuje izpis \texttt{EXPLAIN (ANALYZE,
BUFFERS)}: načrt je \texttt{Seq Scan}, baza pri vsakem vrednotenju prebere
celotno tabelo in zavrže $85.273$ vrstic, ki ne ustrezajo pogojem
(\texttt{Rows Removed by Filter}). Ker so podatki v celoti v predpomnilniku
(\texttt{shared hit}, brez \texttt{read}), je čas stabilen, a raste linearno z
velikostjo tabele.

\begin{center}
\captionof{listing}{\texttt{EXPLAIN (ANALYZE, BUFFERS)} pred dodajanjem indeksa.}
\label{lst:explain-before}
\begin{minted}{text}
Seq Scan on normalized_market_listing
    (cost=0.00..3314.15 rows=12 width=97)
    (actual time=0.555..6.847 rows=573 loops=1)
  Filter: ((kilometers >= 18500) AND (kilometers <= 98500)
       AND (registration_year >= 2019) AND (registration_year <= 2025)
       AND (brand_name = 'Citroën') AND (model_name = 'C3'))
  Rows Removed by Filter: 85273
  Buffers: shared hit=1168
Execution Time: 6.880 ms
\end{minted}
\end{center}

\paragraph{Učinek sestavljenega indeksa.}
Po dodajanju sestavljenega indeksa na
$(\text{brand\_name},\ \text{model\_name},\ \text{registration\_year},\
\text{kilometers})$ se načrt spremeni v \texttt{Bitmap Index Scan}
(izsek~\ref{lst:explain-after}): baza prebere le nekaj strani indeksa in
pripadajoče vrstice, čas izvajanja pa pade z $6{,}88$~ms na $0{,}69$~ms.
Latenca poizvedbe SQL se v mediani zmanjša z $9{,}7$~ms na $2{,}0$~ms
(tabela~\ref{tab:indeks-pred-po}, slika~\ref{fig:indeks-pred-po}). Pomembneje
od trenutnega pohitrenja je, da postane čas poizvedbe neodvisen od velikosti
tabele: namesto $O(\text{tabela})$ postane $O(\text{zadetki})$.

\begin{table}[htbp]
  \centering
  \caption{Latenca poizvedbe SQL pred in po dodajanju indeksa, v milisekundah.}
  \label{tab:indeks-pred-po}
  % vir: benchmarks/report.md (results_algo_before/after.json)
  \begin{tabular}{lrrrr}
    \toprule
    & \textbf{p50} & \textbf{p95} & \textbf{p99} & \textbf{Povprečje} \\
    \midrule
    pred (Seq Scan)  & 9{,}7 & 21{,}0 & 24{,}8 & 11{,}3 \\
    po (Index Scan)  & 2{,}0 & 8{,}6  & 12{,}5 & 3{,}0 \\
    \bottomrule
  \end{tabular}
\end{table}

\begin{figure}[htbp]
  \centering
  \includegraphics[width=0.6\textwidth]{latency_index_before_after.pdf}
  \caption{Latenca poizvedbe SQL pred in po dodajanju sestavljenega indeksa.}
  \label{fig:indeks-pred-po}
  % vir: analysis/out/latency_index_before_after.pdf
\end{figure}

\begin{center}
\captionof{listing}{\texttt{EXPLAIN (ANALYZE, BUFFERS)} po dodajanju indeksa.}
\label{lst:explain-after}
\begin{minted}{text}
Bitmap Heap Scan on normalized_market_listing
    (cost=5.84..50.49 rows=12 width=97)
    (actual time=0.202..0.615 rows=573 loops=1)
  Recheck Cond: (brand_name = 'Citroën' AND model_name = 'C3'
       AND registration_year BETWEEN 2019 AND 2025
       AND kilometers BETWEEN 18500 AND 98500)
  Buffers: shared hit=195
  ->  Bitmap Index Scan on idx_nml_valuation
        (actual time=0.166..0.166 rows=573 loops=1)
        Buffers: shared hit=6
Execution Time: 0.688 ms
\end{minted}
\end{center}

\paragraph{Razčlenitev latence po komponentah.}
Tabela~\ref{tab:latenca-vrednotenja} in slika~\ref{fig:latency-breakdown}
prikazujeta latenco celotnega vrednotenja po dodajanju indeksa. Tipično
vrednotenje traja $2{,}0$~ms, v najslabšem odstotku primerov (p99) pa
$13{,}5$~ms. Skoraj celotna latenca ($\approx 88\,\%$ povprečja) odpade na
poizvedbo SQL; točkovanje v Pythonu ($0{,}4$~ms povprečno) in serializacija
odgovora ($\approx 0{,}01$~ms) sta zanemarljiva. Poizvedba v povprečju vrne
$287$ primerljivih vozil (mediana $175$, največ $1495$).

\begin{table}[htbp]
  \centering
  \caption{Latenca posameznega vrednotenja po dodajanju indeksa
  ($N = 2000$), v milisekundah.}
  \label{tab:latenca-vrednotenja}
  % vir: benchmarks/report.md (results_algo.json)
  \begin{tabular}{lrrrr}
    \toprule
    \textbf{Komponenta} & \textbf{p50} & \textbf{p95} & \textbf{p99} & \textbf{Povprečje} \\
    \midrule
    Skupaj                & 2{,}0  & 9{,}7 & 13{,}5 & 3{,}1 \\
    Poizvedba SQL         & 1{,}8  & 8{,}4 & 11{,}5 & 2{,}7 \\
    Točkovanje (Python)   & 0{,}2  & 1{,}2 & 1{,}8  & 0{,}4 \\
    Serializacija (JSON)  & 0{,}01 & 0{,}02 & 0{,}03 & 0{,}01 \\
    \bottomrule
  \end{tabular}
\end{table}

\begin{figure}[htbp]
  \centering
  \includegraphics[width=0.6\textwidth]{latency_breakdown.pdf}
  \caption{Razčlenitev latence posameznega vrednotenja po komponentah.}
  \label{fig:latency-breakdown}
  % vir: analysis/out/latency_breakdown.pdf
\end{figure}

Vrednotenje se torej izvede daleč pod nefunkcionalno zahtevo $2$~s
(razdelek~\ref{sec:zahteve}); prostora za rast podatkovne baze je ob
prisotnosti indeksa dovolj.

% =====================================================================
\section{Obnašanje pod obremenitvijo}
\label{sec:obremenitev}

Merimo end-to-end latenco (vrsta (c)) in prepustnost končne točke
\texttt{POST /valuate} pri naraščajoči sočasnosti. Obremenitev ustvarja
skripta \texttt{benchmarks/loadtest.py} z več vzporednimi odjemalci, ki v
fiksnem $20$-sekundnem oknu neprekinjeno pošiljajo zahtevke z naključnimi
realnimi vhodi. Test se samodejno prekine ob deležu napak nad $5\,\%$ ali
$p_{95}$ nad $5$~s. Med testom vzorčimo tudi zasedenost baze
(\texttt{pg\_stat\_activity}).

\begin{table}[htbp]
  \centering
  \caption{Prepustnost in latenca pri naraščajoči sočasnosti.
  Latence so v milisekundah, $L$ je ocena po Littlovem zakonu.}
  \label{tab:obremenitev}
  % vir: benchmarks/results_load.json (loadtest.py --db-sample)
  \begin{tabular}{r r r r r r r}
    \toprule
    \textbf{Sočasnost} & \textbf{Prep.\ [z/s]} & \textbf{p50} & \textbf{p95} & \textbf{p99}
      & \textbf{Napake} & \textbf{$L$} \\
    \midrule
    1  & \res{x} & \res{x} & \res{x} & \res{x} & \res{x} & \res{x} \\
    5  & \res{x} & \res{x} & \res{x} & \res{x} & \res{x} & \res{x} \\
    10 & \res{x} & \res{x} & \res{x} & \res{x} & \res{x} & \res{x} \\
    20 & \res{x} & \res{x} & \res{x} & \res{x} & \res{x} & \res{x} \\
    40 & \res{x} & \res{x} & \res{x} & \res{x} & \res{x} & \res{x} \\
    \bottomrule
  \end{tabular}
\end{table}

\paragraph{Rezultat.}
\res{Prepustnost narašča približno linearno do sočasnosti X, nato se ustali pri
približno X zahtevkih na sekundo. Vse latence ostanejo pod nefunkcionalno
zahtevo 2~s.}

\paragraph{Skladnost meritve (Littlov zakon).}
\res{Stolpec $L$ (prepustnost krat povprečna latenca) ostaja blizu nastavljene
sočasnosti, kar potrjuje notranjo skladnost meritve.}

\paragraph{Dokaz ozkega grla.}
Teoretični strop postavitve je $1000 / S \approx 324$ vrednotenj na sekundo,
kjer je $S = 3{,}1$~ms povprečni čas ene evalvacije. Dejanska meja je nižja,
ker ves promet serializira \emph{ena sama deljena povezava do baze}
(\texttt{app/database.py}) in en sam strežniški proces.
\res{To potrjuje vzorčenje baze med testom: število aktivnih poizvedb kljub
naraščajoči sočasnosti odjemalcev ne preseže X.}

% =====================================================================
\section{Preverjanje funkcionalnih zahtev}
\label{sec:funkcionalne}

Tabela~\ref{tab:sledljivost} povezuje vsako funkcionalno zahtevo iz
razdelka~\ref{sec:zahteve} z načinom preverjanja in izidom.

\begin{table}[htbp]
  \centering
  \caption{Sledljivost funkcionalnih zahtev do preverjanja.}
  \label{tab:sledljivost}
  % vir: ročni scenariji + imena Jest testov v car-stock-market-be
  \begin{tabularx}{\textwidth}{@{}X X c@{}}
    \toprule
    \textbf{Zahteva} & \textbf{Način preverjanja} & \textbf{Izid} \\
    \midrule
    Registracija uporabnika              & \res{test / scenarij} & \res{OK} \\
    Prijava in odjava                    & \res{test / scenarij} & \res{OK} \\
    Objava oglasa (le registrirani)      & \res{test / scenarij} & \res{OK} \\
    Vrednotenje vozila po atributih      & \res{test / scenarij} & \res{OK} \\
    Všečkanje oglasov                    & \res{test / scenarij} & \res{OK} \\
    Ažuriranje lastnih/všečkanih oglasov & \res{test / scenarij} & \res{OK} \\
    Pregled oglasov (anonimni)           & \res{test / scenarij} & \res{OK} \\
    \bottomrule
  \end{tabularx}
\end{table}

\paragraph{Pokritost testov.}
Avtomatizirani testi zalednega sistema (Jest + Supertest,
razdelek~\ref{sec:testiranje-be}) obsegajo \res{X} testov v \res{X} zbirkah in
pokrivajo \res{X}\,\% stavkov ter \res{X}\,\% vej. Testi se izvedejo ob vsakem
zahtevku za združitev na cevovodu neprekinjene integracije
(razdelek~\ref{subsec:ci}).

% =====================================================================
\section{Preverjanje nefunkcionalnih zahtev}
\label{sec:nefunkcionalne}

Tabela~\ref{tab:nfr} za vsako nefunkcionalno zahtevo iz
razdelka~\ref{sec:zahteve} navaja merilo, izmerjeno vrednost in sodbo.

\begin{table}[htbp]
  \centering
  \caption{Preverjanje nefunkcionalnih zahtev.}
  \label{tab:nfr}
  \setlength{\tabcolsep}{4pt}
  \begin{tabularx}{\textwidth}{@{}>{\raggedright\arraybackslash}X
                                 >{\raggedright\arraybackslash}X
                                 >{\raggedright\arraybackslash}X c@{}}
    \toprule
    \textbf{Zahteva} & \textbf{Merilo} & \textbf{Izmerjeno} & \textbf{Sodba} \\
    \midrule
    Odziv vmesnika na dejanje uporabnika & $< 100$~ms & \res{x}~ms  & \res{OK} \\
    Zahtevki HTTP (razen vrednotenja)    & $< 300$~ms & \res{x}~ms ($p_{95}$) & \res{OK} \\
    Vrednotenje vozila                   & $< 2$~s    & $13{,}5$~ms ($p_{99}$, brez omrežja) & OK \\
    Vizualna privlačnost / dostopnost    & Lighthouse & \res{x} / \res{x} & \res{OK} \\
    Prilagodljiva postavitev             & preizkus na več velikostih zaslona & \res{da} & \res{OK} \\
    Varovanje uporabniških podatkov      & bcrypt, JWT z rokom & \res{opis} & \res{OK} \\
    Odzivnost pri zahtevnejših opravilih & asinhrona obdelava & \res{opis} & \res{OK} \\
    \bottomrule
  \end{tabularx}
\end{table}

\paragraph{Meritev odzivnosti vmesnika in endpointov.}
\res{Odzivnost uporabniškega vmesnika smo izmerili z orodjem Lighthouse na treh
ključnih straneh; ocene za performanco so X, za dostopnost pa X. Latenco
glavnih endpointov zalednega sistema smo izmerili z orodjem autocannon:
$p_{95}$ je v vseh primerih pod X~ms, torej pod zahtevo 300~ms.}

\paragraph{Varnost.}
Gesla so pred shranjevanjem zgoščena z algoritmom bcrypt
(razdelek~\ref{sec:avtentikacija}), žetoni JWT imajo rok veljavnosti $7$~dni,
poti vmesnika API pa ščiti sredinski preverjevalnik žetonov. \res{Orodje
\texttt{npm audit} ob zadnjem pregledu ni javilo kritičnih ranljivosti v
odvisnostih.}

% =====================================================================
\section{Uporabniški preizkus}
\label{sec:uporabniki}

Izvedli smo formativni uporabniški preizkus~\cite{nielsen1994} z namenom
odkriti težave uporabnosti, ne pa statistično ovrednotiti izdelka.

\paragraph{Udeleženci in postopek.}
V preizkus smo vključili \res{N} udeležencev iz kroga družine in bližnjih
prijateljev, katerih starost in računalniška pismenost sta se med seboj
razlikovali (tabela~\ref{tab:udelezenci}). Vsak udeleženec je prejel enak
seznam nalog: registracija, objava oglasa, všečkanje oglasa, izbira
slovenskega jezika v nastavitvah in vrednotenje poljubnega vozila. Med
izvajanjem je udeleženec razmišljal na glas, opazovalec pa je beležil, ali je
nalogo dokončal samostojno, in vsako opaženo težavo.

\begin{table}[htbp]
  \centering
  \caption{Udeleženci uporabniškega preizkusa.}
  \label{tab:udelezenci}
  \begin{tabular}{c c c l}
    \toprule
    \textbf{\#} & \textbf{Starost (razpon)} & \textbf{Rač.\ pismenost (1--5)} & \textbf{Naprava} \\
    \midrule
    1 & \res{x} & \res{x} & \res{x} \\
    2 & \res{x} & \res{x} & \res{x} \\
    3 & \res{x} & \res{x} & \res{x} \\
    \res{\dots} & & & \\
    \bottomrule
  \end{tabular}
\end{table}

\paragraph{Rezultati.}
Vsi udeleženci so vse naloge dokončali brez pomoči, kar potrjuje, da so glavne
poti v vmesniku dovolj jasne (RV4). Kljub temu so odkrili nekaj
pomanjkljivosti; tabela~\ref{tab:tezave} jih razvršča po resnosti (lestvica
$0$--$4$ po Nielsenu~\cite{nielsen1994}) in navaja, ali so bile odpravljene.

\begin{table}[htbp]
  \centering
  \caption{Težave, odkrite med uporabniškim preizkusom
  (resnost po lestvici $0$--$4$).}
  \label{tab:tezave}
  \setlength{\tabcolsep}{5pt}
  \begin{tabularx}{\textwidth}{@{}>{\raggedright\arraybackslash}X c c@{}}
    \toprule
    \textbf{Težava} & \textbf{Resnost} & \textbf{Status} \\
    \midrule
    Premajhni naslovi podstrani v glavi glavne strani       & \res{x} & odpravljeno \\
    Manjkajoče polje za ponovni vnos gesla pri registraciji & \res{x} & odpravljeno \\
    Velika števila brez ločil za tisočice                   & \res{x} & odpravljeno \\
    Manjkajoč opis pod slikami oblik vozil v obrazcu        & \res{x} & odpravljeno \\
    Premalo vozil, dostopnih za vrednotenje                 & \res{x} & \res{v načrtu} \\
    \bottomrule
  \end{tabularx}
\end{table}

Omejitve preizkusa (majhen in priložnostni vzorec, brez standardiziranega
vprašalnika, brez merjenja časa po nalogah) obravnavamo v
razdelku~\ref{sec:groznje}.

% =====================================================================
\section{Simulacija uporabe}
\label{sec:simulacija}

Za ponazoritev delovanja rešitve povzemamo tipičen potek uporabe. Uporabnik je
ob vstopu usmerjen na glavno stran (slika~\ref{fig:glavna-stran}), kjer sta
izpostavljeni prijava in registracija ter pregled novih in vseh objavljenih
oglasov. Registriran uporabnik lahko objavi oglas prek vodene forme
(slika~\ref{fig:objava-oglasa}), ovrednoti poljubno vozilo, všečka oglase in si
ogleda svoj profil z všečkanimi vozili (slika~\ref{fig:liked-cars}).

\begin{figure}[H]
  \centering
  \includegraphics[width=0.9\textwidth]{glavna-stran.png}
  \caption{Glavna stran prototipa}
  \label{fig:glavna-stran}
\end{figure}
\begin{figure}[H]
  \centering
  \includegraphics[width=0.9\textwidth]{objava-oglasa.png}
  \caption{Vodena forma za objavo oglasa}
  \label{fig:objava-oglasa}
\end{figure}
\begin{figure}[H]
  \centering
  \includegraphics[width=0.9\textwidth]{liked-cars.png}
  \caption{Stran uporabnika s prikazanimi všečkanimi vozili}
  \label{fig:liked-cars}
\end{figure}

\paragraph{Primerjava ocen z obstoječimi portali.}
Za ilustracijo smo delovanje sistema primerjali s ponudbo na dveh spletnih
portalih na treh vozilih (tabela~\ref{tab:vrednotenje}). Sistematično in
kvantitativno oceno natančnosti podaja razdelek~\ref{sec:natancnost}; ta
primerjava je zgolj kvalitativna ponazoritev.

\begin{table}[htbp]
  \centering
  \caption{Primerjava ocenjenih cen s ponudbo na spletnih portalih (v evrih).}
  \label{tab:vrednotenje}
  \setlength{\tabcolsep}{5pt}
  \begin{tabular}{lccc}
    \toprule
    & \textbf{Audi A3} & \textbf{BMW ser.\ 2} & \textbf{Citroën Berlingo} \\
    \midrule
    slov.\ portal (razpon)     & 4.000--6.000  & 9.000--12.000  & 7.500--10.000 \\
    nem.\ portal (razpon)      & 4.500--9.000  & pribl.\ 11.100 & 8.000--12.000 \\
    \midrule
    Sistem --- dobra cena      & 4.500  & 7.600   & 7.600  \\
    Sistem --- sredinska cena  & 6.350  & 9.300   & 9.941  \\
    Sistem --- slaba cena      & 7.600  & 10.500  & 12.100 \\
    Raznolikost                & $49\,\%$ & $31\,\%$ & $45\,\%$ \\
    \bottomrule
  \end{tabular}
\end{table}

Pri vozilih z nizko raznolikostjo ponudbe (Audi~A3, BMW serije~2) se sredinska
ocena dobro ujema s cenami na obeh portalih. Pri modelu Citroën Berlingo je
raznolikost bistveno višja, zato je razpon med dobro in slabo oceno širši,
sredinska ocena pa manj zanesljiva --- kar je skladno z ugotovitvijo iz
razdelka~\ref{sec:natancnost}, da napaka strmo narašča z raznolikostjo ponudbe.
Pri interpretaciji je treba upoštevati, da sistem ocenjuje pričakovano prodajno
in ne oglaševano ceno.

% =====================================================================
\section{Grožnje veljavnosti}
\label{sec:groznje}

\paragraph{Notranja veljavnost.}
Kot merilo resnice pri merjenju natančnosti uporabljamo \emph{oglaševano} ceno,
ne dejanske prodajne cene, ki ni javno dostopna. Namerni popust, s katerim
sistem cilja pričakovano prodajno ceno, je zato ocenjen in ne umerjen na
dejanskih prodajah. Isti fizični oglas se lahko pojavi pri več virih; prečno
preverjanje izloči zapise z istim izvornim oglasom, ne more pa zaznati
podvojitev med viri.

\paragraph{Zunanja veljavnost.}
Podatki pokrivajo $10$ znamk in pretežno slovenski trg (z manjšim nemškim
vzorcem); znamka z največ oglasi na trgu (okoli $8000$ oglasov) v šifrantu ni
zajeta in je iz analize izpuščena. Uporabniški preizkus je temeljil na majhnem
in priložnostnem vzorcu.

\paragraph{Konstruktna veljavnost.}
Pojem \enquote{tržna vrednost} operacionaliziramo kot uteženo povprečje
oglaševanih cen primerljivih vozil, znižano za pričakovani pogajalski popust;
\enquote{raznolikost ponudbe} pa kot relativno širino interkvartilnega razpona
teh cen. Drugačna operacionalizacija bi lahko dala drugačne rezultate.

\paragraph{Zanesljivost.}
Meritve latence in obremenitve so bile izvedene na deljenem produkcijskem
strežniku, kjer lahko na rezultate vpliva ozadje. Meritve natančnosti so
ponovljive nad fiksnim posnetkom baze.

% =====================================================================
\section{Povzetek ugotovitev}
\label{sec:povzetek-analize}

Tabela~\ref{tab:rv-povzetek} povzema odgovore na raziskovalna vprašanja glede
na vnaprej postavljena merila.

\begin{table}[htbp]
  \centering
  \caption{Povzetek odgovorov na raziskovalna vprašanja.}
  \label{tab:rv-povzetek}
  \begin{tabularx}{\textwidth}{@{}l X c@{}}
    \toprule
    \textbf{RV} & \textbf{Glavni rezultat} & \textbf{Sodba} \\
    \midrule
    RV1 natančnost & MdAPE $13{,}5\,\%$ ($9$--$10\,\%$ za dobro pokrita vozila),
                     bistveno bolje od naivne ocene ($37{,}8\,\%$), primerljivo z
                     drugimi metodami & izpolnjeno \\
    RV2 odvisnost  & napaka pada z velikostjo vzorca (ustali se pri $n\approx20$),
                     narašča s starostjo, kilometrino in raznolikostjo ponudbe
                     ($7{,}5\,\%$ do $29{,}7\,\%$) & pojasnjeno \\
    RV3 hitrost    & vrednotenje $2$~ms (p50), $13{,}5$~ms (p99); indeks odpravi
                     odvisnost od velikosti baze; \res{meja X~z/s} & \res{izpolnjeno} \\
    RV4 uporabnost & vse naloge dokončane brez pomoči; \res{X} težav odpravljenih & izpolnjeno \\
    \bottomrule
  \end{tabularx}
\end{table}

\paragraph{Sodba o zastavljenih izhodiščih.}
Rezultati \emph{potrjujejo} izhodišča, ki smo si jih zastavili na začetku. Sistem na podlagi štirih
atributov vozila vrne oceno vrednosti z mediano odstotne napake $13{,}5\,\%$
(brez namernega popusta), pri dobro pokritih vozilih pa $9$--$10\,\%$, kar je
znotraj zastavljenega praga $15\,\%$ in bistveno bolje od naivne metode.
Predlagani pristop se med primerljivo natančnimi metodami ne uvršča na prvo
mesto --- preprost $k$-najbližjih sosedov doseže $11{,}2\,\%$ --- je pa edini,
ki poleg točkovne ocene vrne tudi razložljiv cenovni razpon in indeks
zaupanja, ta pa se izkaže za dober pokazatelj dejanske zanesljivosti
posamezne ocene. Glavne omejitve so odvisnost od enega prevladujočega
podatkovnega vira, nižja natančnost pri starejših vozilih in modelih z
razpršeno ponudbo ter merilo resnice, ki temelji na oglaševanih in ne
prodajnih cenah. Konkretne smeri za izboljšave --- vključno s prehodom na
metode strojnega učenja, ki jih ta analiza dodatno utemeljuje --- navajamo v
poglavju~\ref{ch:sklep}.
